%%%%%%%%%%%%%%%%%%%%%%%%%%%%%%%%%%%%%%%%%%%%%%%%%%%%%%%%%%%%%%%%%%%%%%
% Lớp 11 | Giới hạn. Hàm số liên tục | Hàm số liên tục
% Dạng 1D3?3-4: Hàm số liên tục trên khoảng, đoạn
% Số câu: 137  (sinh tự động bởi tach_cau.py; từ nay thầy sửa trực tiếp file này)
%%%%%%%%%%%%%%%%%%%%%%%%%%%%%%%%%%%%%%%%%%%%%%%%%%%%%%%%%%%%%%%%%%%%%%

%%%%%%%%%%%%%%%% TRẮC NGHIỆM 4 PHƯƠNG ÁN %%%%%%%%%%%%%%%%

% nguồn: Ok_TracNghiemLop11_Lan1
\begin{ex}%[1D3N3-4]%[ma:f66b4fa0]
	Cho hàm số $f(x)=\dfrac{3}{x+3}$. Khẳng  định nào sau đây đúng?
	 \choice
		{\True Hàm số liên tục trên mỗi khoảng $(-\infty;-3)$ và $(-3;+\infty)$} 
		{Hàm số liên tục trên mỗi khoảng $(-\infty;0)$ và  $(0;+\infty)$}
		{Hàm số liên tục  trên $\mathbb{R}$}
		 {Hàm số liên tục trên mỗi khoảng $(-\infty;3)$ và $(3;+\infty)$}
	\loigiai{ 
		Hàm số $f(x)=\dfrac{3}{x+3}$ có tập xác định là $x+3\ne 0\Rightarrow x\ne -3$.\\ 
		Nên hàm số liên tục trên mỗi khoảng $(-\infty;-3)$ và $(-3;+\infty)$. 
	} 
\end{ex}

% nguồn: Ok_TracNghiemLop11_Lan1
\begin{ex}%[1D3N3-4]%[ma:51a41679]
	Cho hàm số  $f(x)=\dfrac{x^2+1}{x^2-6x+5}$. Hàm số $f(x)$ liên tục trên khoảng nào sau đây?
	\choice 
		 {($-\infty;3)$}
		{\True $(2;3)$} 
		{$(2;+\infty)$} 
		{$\mathbb{R}$} 
	\loigiai{ 
		Tập xác định $\mathscr{D} =  \mathbb{R} \setminus \{2;3\}$.\\
		Do $f(x)$ là hàm số dạng phân  thức hữu tỷ nên hàm số liên tục trên từng khoảng của tập xác định, tức là hàm số $f(x)$ liên tục trên các khoảng $(-\infty;2)$, $(2;3)$ và $(3;+\infty)$.
	 }
\end{ex}

% nguồn: Ok_Toan11_TracNghiem_Lan2
\begin{ex}%[1D3H3-4]%[ma:e4553a3e]
	Cho hàm số $f(x)=\heva{&\dfrac{\sqrt[3]{4x}-2}{x-2} &\text{ khi } x\ne 2 \\ &ax+3 &\text{ khi } x=2}$. Xác định $a$  để hàm số liên tục trên $ \mathbb{R}$.
	\choice 
		 {$a=-1$}
		 {$a=\dfrac{1}{6}$}
		 {$a=\dfrac{4}{3}$}
		{\True $a=-\dfrac{4}{3}$} 
	 \loigiai{
		Tập xác định của hàm số  là $\mathscr{D}= \mathbb{R}$.
		Nếu $x \ne 2$, ta  có $f(x)=\dfrac{\sqrt[3]{4x}-2}{x-2}$.
		Hàm số $f(x)=\dfrac{\sqrt[3]{4x}-2}{x-2}$ xác định và liên tục trên mỗi  khoảng $\left(-\infty; 2\right)$ và $\left(2;+\infty \right)$.
		Tại $x=2$, ta có: $f(2)=2a+3$. 
		\begin{align*} \lim\limits_{x\to 2} f(x)
			&=\lim\limits_{x\to  2} \dfrac{\sqrt[3]{4x}-2}{x-2} \\
			&=\lim\limits_{x\to 2} \dfrac{\left(\sqrt[3]{4x}-2\right)\left[\left(\sqrt[3]{4x} \right)^2+2\sqrt[3]{4x}+4\right]}{\left(x-2\right)\left[\left(\sqrt[3]{4x}  \right)^2+2\sqrt[3]{4x}+4\right]} \\
			&=\lim\limits_{x\to 2} \dfrac{4\left(x-2\right)}{\left(x-2\right)\left[\left(\sqrt[3]{4x} \right)^2+2\sqrt[3]{4x}+4\right]} \\ 
			&=\lim\limits_{x\to  2} \dfrac{4}{\left(\sqrt[3]{4x} \right)^2+2\sqrt[3]{4x}+4}\\
			 &=\dfrac{1}{3}.
		\end{align*}
		Hàm số liên  tục tại $x=2$ khi và chỉ khi $\lim\limits_{x\to 2} f(x)=f(2)\Leftrightarrow 2a+3=\dfrac{1}{3} \Leftrightarrow a=-\dfrac{4}{3}$.
		 Vậy hàm số liên tục trên $ \mathbb{R}$ khi và chỉ khi $a=-\dfrac{4}{3}$.
	 }
\end{ex}

% nguồn: Ok_Toan11_TracNghiem_Lan2
\begin{ex}%[1D3H3-4]%[ma:0005c3ca]
	Tìm tất cả các  giá trị thực của $m$ để hàm số $f(x)=\heva{
	 &\dfrac{\sqrt{x+1}-1}{x} & \text{nếu } x > 0 \\
	&\sqrt{x^2+1}-m & \text{nếu  } x \le 0
	 }$ liên tục trên $\mathbb{R}$.
	 \choice
		{$m=\dfrac{3}{2}$} 
		 {\True $m=\dfrac{1}{2}$}
		{$m=-2$} 
		 {$m=-\dfrac{1}{2}$}
	\loigiai{ 
		Tập xác định của hàm số là $\mathscr{D}  = \mathbb{R}$.\\
		Xét trên $(0;  +\infty)$, ta có $f(x) = \dfrac{\sqrt{x+1}-1}{x}$ liên tục trên $(0; +\infty)$.\\
		Xét trên $(-\infty; 0)$, ta có $f(x) = \sqrt{x^2+1} - m$ liên tục  trên $(-\infty; 0)$.\\
		Hàm số liên tục trên $\mathbb{R}$ khi và  chỉ khi hàm số liên tục tại $x = 0$.\\
		Ta có 
		\begin{align*}
			\lim\limits_{x \to 0^{+}} f(x) &= \lim\limits_{x \to 0^{+}} \dfrac{\sqrt{x+1}-1}{x} = \lim\limits_{x \to 0^{+}} \dfrac{1}{\sqrt{x+1}+1} = \dfrac{1}{2},  \\
			\lim\limits_{x \to 0^{-}} f(x) &= \lim\limits_{x \to 0^{-}} \left(\sqrt{x^2+1}-m\right) = 1  - m = f(0).
		\end{align*}
		 Do đó, hàm số liên tục tại $x = 0$ khi và chỉ khi
		\[ 
		 \dfrac{1}{2} = 1 - m \Leftrightarrow m = \dfrac{1}{2}.
		 \]
		Vậy hàm số liên tục trên $\mathbb{R}$  khi và chỉ khi $m = \dfrac{1}{2}$.
	} 
\end{ex}

% nguồn: Ok_Toan11_TracNghiem_Lan2
\begin{ex}%[1D3H3-4]%[ma:28462a6b]
	Nếu hàm số $f(x)=\heva{ 
	&x^2 + ax + b &&  \text{nếu } x < -5 \\
	&x + 17 && \text{nếu } -5 \le x \le  10 \\
	&ax + b +  10 && \text{nếu } x > 10
	 }$ liên tục trên $\mathbb{R}$ thì $a + b$ bằng
	\choice 
		 {\True $-1$}
		{$0$} 
		 {$1$}
		 {$2$}
	 \loigiai{
		Tập xác định của hàm  số là $\mathscr{D} = \mathbb{R}$.\\
		Khi $x < -5$, ta có $f(x) = x^2 + ax + b$,  là hàm đa thức, nên liên tục trên $\left(-\infty; -5\right)$.\\
		Khi  $-5 \le x \le 10$, ta có $f(x) = x + 17$, là hàm đa thức, nên liên tục trên $(-5; 10)$.\\
		Khi $x > 10$, ta có  $f(x) = ax + b + 10$, là hàm đa thức, nên liên tục trên $(10; +\infty)$.\\
		Để hàm số liên tục trên $\mathbb{R}$, thì hàm số phải liên tục tại $x = -5$ và $x  = 10$.\\
		 Ta có
		\begin{align*}
			 f(-5) &= 12,\\
			f(10) &= 17,  \\
			\lim\limits_{x \to -5^{-}} f(x) &=  -5a + b + 25, \\
			\lim\limits_{x  \to -5^{+}} f(x) &= 12, \\
			\lim\limits_{x \to 10^{-}} f(x)  &= 27, \\
			\lim\limits_{x \to 10^{+}} f(x)  &= 10a + b + 10.
		\end{align*}
		Hàm số liên tục tại $x = -5$  và $x = 10$ khi:
		 \[
		 \heva{
		&-5a + b +  25 = 12 \\
		&10a  + b + 10 = 27
		 }
		\Leftrightarrow 
		\heva{ 
		&-5a + b  = -13 \\
		&10a  + b = 17
		} 
		\Leftrightarrow 
		 \heva{
		&a = 2 \\ 
		 &b = -3.
		} 
		\] 
		Suy ra $a +  b = -1$.\\
		Vậy hàm số liên tục trên $\mathbb{R}$  khi và chỉ khi $a + b = -1$.
	 }
\end{ex}

% nguồn: Ok_Toan11_TracNghiem_Lan2
\begin{ex}%[1D3H3-4]%[ma:ede01a79]
	Hàm số  $f(x)=\dfrac{2x}{x^2-3x+2}$ liên tục trên khoảng
	\choice 
		{$(0; 2)$} 
		{\True  $(-2; 0)$}
		{$(-\infty;+\infty)$} 
		{$(1;  3)$}
	\loigiai{Hàm số $f(x)=\dfrac{2x}{x^2-3x+2}$ có nghĩa khi $x^2-3x+2\neq  0\Leftrightarrow\heva{&x \neq 1\\
		 &x \neq 2.}$\\
		Do đó, $f(x)$ liên tục trên  từng khoảng $(-\infty;1)$, $(1;2)$, $(2;+\infty)$.\\
		 Vì $(-2;0)\subset(-\infty;1)$ nên chọn đáp án \circled{B}.
	} 
\end{ex}

% nguồn: Ok_Toan11_TracNghiem_Lan2
\begin{ex}%[1D3H3-4]%[ma:8ca988e4]
	Cho hàm số  $y=f(x)=\heva{
	&\dfrac{\sqrt{x^2+16}-5}{x-3} & \text{nếu  } x \neq 3 \\
	&a & \text{nếu } x = 3. 
	 }$\\
	Tập các giá trị của $a$ để hàm số đã cho liên  tục trên $S$ là
	\choice 
		{$\left\{\dfrac{2}{5}  \right\}$}
		{$\left\{\dfrac{1}{5} \right\}$} 
		{$\{0\}$} 
		{\True $\left\{\dfrac{3}{5}  \right\}$}
	\loigiai{ 
		Tập xác định $\mathscr{D}=\mathbb{R}$.\\ 
		Khi $x  \neq 3$, ta có $f(x) = \dfrac{\sqrt{x^2+16}-5}{x-3}$ xác định và liên tục trên các khoảng $\left(-\infty; 3\right)$ và $\left(3; +\infty \right)$.\\
		Khi  $x = 3$, ta có $f(3) = a$ và
		\[ 
		\lim\limits_{x \to 3} f(x) = \lim\limits_{x \to 3} \dfrac{\sqrt{x^2+16}-5}{x-3} =  \lim\limits_{x \to 3} \dfrac{x+3}{\sqrt{x^2+16}+5} = \dfrac{3}{5}.
		 \]
		Hàm số liên tục trên $\mathbb{R}$ khi và chỉ khi hàm số liên  tục tại điểm $x = 3$.\\
		Khi  đó $\lim\limits_{x \to 3} f(x) = f(3) \Leftrightarrow a = \dfrac{3}{5}$.\\
		 Vậy hàm số liên tục trên $\mathbb{R}$ khi và chỉ khi $a = \dfrac{3}{5}$.
	 }
\end{ex}

% nguồn: Ok_Toan11_TracNghiem_Lan2
\begin{ex}%[1D3H3-4]%[ma:a9f2cfbc]
	Cho hàm số $f(x)=\dfrac{x^2+2}{x^2+5x+6}$. Hàm số  đã cho liên tục trên khoảng nào sau đây?
	 \choice
		{$(-3;3)$} 
		{$(-\infty;3)$} 
		 {$(-3;2)$}
		{\True  $(-2;+\infty)$}
	\loigiai{ 
		Tập xác định $\mathscr{D}=\mathbb{R}\setminus  \left\{ -2;-3 \right\}$.\\
		Hàm  số $f(x)=\dfrac{x^2+2}{x^2+5x+6}$ là hàm phân thức, liên tục trên từng khoảng xác định tập xác định, tức liên tục trên $(-\infty;-3)$; $(-3;-2)$ và $(-2;+\infty)$.
	 }
\end{ex}

% nguồn: Ok_Toan11_TracNghiem_Lan2
\begin{ex}%[1D3H3-4]%[ma:3fb2a466]
	Trong các hàm số sau, hàm số nào  liên tục trên $\mathbb{R}$?
	\choice 
		 {$f(x)=\sqrt{x-5}$}
		 {\True $f(x)=\dfrac{x+5}{x^2+4}$}
		{$f(x)=\cot  x+3$}
		 {$f(x)=\dfrac{x^2+3}{2-x}$}
	 \loigiai{
		Hàm số $f(x)=\dfrac{x+5}{x^2+4}$ là hàm  phân thức hữu tỉ và có tập xác định là $\mathscr{D}=\mathbb{R}$ do đó hàm số
		$f(x)=\dfrac{x+5}{x^2+4}$ liên tục  trên $\mathbb{R}$.
	 }
\end{ex}

% nguồn: Ok_Toan11_TracNghiem_Lan2
\begin{ex}%[1D3H3-4]%[ma:f4f08655]
	Trong các hàm số  sau, hàm số nào liên tục trên $\mathbb{R}$?
	 \choice
		 {\True $y=x^3-x$}
		{$y=\cot x$} 
		 {$y=\dfrac{2x-1}{x-1}$}
		{$y=\sqrt{x^2-1}$} 
	\loigiai{Vì $y=x^3-x$ là đa thức nên  nó liên tục trên $\mathbb{R}$.
	} 
\end{ex}

% nguồn: Ok_Toan11_TracNghiem_Lan2
\begin{ex}%[1D3H3-4]%[ma:24e73ed9]
	 Hàm số nào dưới đây liên tục trên khoảng $(-\infty;+\infty)$?
	\choice 
		 {$y=x+\dfrac{1}{x}$}
		{$y=\sqrt{2-x}$} 
		{$y=\dfrac{2x+1}{x-7}$} 
		 {\True $y=x+7$}
	 \loigiai{
		Hàm số đa thức $y=x+7$ liên  tục trên toàn bộ tập số thực $\mathbb{R}$.
	} 
\end{ex}

% nguồn: Ok_Toan11_TracNghiem_Lan2
\begin{ex}%[1D3H3-4]%[ma:e10784d0]
	Cho hàm số $y=\heva{&-x^2+x+3  &\text{khi } x \geq 2\\
	&5x+2 &\text{khi } x < 2}$. Chọn mệnh đề sai  trong các mệnh đề sau.
	\choice 
		 {Hàm số liên tục tại $x_0=1$}
		{\True Hàm số liên  tục trên $\mathbb{R}$}
		{Hàm số  liên tục trên các khoảng $(-\infty; 2)$, $(2;+\infty)$}
		{Hàm  số gián đoạn tại $x_0=2$}
	\loigiai{ 
		Với $x > 2$, ta có $f(x)=-x^2+x+3$ là hàm đa thức  nên hàm số $f(x)$ liên tục trên khoảng $(2;+\infty)$.\\
		Với $x < 2$, ta có $f(x)=5x+2$ là hàm đa thức nên hàm số  $f(x)$ liên tục trên khoảng $(-\infty; 2)$.\\
		Tại $x=2$,  ta có
		\begin{align*}
			&\lim\limits_{x \to 2^{+}}  f(x)=\lim\limits_{x \to 2^{+}}\left(-x^2+x+3\right)=1\\
			&\lim\limits_{x  \to 2^{-}} f(x)=\lim\limits_{x \to 2^{-}}(5x+2)=12\\
			&\Rightarrow \lim\limits_{x \to 2^{+}} f(x) \neq \lim\limits_{x \to 2^{-}}  f(x).
		\end{align*}
		Do đó không tồn tại $\lim\limits_{x \to 2} f(x)$ nên hàm số gián đoạn tại $x_0=2$, tức là hàm  số không liên tục trên $\mathbb{R}$.
	 }
\end{ex}

% nguồn: Ok_Toan11_TracNghiem_Lan2
\begin{ex}%[1D3H3-4]%[ma:d4e4fa58]
	[Dự án 11 HVA  2024-2025]
	Tìm $a$ để hàm số liên tục trên $\mathbb{R}$: $f(x)=\heva{&2x+a &\text{khi  \,}x\le 1\\&\dfrac{x^3-x^2+2x-2}{x-1}&\text{khi \,}x > 1}$.
	 \choice
		{$a=-2$} 
		 {\True $a=1$}
		{$a=2$} 
		{$a=-1$} 
	\loigiai{Tập xác  định $\mathscr{D}=\mathbb{R}$.
		\begin{itemize}
			\item Khi $x < 1$ thì $f(x)=2x+a$ là hàm đa thức nên liên tục trên khoảng $(-\infty; 1)$. 
			\item Khi $x > 1$ thì $f(x)=\dfrac{x^3-x^2+2x-2}{x-1}$ là hàm phân thức hữu tỉ xác định trên khoảng $(1;+\infty)$ nên liên  tục trên khoảng $(1;+\infty)$.
		\end{itemize}
		Hàm số $f(x)$ liên tục trên $\mathbb{R}$ khi và chỉ khi hàm  số $f(x)$ liên tục tại $x=1$.\\
		Xét tính liên tục  của hàm số tại điểm $x=1$, ta có
		\begin{itemize}
			\item $f(1)=2+a$. 
			\item  $\lim\limits_{x\to 1^-}f(x)=\lim\limits_{x\to 1^-}(2x+a)=2+a$.
			 \item $\lim\limits_{x\to 1^+}f(x)=\lim\limits_{x\to 1^+}\dfrac{x^3-x^2+2x-2}{x-1}=\lim\limits_{x\to 1^+}\dfrac{(x-1)(x^2+2)}{x-1}=\lim\limits_{x\to 1^+}(x^2+2)=3$.
		\end{itemize}
		Hàm  số $f(x)$ liên tục tại $x=1$ khi và chỉ khi
		\begin{align*}
			\lim\limits_{x\to  1^-}f(x)=\lim\limits_{x\to 1^+}f(x)=f(1)\Leftrightarrow 2+a=3\Leftrightarrow a=1.
		\end{align*}
	 }
\end{ex}

% nguồn: Ok_Toan11_TracNghiem_Lan2
\begin{ex}%[1D3H3-4]%[ma:4fc6f365]
	Cho hàm số $y=\heva{& 3x+1&\text{ khi } x\ge-1 \\ & x+m &\text{ khi } x <-1}$, $m$ là tham số. Tìm $m$ để hàm số liên tục  trên $\mathbb{R}$.
	 \choice
		 {$m=5$}
		{\True  $m=-1$}
		{$m=3$} 
		 {$m=-3$}
	 \loigiai{
		Ta có  hàm số liên tục trên các khoảng $\left(-\infty;-1\right)$ và $\left(-1;+\infty \right)$.
		Xét tính  liên tục của hàm số tại $x=-1$.
		Ta có $y\left(-1\right)=-2=\lim\limits_{x\to-1^{+}} y$ và $\lim\limits_{x\to-1^{-}} y=-1+m$. 
		Để hàm số liên tục trên $\mathbb{R}$ thì $y\left(-1\right)=\lim\limits_{x\to-1^{+}} y=\lim\limits_{x\to-1^{-}}  y\Leftrightarrow-2=-1+m\Leftrightarrow m=-1$.
	 }
\end{ex}

% nguồn: Ok_Toan11_TracNghiem_Lan2
\begin{ex}%[1D3H3-4]%[ma:cdb2ab27]
	Cho hàm số $f(x)=\dfrac{2x+3}{x-2}$. Mệnh đề nào sau đây đúng? 
	\choice 
		{Hàm số liên tục  trên khoảng $(1; 5)$}
		{Hàm số gián đoạn  tại $x=2020$}
		{Hàm số liên  tục tại $x=2$}
		{\True  Hàm số gián đoạn tại $x=2$}
	 \loigiai{Hàm số $f(x)=\dfrac{2x+3}{x-2}$ không xác định tại $x=2$ nên hàm số này gián đoạn tại $x=2$.
	 }
\end{ex}

% nguồn: Ok_Toan11_TracNghiem_Lan2
\begin{ex}%[1D3H3-4]%[ma:cd7fb0e0]
	Cho hàm số $f(x)=\heva{& \dfrac{x^2-3}{x-\sqrt{3}} &\quad& \text{ khi } x\ne \sqrt{3} \\ &2\sqrt{3} &\quad& \text{ khi } x=\sqrt{3}}$. Hàm số đã cho liên tục trên tập nào  sau đây?
	\choice 
		{\True $\mathbb{R}$} 
		 {$(-\infty;\sqrt{3}]$}
		{$(0;+\infty)$} 
		{$[1;+\infty)$} 
	 \loigiai{
		\begin{itemize}
			 \item Nếu $x\ne \sqrt{3}$ thì $f(x)=\dfrac{x^2-3}{x-\sqrt{3}}$ là hàm phân thức hữu tỉ nên liên tục trên $(-\infty;\sqrt{3})$ và $(\sqrt{3};+\infty)$.
			\item Tại $x=\sqrt{3}$, ta có $f\left(\sqrt{3}\right)=2\sqrt{3}$. 
			\item  Xét $\lim\limits_{x\to \sqrt{3}}\dfrac{x^2-3}{x-\sqrt{3}}=\lim\limits_{x\to \sqrt{3}}\dfrac{\left(x-\sqrt{3}\right)\left(x+\sqrt{3}\right)}{x-\sqrt{3}}=\lim\limits_{x\to \sqrt{3}}(x+\sqrt{3})=2\sqrt{3}$.\\
			Suy ra hàm số liên  tục tại $x=\sqrt{3}$.
		\end{itemize}
		Vậy  hàm số liên tục trên $\mathbb{R}$.
	} 
\end{ex}

% nguồn: Ok_Toan11_TracNghiem_Lan2
\begin{ex}%[1D3H3-4]%[ma:19e3c98d]
	Hàm  số $f(x)=\sqrt{3+x}+\sqrt{4-x}$ liên tục trên
	\choice 
		 {$(3;10)$}
		{\True $[-3;4]$} 
		{$[-3;+\infty)$} 
		 {$(-\infty;4]$}
	\loigiai{ 
		Tập  xác định của hàm số $\mathscr{D}=[-3;4]$.\\
		Vì với mọi $x_0\in (-3;4)$ ta có $\lim\limits_{x\to x_0}f(x)=\lim\limits_{x\to x_0}\left(\sqrt{3+x}+\sqrt{4-x}\right)=f(x_0)$ nên hàm số liên tục trên  khoảng $(-3;4)$.\\
		Ngoài ra  ta có
		\begin{itemize}
			\item $\lim\limits_{x\to (-3)^+}f(x)=\lim\limits_{x\to (-3)^+}\left(\sqrt{3+x}+\sqrt{4-x}\right)=\sqrt{7}=f(-3)$ nên  hàm số liên tục phải tại $x=-3$.
			\item $\lim\limits_{x\to 4^-}f(x)=\lim\limits_{x\to 4^-}\left(\sqrt{3+x}+\sqrt{4-x}\right)=\sqrt{7}=f(4)$ nên hàm số liên  tục trái tại $x=4$.
		\end{itemize}
		Do đó hàm số  đã cho liên tục trên đoạn $[-3;4]$.
	} 
\end{ex}

% nguồn: Ok_Toan11_TracNghiem_Lan2
\begin{ex}%[1D3H3-4]%[ma:593ff058]
	Tìm $P$ để  hàm số $y=\heva{
	&\dfrac{x^2  - 4x + 3}{x - 1} & \text{nếu } x > 1 \\
	&6Px  - 3 & \text{nếu } x \le 1}$ liên tục trên $\mathbb{R}$.
	 \choice
		{$P = \dfrac{5}{6}$} 
		{$P  = \dfrac{1}{2}$}
		{\True $P = \dfrac{1}{6}$} 
		{$P = \dfrac{1}{3}$} 
	\loigiai{ 
		Tập xác  định của hàm số là $\mathscr{D} = \mathbb{R}$.\\
		Xét trên $(-\infty; 1)$, ta có $f(x)  = 6Px - 3$ liên tục trên $(-\infty; 1)$.\\
		Xét  trên $(1; +\infty)$, ta có $f(x) = \dfrac{x^2 - 4x + 3}{x - 1}$ liên tục trên $(1; +\infty)$.\\
		Xét tại $x = 1$, ta có $\lim\limits_{x \to 1^{+}} f(x) = -2$, $\lim\limits_{x \to  1^{-}} f(x) = 6P - 3$, $f(1) = 6P - 3$.\\
		Hàm  số liên tục tại $x = 1$ khi và chỉ khi
		\[ 
		\lim\limits_{x \to 1^{+}} f(x) = \lim\limits_{x \to 1^{-}} f(x) = f(1) \Leftrightarrow 6P - 3 = -2 \Leftrightarrow  P = \dfrac{1}{6}.
		\] 
		Vậy hàm số liên tục trên $\mathbb{R}$ khi và chỉ khi $P = \dfrac{1}{6}$. 
	} 
\end{ex}

% nguồn: Ok_Toan11_TracNghiem_Lan2
\begin{ex}%[1D3H3-4]%[ma:85097f02]
	 Cho hàm số $f(x)=\dfrac{x^2+1}{x^2+5x+6}$. Khi đó hàm số $y=f(x)$ liên tục trên các khoảng nào sau đây?
	\choice 
		 {$(-3; 2)$}
		{\True $(-2;+\infty)$} 
		 {$(-\infty; 3)$}
		 {$(-3; 3)$}
	\loigiai{Hàm số  có nghĩa khi $x^2+5x+6\neq 0\Leftrightarrow\heva{&x \neq-3\\
		 &x \neq-2.}$\\
		Vậy hàm  số $f(x)=\dfrac{x^2+1}{x^2+5x+6}$ liên tục trên khoảng $(-\infty;-3)$; $(-3;-2)$
		và $(-2;+\infty)$. 
	 }
\end{ex}

% nguồn: Ok_Toan11_TracNghiem_Lan2
\begin{ex}%[1D3H3-4]%[ma:3a593bce]
	Hàm số nào trong các hàm số dưới đây không liên tục trên  $\mathbb{R}$?
	\choice 
		 {$y=|x|$}
		{\True $y=\dfrac{x}{x+1}$} 
		{$y=\sin x$} 
		 {$y=\dfrac{x}{|x|+1}$}
	\loigiai{Tập xác định của hàm số  $y=\dfrac{x}{x+1}$ là $\mathbb{R} \setminus\{1\}$. Do đó, hàm số liên tục trên từng khoảng $(-\infty; 1)$ và $(1;+\infty)$ nên hàm số không liên tục trên $\mathbb{R}$.}
\end{ex}

% nguồn: Ok_Toan11_TracNghiem_Lan2
\begin{ex}%[1D3H3-4]%[ma:859abd1e]
	Hàm số nào trong các hàm số sau đây  không liên tục trên $\mathbb{R}$?
	\choice 
		{$y=|x|$} 
		{\True  $y=\dfrac{x}{x-1}$}
		{$y=\sin  x$}
		 {$y=\dfrac{x}{|x|+1}$}
	 \loigiai{
		Ta có tập xác định của hàm số $y=\dfrac{x}{x-1}$ là $\mathbb{R}\setminus \left\{ 1 \right\}$ nên hàm  số không liên tục trên $\mathbb{R}$.
	 }
\end{ex}

% nguồn: Ok_Toan11_TracNghiem_Lan2
\begin{ex}%[1D3H3-4]%[ma:bca307b4]
	Tìm $m$ để hàm số $y =  f(x) = \heva{
	&x^2 + 2\sqrt{x - 2}  & \text{nếu } x \ge 2 \\
	&5x - 5m + m^2 & \text{nếu } x < 2}$ liên  tục trên $\mathbb{R}$?
	\choice 
		{\True $m = 2$; $m  = 3$}
		{$m = -2$; $m = -3$} 
		{$m = 1$; $m  = 6$}
		{$m  = -1$; $m = -6$}
	\loigiai{ 
		 Tập xác định của hàm số là $\mathscr{D} = \mathbb{R}$.\\
		Xét trên $(2; +\infty)$, khi đó $f(x) = x^2 + 2\sqrt{x - 2}$. Hàm  số xác định và liên tục trên $(2; +\infty)$.\\
		Xét trên $(-\infty; 2)$,  khi đó $f(x) = 5x - 5m + m^2$. Hàm số xác định và liên tục trên $(-\infty; 2)$.\\
		Xét tại $x = 2$, ta có $f(2)  = 4$.\\
		Mặt  khác
		\begin{align*}
			\lim\limits_{x \to 2^{+}} f(x) &= \lim\limits_{x \to  2^{+}} \left(x^2 + 2\sqrt{x - 2} \right) = 4,\\
			\lim\limits_{x \to 2^{-}} f(x) &= \lim\limits_{x \to 2^{-}} \left(5x - 5m + m^2 \right) = m^2 - 5m + 10. 
		\end{align*}
		Để hàm số liên tục trên $\mathbb{R}$, nó  phải liên tục tại $x = 2$, tức là
		\[ 
		\lim\limits_{x \to 2^{+}} f(x) = \lim\limits_{x \to 2^{-}} f(x) = f(2) \Leftrightarrow m^2 - 5m + 10 = 4 \Leftrightarrow m^2 - 5m + 6 = 0 \Leftrightarrow  \hoac{&m=2\\&m=3.}
		 \]
		Vậy $m = 2$ hoặc $m = 3$ để hàm  số liên tục trên $\mathbb{R}$.
	} 
\end{ex}

% nguồn: Ok_Toan11_TracNghiem_Lan2
\begin{ex}%[1D3H3-4]%[ma:e9eac5f3]
	Cho  hàm số $f(x)=\heva{
	&\sqrt{x}  - m & \text{nếu } x \ge 0 \\
	&mx + 1 & \text{nếu }  x < 0}$.
	 Tìm tất cả các giá trị của $m$ để $f(x)$ liên tục trên $\mathbb{R}$.
	\choice 
		 {$m=1$}
		{$m=0$} 
		 {\True $m=-1$}
		{$m=-2$} 
	\loigiai{ 
		Tập xác  định của hàm số là $\mathscr{D} = \mathbb{R}$.\\
		Xét trên $(-\infty; 0)$,  ta có $f(x) = mx + 1$ liên tục trên $(-\infty; 0)$.\\
		Xét trên $(0; +\infty)$, ta có $f(x) = \sqrt{x} - m$ liên tục trên $(0;  +\infty)$.\\
		Xét tại $x  = 0$, ta có $\lim\limits_{x \to 0^{+}} f(x) = -m$, $\lim\limits_{x \to 0^{-}} f(x) = 1$, $f(0) = -m$.\\
		Hàm số liên tục tại $x = 0$ khi  và chỉ khi
		 \[
		\lim\limits_{x \to 0^{+}} f(x) = \lim\limits_{x \to 0^{-}} f(x) = f(0) \Leftrightarrow -m = 1 \Leftrightarrow m = -1. 
		 \]
		Vậy hàm số liên  tục trên $\mathbb{R}$ khi và chỉ khi $m = -1$.
	 }
\end{ex}

% nguồn: Ok_Toan11_TracNghiem_Lan2
\begin{ex}%[1D3H3-4]%[ma:c03f0e72]
	Cho hàm số $f(x)=\heva{&\dfrac{x^2-1}{x-1} &\text{khi } x \neq  1\\
	&2 &\text{khi } x=1.}$ Khẳng  định nào sau đây đúng?
	 \choice
		{Hàm số liên tục tại mọi điểm $x \neq 1$ và gián đoạn tại  $x=1$}
		{\True Hàm số liên tục trên  $\mathbb{R}$}
		{Hàm số không liên tục trên  $(1;+\infty)$}
		{Hàm số gián  đoạn tại điểm $x=1$}
	\loigiai{TXĐ:  $\mathscr{D}=\mathbb{R}$.\\
		Hàm số $y=\dfrac{x^2-1}{x-1}(x \neq 1)$ là hàm phân  thức nên liên tục trên $(-\infty; 1)$ và $(1;+\infty)$.
		Ta xét tính liên tục  của $f(x)$ tại $x=1$.
		Ta có 
		\begin{align*}
			&\lim\limits_{x \to 1} f(x)=\lim\limits_{x \to  1} \dfrac{x^2-1}{x-1}=\lim\limits_{x \to 1} \dfrac{(x-1)(x+1)}{x-1}=\lim\limits_{x \to 1}(x+1)=2\\
			&f(1)=2. 
		\end{align*}
		Do đó $\lim\limits_{x \to 1} f(x)=f(1)=2$ nên hàm số $f(x)$  liên tục tại $x=1$.
		Vậy hàm số đã  cho liên tục trên $\mathbb{R}$.
	 }
\end{ex}

% nguồn: Ok_Toan11_TracNghiem_Lan2
\begin{ex}%[1D3H3-4]%[ma:b2db5484]
	Cho hàm số $f(x)=\heva{&-\cos x,& &\text{khi } x<0\\ 
	&\dfrac{x^2}{1+x},&&\text{khi }  0\leq x<1\\
	&x^2,&&\text{khi } x\geq 1}$. Khẳng định nào sau đây  đúng?
	\choice 
		{Hàm số $f(x)$ liên  tục tại mọi điểm $x$ thuộc $\mathbb{R}$}
		{Hàm số $f(x)$ bị gián đoạn tại  điểm $x=0$}
		{\True Hàm số $f(x)$ bị gián  đoạn tại điểm $x=1$}
		{Hàm số $f(x)$ bị gián đoạn tại  điểm $x=0$ và $x=1$}
	 \loigiai{
		Ta có $f(x)$ liên tục tại  $x\ne 0$ và $x\ne 1$.\\
		 Tại $x=0$ thì
		$\lim\limits_{x\to 0^-} f(x)=\lim\limits_{x\to 0^-}  (-\cos x)=-1$, $\lim\limits_{x\to 0^+} f(x)=\lim\limits_{x\to 0^+} \dfrac{x^2}{1+x}=0$ và $f(0)=0$ nên $\lim\limits_{x\to 0^-} f(x)\ne \lim\limits_{x\to 0^+} f(x)=f(0)$. Do đó, hàm số liên tục tại $x=0$.\\
		Tại  $x=1$ thì
		$\lim\limits_{x\to 1^-} f(x)=\lim\limits_{x\to 1^-}  \dfrac{x^2}{1+x}=\dfrac{1}{2}$, $\lim\limits_{x\to 1^+} f(x)=\lim\limits_{x\to 1^+} x^2=1$.\\
		Suy ra $\lim\limits_{x\to 1^-} f(x)\ne  \lim\limits_{x\to 1^+} f(x)$.\\
		 Do đó, hàm số gián đoạn tại $x=1$.
	 }
\end{ex}

% nguồn: Ok_Toan11_TracNghiem_Lan2
\begin{ex}%[1D3H3-4]%[ma:1a132c34]
	 Cho hàm số $f(x)=\heva{
	&3x + a - 1  & \text{nếu } x \le 0 \\
	&\dfrac{\sqrt{1 + 2x} - 1}{x} &  \text{nếu } x > 0}$.
	Tìm tất cả giá trị thực  của $a$ để hàm số đã cho liên tục trên $\mathbb{R}$.
	 \choice
		 {$a = 1$}
		{$a =  3$}
		{$a =  4$}
		{\True $a = 2$} 
	 \loigiai{
		Tập xác  định của hàm số là $\mathscr{D} = \mathbb{R}$.\\
		Xét trên $(-\infty; 0)$, ta có $f(x) = 3x + a  - 1$. Hàm này liên tục trên $(-\infty; 0)$.\\
		Xét  trên $(0; +\infty)$, ta có $f(x) = \dfrac{\sqrt{1 + 2x} - 1}{x}$. Hàm này liên tục trên $(0; +\infty)$.\\
		Xét tại $x = 0$, ta có  $f(0) = a - 1$.\\
		 Mặt khác
		\begin{align*}
			\lim\limits_{x \to  0^{-}} f(x) &= a - 1,\\
			\lim\limits_{x \to 0^{+}} f(x)  &= \lim\limits_{x \to 0^{+}} \dfrac{\sqrt{1 + 2x} - 1}{x} = 1.
		\end{align*}
		Để hàm số liên tục tại  $x = 0$ thì
		\[ 
		\lim\limits_{x \to 0^{-}} f(x)  = \lim\limits_{x \to 0^{+}} f(x) = f(0) \Leftrightarrow a - 1 = 1 \Leftrightarrow a = 2.
		\] 
		Vậy  $a = 2$ để hàm số liên tục trên $\mathbb{R}$.
	 }
\end{ex}

% nguồn: Ok_Toan11_TracNghiem_Lan2
\begin{ex}%[1D3H3-4]%[ma:3fdef785]
	Tìm tất cả các giá  trị thực của tham số $m$ để hàm số $f(x)=\heva{
	&\dfrac{x^2-16}{x-4} & \text{nếu }  x > 4 \\
	&mx + 1 & \text{nếu } x  \le 4
	}$ liên tục trên  $\mathbb{R}$.
	\choice 
		 {$m=8$ hoặc $m=-\dfrac{7}{4}$}
		{$m=\dfrac{7}{4}$} 
		 {$m=-\dfrac{7}{4}$}
		 {\True $m=-8$ hoặc $m=\dfrac{7}{4}$}
	\loigiai{ 
		 Tập xác định của hàm số là $\mathscr{D} = \mathbb{R}$.\\
		Khi $x > 4$, ta có $f(x) = \dfrac{x^2 - 16}{x - 4}$ là hàm phân thức, nên liên tục trên tập xác định của nó $\Rightarrow  f(x)$ liên tục trên $\left(4; +\infty \right)$.\\
		Khi $x \le 4$, ta có $f(x) = mx + 1$ là hàm đa thức, nên liên tục  trên $\mathbb{R} \Rightarrow f(x)$ liên tục trên $\left(-\infty; 4\right)$.\\
		Do đó, hàm số $f(x)$ liên tục trên $\mathbb{R}$ khi và  chỉ khi hàm số liên tục tại $x = 4$.\\
		 Ta có
		\begin{align*}
			\lim\limits_{x \to 4^{+}} f(x) &= \lim\limits_{x \to 4^{+}} \dfrac{x^2 - 16}{x - 4} = \lim\limits_{x \to 4^{+}} (x  + 4) = 8, \\
			\lim\limits_{x \to 4^{-}} f(x) &= \lim\limits_{x \to 4^{-}} (mx + 1) = 4m + 1. 
		\end{align*}
		Hàm số liên tục tại $x  = 4$ khi
		 \[
		\lim\limits_{x \to 4^{+}} f(x) = \lim\limits_{x \to 4^{-}} f(x) = f(4) \Leftrightarrow 8 = 4m + 1 \Leftrightarrow m  = \dfrac{7}{4}.
		\] 
		Vậy hàm số liên tục  trên $\mathbb{R}$ khi và chỉ khi $m = \dfrac{7}{4}$.
	 }
\end{ex}

% nguồn: Ok_Toan11_TracNghiem_Lan2
\begin{ex}%[1D3H3-4]%[ma:fd9fb419]
	Hàm số nào dưới  đây liên tục trên khoảng $(0; 5)$?
	 \choice
		 {$y=\dfrac{3x-2}{x-3}$}
		{\True $y=\dfrac{x+1}{x+2}$} 
		 {$y=\dfrac{5x+1}{x-4}$}
		 {$y=\dfrac{1}{x^2-1}$}
	\loigiai{Hàm số $y=\dfrac{x+1}{x+2}$ có tập xác định là $\mathbb{R} \setminus\{-2\}$ nên liên tục trên mỗi khoảng  $(-\infty;-2)$ và $(-2;+\infty)$. Vậy hàm số đó liên tục trên khoảng $(0; 5)$.
	} 
\end{ex}

% nguồn: Ok_Toan11_TracNghiem_Lan2
\begin{ex}%[1D3H3-4]%[ma:90fd6296]
	Hàm số $f(x)=\heva{ 
	&ax  + b + 1 & \text{nếu } x > 0 \\
	&a\cos x  + b\sin x & \text{nếu } x \le 0
	}$ liên tục trên $\mathbb{R}$ khi và chỉ khi 
	 \choice
		{\True $a - b = 1$} 
		 {$a - b = -1$}
		 {$a + b = 1$}
		{$a  + b = -1$}
	\loigiai{ 
		Tập xác định của hàm số là $\mathscr{D}  = \mathbb{R}$.\\
		Xét trên $(0; +\infty)$,  ta có $f(x) = ax + b + 1$ liên tục trên $(0; +\infty)$.\\
		Xét trên $(-\infty; 0)$, ta có $f(x) =  a\cos x + b\sin x$ liên tục trên $(-\infty; 0)$.\\
		Xét tại $x = 0$, ta có 
		\begin{align*}
			f(0) &= a,\\ 
			\lim\limits_{x \to 0^{+}} f(x) &= \lim\limits_{x \to 0^{+}} (ax +  b + 1) = b + 1,\\
			\lim\limits_{x \to 0^{-}} f(x) &=  \lim\limits_{x \to 0^{-}} (a\cos x + b\sin x) = a.
		\end{align*}
		Để hàm số liên tục tại $x = 0$, ta  cần
		\[ 
		\lim\limits_{x \to 0^{+}} f(x) = \lim\limits_{x \to 0^{-}} f(x) = f(0) \Leftrightarrow b + 1 = a \Leftrightarrow a -  b = 1.
		\] 
		Vậy hàm số liên tục trên $\mathbb{R}$ khi và chỉ khi  $a - b = 1$.
	 }
\end{ex}

% nguồn: Ok_Toan11_TracNghiem_Lan2
\begin{ex}%[1D3H3-4]%[ma:a61aa39e]
	Trong các hàm số sau, hàm  số nào liên tục trên $\mathbb{R}$?
	 \choice
		{$f(x)=\tan  x+5$}
		 {$f(x)=\dfrac{x^2+3}{5-x}$}
		 {$f(x)=\sqrt{x-6}$}
		{\True  $f(x)=\dfrac{x+5}{x^2+4}$}
	\loigiai{Hàm số $f(x)=\dfrac{x+5}{x^2+4}$ là hàm phân thức hữu tỉ và có TXĐ là $\mathscr{D}=\mathbb{R}$ do đó hàm số $f(x)=\dfrac{x+5}{x^2+4}$ liên tục  trên $\mathbb{R}$.
	 }
\end{ex}

% nguồn: Ok_Toan11_TracNghiem_Lan2
\begin{ex}%[1D3H3-4]%[ma:12507a6f]
	Cho bốn hàm số $f_1(x)=2x^3-3x+1$,  $f_2(x)=\dfrac{3x+1}{x-2}$, $f_3(x)=\cos x+3$ và $f_4(x)=\log_3 x$. Hỏi có bao nhiêu hàm số liên tục trên tập $\mathbb{R}$?
	\choice 
		{$1$} 
		 {$3$}
		 {$4$}
		 {\True $2$}
	\loigiai{Ta có hai hàm số $f_2(x)=\dfrac{3x+1}{x-2}$  và $f_4(x)=\log_3 x$ có tập xác định không phải là $\mathbb{R}$ nên không thỏa yêu cầu.\\
		Cả hai hàm số $f_1(x)=2x^3-3x+1$ và $f_3(x)=\cos x+3$ đều có tập xác định là  $\mathbb{R}$ đồng thời liên tục trên $\mathbb{R}$.
	 }
\end{ex}

% nguồn: Ok_Toan11_TracNghiem_Lan2
\begin{ex}%[1D3H3-4]%[ma:cc20094b]
	Cho  hàm số $y=\heva{
	&3x +  1 & \text{nếu } x \ge -1 \\
	&x + m &  \text{nếu } x < -1,
	}$ $m$ là tham số. Tìm $m$ để hàm số liên tục  trên $\mathbb{R}$.
	 \choice
		{$m  = 5$}
		 {\True $m = -1$}
		{$m =  3$}
		 {$m = -3$}
	\loigiai{ 
		Tập xác định của hàm số là $\mathscr{D}  = \mathbb{R}$.\\
		Xét trên $(-\infty; -1)$, ta có $y = x  + m$ liên tục trên $(-\infty; -1)$.\\
		Xét trên $(-1; +\infty)$, ta có  $y = 3x + 1$ liên tục trên $(-1; +\infty)$.\\
		Xét tại $x =  -1$, ta có
		\begin{align*}
			y(-1) &= 3(-1) + 1 =  -2,\\
			 \lim\limits_{x \to -1^{+}} y &= 3(-1) + 1 = -2,\\
			\lim\limits_{x \to -1^{-}} y  &= (-1) + m.
		\end{align*}
		Hàm số liên tục tại $x = -1$ khi và chỉ  khi
		 \[
		y(-1) = \lim\limits_{x \to -1^{+}} y = \lim\limits_{x \to  -1^{-}} y \Leftrightarrow -2 = -1 + m \Leftrightarrow m = -1.
		 \]
		Vậy hàm số liên tục trên $\mathbb{R}$ khi và chỉ  khi $m = -1$.
	 }
\end{ex}

% nguồn: Ok_Toan11_TracNghiem_Lan2
\begin{ex}%[1D3H3-4]%[ma:9a5dedb3]
	Hàm số nào sau  đây liên tục trên $\mathbb{R}$?
	 \choice
		 {$f(x)=\sqrt{x}$}
		{\True  $f(x)=x^4-4x^2$}
		{$f(x)=\sqrt{\dfrac{x^4-4x^2}{x+1}}$} 
		{$f(x)=\dfrac{x^4-4x^2}{x+1}$} 
	\loigiai{Vì hàm số $f(x)=x^4-4x^2$ có dạng đa thức nên  hàm số này liên tục trên $\mathbb{R}$.
	 }
\end{ex}

% nguồn: Ok_Toan11_TracNghiem_Lan2
\begin{ex}%[1D3H3-4]%[ma:51462663]
	Cho hàm số  $f(x)=\heva{& \dfrac{x^2-1}{x-1} &\text{ khi } x\ne 1 \\
	&m-2&\text{ khi }  x=1}$. Tìm $m$ để hàm số $f(x)$ liên tục trên $\mathbb{R}$.
	\choice 
		{$m=1$} 
		 {$m=2$}
		{\True  $m=4$}
		 {$m=-4$}
	\loigiai{ 
		Do $\lim\limits_{x\to  1} f(x)=\lim\limits_{x\to 1} \dfrac{x^2-1}{x-1}=\lim\limits_{x\to 1} \left(x+1\right)=2$ nên hàm số liên tục tại $x=1$ khi
		\[\lim\limits_{x\to  1} f(x)=f(1)\Leftrightarrow m-2=2\Leftrightarrow m=4.\]
		Khi đó  hàm số liên tục trên $\mathbb{R}$.
	} 
\end{ex}

% nguồn: Ok_Toan11_TracNghiem_Lan2
\begin{ex}%[1D3H3-4]%[ma:b0f7e748]
	Cho hàm số $y=\heva{& -x^2+x+3 &\quad& \text{ khi } x\ge 2  \\ & 5x+2 &\quad& \text{ khi } x<2}$. Chọn mệnh đề sai trong các mệnh đề sau.
	\choice 
		{Hàm số liên tục  tại $x_0=1$}
		{\True  Hàm số liên tục trên $\mathbb{R}$}
		{Hàm số liên tục trên các khoảng $(-\infty;2)$  và $(2;+\infty)$}
		{Hàm số  gián đoạn tại $x_0=2$}
	 \loigiai{
		\begin{itemize}
			\item Với $x>2$, ta có $f(x)=-x^2+x+3$ là hàm  đa thức $\Rightarrow$ hàm số $f(x)$ liên tục trên khoảng $(2;+\infty)$.
			 \item Với $x<2$, ta có $f(x)=5x+2$ là hàm đa thức $\Rightarrow$ hàm số $f(x)$ liên tục trên khoảng $(-\infty;2)$.
			 \item Tại $x=2$.
			\begin{itemize}
				\item  $\lim\limits_{x\to 2^+}f(x)=\lim\limits_{x\to 2^+}(-x^2+x+3)=1$;
				\item $\lim\limits_{x\to 2^-}f(x)=\lim\limits_{x\to  2^-}(5x+2)=12$;
			\end{itemize}
			Suy ra $\lim\limits_{x\to 2^+}f(x)\ne \lim\limits_{x\to 2^-}f(x) \Rightarrow $ không tồn  tại $\lim\limits_{x\to 2}f(x)\Rightarrow $ hàm số gián đoạn tại $x_0=2$.
		\end{itemize}
	} 
\end{ex}

% nguồn: Ok_Toan11_TracNghiem_Lan2
\begin{ex}%[1D3H3-4]%[ma:9805b4b1]
	Cho  hàm số $f(x)=\dfrac{x+5}{x^2-3x+2}$. Khi đó hàm số $y=f(x)$ liên tục trên các khoảng nào sau đây?
	\choice 
		{\True $(1;2)$} 
		{$(-\infty;2)$} 
		{$(1;+\infty)$} 
		{$(1;3)$} 
	\loigiai{ 
		Hàm số có nghĩa khi $x^2-3x+2 \ne 0\Leftrightarrow \heva{&  x\ne 2 \\ & x\ne 1.}$\\
		Vậy theo định lí ta có hàm số $f(x)=\dfrac{x+5}{x^2+3x+2}$ liên tục trên  khoảng $(-\infty;1)$; $(1;2)$ và $(2;+\infty)$.
	} 
\end{ex}

% nguồn: Ok_Toan11_TracNghiem_Lan2
\begin{ex}%[1D3H3-4]%[ma:96f1aa0a]
	Có bao nhiêu giá trị thực của tham số $m$ để hàm số $f(x)=\heva{ 
	&m^2x^2 & \text{nếu } x \le 2 \\ 
	&(1 - m)x & \text{nếu } x > 2}$ liên tục  trên $\mathbb{R}$?
	 \choice
		{$0$} 
		{\True $2$} 
		 {$3$}
		 {$4$}
	 \loigiai{
		Tập xác định của hàm số là $\mathscr{D} = \mathbb{R}$.\\ 
		Xét  trên $(-\infty; 2)$, ta có $f(x) = m^2x^2$ liên tục trên $(-\infty; 2)$.\\
		Xét trên $(2; +\infty)$, ta có $f(x) = (1  - m)x$ liên tục trên $(2; +\infty)$.\\
		Xét tại $x = 2$, ta có $\lim\limits_{x \to 2^{+}} f(x) = (1 - m)2$, $\lim\limits_{x \to  2^{-}} f(x) = 4m^2$, $f(2) = 4m^2$.\\
		Hàm số  liên tục tại $x = 2$ khi và chỉ khi
		 \[
		4m^2 = (1 - m)2 \Leftrightarrow 4m^2 + 2m - 2  = 0 \Leftrightarrow \hoac{&m = \dfrac{1}{2} \\&m = -1.}
		 \]
		Vậy  có hai giá trị của $m$.
	 }
\end{ex}

% nguồn: Ok_Toan11_TracNghiem_Lan2
\begin{ex}%[1D3H3-4]%[ma:2bbd6f05]
	Hàm số nào sau  đây không liên tục trên $\mathbb{R}$?
	\choice 
		{$y=x^3+1$} 
		 {$y=x^4-2x^2+1$}
		{$y=\sin  x$}
		{\True $y=\dfrac{3x-4}{x-2}$} 
	 \loigiai{
		Do các hàm số $y=x^3+1$, $y=x^4-2x^2+1$  là các hàm đa thức có tập xác định $\mathscr{D}=\mathbb{R}$ nên đều liên tục trên $\mathbb{R}$.\\
		Hàm số ở đáp án $y=\sin x$ là hàm  số lượng giác có tập xác định $\mathscr{D}=\mathbb{R}$ nên cũng liên tục trên $\mathbb{R}$.\\
		Xét hàm số $y=\dfrac{3x-4}{x-2}$ có tập xác định $\mathscr{D}=\mathbb{R}\setminus \left\{ 2  \right\}$.\\
		Ta có $\lim\limits_{x\to 2^+}\dfrac{3x-4}{x-2}=+\infty $ và $\lim\limits_{x\to 2^-}\dfrac{3x-4}{x-2}=-\infty $ nên hàm số gián đoạn tại  $x=2$.\\
		Vậy hàm số $y=\dfrac{3x-4}{x-2}$ không liên tục  trên $\mathbb{R}$.
	 }
\end{ex}

% nguồn: Ok_Toan11_TracNghiem_Lan2
\begin{ex}%[1D3H3-4]%[ma:0b22af53]
	Cho hàm số $y=\dfrac{x^2+1}{x^2+5x+4}$. Khi đó, hàm số liên tục trên khoảng  nào dưới đây?
	\choice 
		{$(-3; 2)$} 
		{$(-\infty; 3)$} 
		{$(-5;  3)$}
		{\True  $(-1;+\infty)$}
	 \loigiai{Hàm số xác định khi và chỉ khi $x^2+5x+4\neq 0\Leftrightarrow\heva{&x \neq-1\\
		&x  \neq-4.}$\\
		Tập xác định của làm số là $\mathscr{D}=(-\infty;-4) \cup(-4;-1) \cup(-1;+\infty)$. Do đó,  hàm số $y=\dfrac{x^2+1}{x^2+5x+4}$ liên tục trên $(-\infty;-4)$, $(-4;-1)$ và $(-1;+\infty)$.
	 }
\end{ex}

% nguồn: Ok_Toan11_TracNghiem_Lan2
\begin{ex}%[1D3H3-4]%[ma:e8bf0089]
	Hàm số nào dưới  đây liên tục trên $\mathbb{R}$?
	\choice 
		{\True $y=x+\cos  x$}
		{$y=x-\tan x$} 
		{$y=1+\cot  x$}
		{$y=\dfrac{1}{\cos  x}$}
	\loigiai{Hàm số $y=x+\cos x$ có tập xác định là $\mathbb{R}$ nên liên tục  trên $\mathbb{R}$.
	 }
\end{ex}

% nguồn: Ok_Toan11_TracNghiem_Lan2
\begin{ex}%[1D3H3-4]%[ma:c4adfa6e]
	Cho hàm số $f(x)=\heva{& \sqrt{x}-m &\text{ khi } x\ge 0 \\& mx+1&\text{ khi } x < 0}$. Tìm tất cả các giá  trị của $m$ để $f(x)$ liên tục trên $\mathbb{R}$.
	\choice 
		 {$m=1$}
		 {$m=0$}
		{\True  $m=-1$}
		{$m=-2$} 
	\loigiai{ 
		Để hàm số $f(x)$ liên  tục trên $\mathbb{R} $ thì $ f(x)$ liên tục tại $x=0$.
		 \[\lim\limits_{x\to 0^{+}}f(x)=\lim\limits_{x\to 0^{+}}\left(\sqrt{x}-m\right)=-m; \lim\limits_{x\to 0^{-}}f(x)=\lim\limits_{x\to 0^{-}}\left(mx+1\right)=1; f(0)=-m.\]
		Hàm số $f(x)$ liên tục tại $x=0$ $\Leftrightarrow$ $\lim\limits_{x\to 0^{+}}f(x)=\lim\limits_{x\to 0^{-}}f(x)=f(0)\Leftrightarrow-m=1\Leftrightarrow  m=-1$.
	 }
\end{ex}

% nguồn: Ok_Toan11_TracNghiem_Lan2
\begin{ex}%[1D3H3-4]%[ma:ed996312]
	Hàm số nào trong các hàm số dưới đây không liên tục trên  $\mathbb{R}$?
	 \choice
		{$y=|x|$} 
		 {\True $y=\dfrac{x}{x-1}$}
		{$y=\sin x$} 
		{$y=\dfrac{x}{|x|+1}$} 
	 \loigiai{
		Tập xác định của  hàm số $y=\dfrac{x}{x-1}$ là $\mathscr{D} = \mathbb{R}\setminus\left\{ 1 \right\}$.\\
		Hàm số liên tục trên từng khoảng $(-\infty;1)$ và $(1;+\infty)$ nên hàm số không liên tục trên  $\mathbb{R}$.
	} 
\end{ex}

% nguồn: Ok_Toan11_TracNghiem_Lan2
\begin{ex}%[1D3H3-4]%[ma:b5982e1b]
	 Hàm số nào dưới đây liên tục trên $\mathbb{R}$?
	\choice 
		{\True  $y=2x-3\cos x$}
		{$y=1+\tan x$} 
		{$y=x-\cot x$} 
		{$y=\dfrac{1}{\cos x}$} 
	 \loigiai{Do hàm số $y=1+\tan x$ và hàm số $y=\dfrac{1}{\cos x}$ có tập xác định là
		$\mathscr{D}=\mathbb{R} \setminus\left\{\dfrac{\pi}{2}+k \pi,k  \in \mathbb{Z}\right\}$ nên hàm số không liên tục trên $\mathbb{R}$.\\
		Do hàm số $y=x-\cot x$ có tập xác định là  $D=\mathbb{R} \setminus\{k \pi$, $k \in \mathbb{Z}\}$ nên hàm số không liên tục trên $\mathbb{R}$.\\
		Do hàm số $y=2x-3\cos x$ là hàm số sơ cấp có tập xác định là $\mathbb{R}$ nên liên  tục trên $\mathbb{R}$.
	} 
\end{ex}

% nguồn: Ok_Toan11_TracNghiem_Lan2
\begin{ex}%[1D3H3-4]%[ma:d964a5e9]
	Cho hàm số $f(x)=\heva{ 
	&\dfrac{x^2-1}{x-1} & \text{nếu } x \neq  1 \\
	&m - 2 & \text{nếu } x = 1.}$\\ 
	Tìm $m$ để hàm số  $f(x)$ liên tục trên $\mathbb{R}$.
	\choice 
		{$m =  1$}
		{$m  = 2$}
		{\True $m = 4$} 
		{$m =  -4$}
	 \loigiai{
		Tập xác định của hàm số là $\mathscr{D} = \mathbb{R}$.\\ 
		 Nếu $x \neq 1$, ta có
		 \[
		f(x) = \dfrac{x^2 - 1}{x - 1} = \dfrac{(x - 1)(x + 1)}{x - 1}  = x + 1, \quad \text{với} \, x \neq 1.
		\] 
		Hàm số này xác định và liên tục trên mỗi  khoảng $(-\infty; 1)$ và $(1; +\infty)$.\\
		 Tại $x = 1$, ta có $f(1) = m - 2$.\\
		Mặt khác,  ta tính giới hạn của hàm số tại $x = 1$:
		 \[
		\lim\limits_{x \to 1} f(x)  = \lim\limits_{x \to 1} (x + 1) = 2.
		 \]
		Hàm số liên tục tại $x = 1$ khi và  chỉ khi:
		\[ 
		\lim\limits_{x \to  1} f(x) = f(1) \Leftrightarrow 2 = m - 2 \Leftrightarrow m = 4.
		\] 
		Vậy hàm số  liên tục trên $\mathbb{R}$ khi và chỉ khi $m = 4$.
	 }
\end{ex}

% nguồn: Ok_Toan11_TracNghiem_Lan2
\begin{ex}%[1D3H3-4]%[ma:449ee57b]
	Hàm  số nào trong các hàm số dưới đây không liên tục trên $\mathbb{R}$?
	 \choice
		 {$y=x^2+6x+20$}
		 {$y=\cos x$}
		 {$y=\dfrac{x}{x^2+x+2}$}
		{\True  $y=\dfrac{x}{x+1}$}
	\loigiai{Hàm số $y=\dfrac{x}{x+1}$ không  xác định tại điểm $x=-1$ nên không liên tục tại điểm đó.
	} 
\end{ex}

% nguồn: Ok_TracNghiemLop11_Lan1
\begin{ex}%[1D3H3-4]%[ma:c9cf4009]
	Cho hàm số $f(x)=\heva{&x&\text{khi}&\,\,x<1,x\ne 0\\&0&\text{khi}&\,\,x=0\\&\sqrt{x}&\text{khi}&\,\,x\ge  1}$. Hàm số $f(x)$ liên tục tại:
	 \choice
		{\True mọi điểm thuộc  $\mathbb{R}$}
		{mọi  điểm trừ $x=1$}
		{mọi  điểm trừ $x=0$}
		 {mọi điểm trừ $x=0$ và $x=1$}
	\loigiai{ 
		Tập xác định  $\mathscr{D}=\mathbb{R}$.\\
		Hàm số  liên tục trên mỗi khoảng $(-\infty;0)$, $(0;1)$ và $(1;+\infty)$.
		\begin{enumerate}[$\bullet$]
			\item Tại $x=0$,  ta có $\heva{&f(0)=0\\&\lim\limits_{x\to 0}f(x)=\lim\limits_{x\to 0} x=0.}$\\
			Vậy hàm số $f(x)$  liên tục tại $x=0$.
			\item Tại $x=1$, ta có $\heva{&f(1)=1\\&\lim\limits_{x\to  1^{-}}f(x)=\lim\limits_{x\to 1^{-}}x=1\\&\lim\limits_{x\to 1^{+}}f(x)=\lim\limits_{x\to 1^{+}}\sqrt{x}=1.}$\\
			Vậy hàm số $f(x)$ liên tục tại  $x=1$.
		\end{enumerate}
	} 
\end{ex}

% nguồn: Ok_TracNghiemLop11_Lan1
\begin{ex}%[1D3H3-4]%[ma:17773124]
	Xét tính liên tục của hàm số $f(x)=\heva{&\dfrac{x-1}{\sqrt{2-x}-1}&\text{khi}&\,\,x<1\\&-2x&\text{khi}&\,\,x\ge 1}$. Khẳng định nào  dưới đây \textbf{đúng}?
	 \choice
		 {$f(x)$ không liên tục trên khoảng $(-\infty;1)$}
		{$f(x)$ không liên tục  trên khoảng $(1;+\infty)$}
		{$f(x)$ gián đoạn tại $x=1$} 
		{\True $f(x)$ liên tục trên $\mathbb{R}$} 
	\loigiai{ 
		Tập xác định $\mathscr{D}=(-\infty;2]$.\\ 
		Xét $x\in (-\infty;1)$, khi đó hàm số $f(x)=\dfrac{x-1}{\sqrt{2-x}-1}$ xác định và liên tục trên khoảng $(-\infty;1)$.\\ 
		Xét $x\in (1;+\infty)$,  khi đó hàm số $f(x)=-2x$ xác định và liên tục trên $(1;+\infty)$.\\
		 Ta có $\heva{&f(1)=-2\\&\lim\limits_{x\to 1^{+}}f(x)=\lim\limits_{x\to 1^{+}}(-2x)=-2\\&\lim\limits_{x\to 1^{-}}f(x)=\lim\limits_{x\to 1^{-}}\dfrac{x-1}{\sqrt{2-x}-1}=\lim\limits_{x\to 1^{-}}\left[-\left(\sqrt{2-x}+1 \right)\right]=-2.}$\\
		 Vậy hàm số $f(x)$ liên trục trên $\mathbb{R}$.
	} 
\end{ex}

% nguồn: Ok_TracNghiemLop11_Lan1
\begin{ex}%[1D3H3-4]%[ma:1cb7ddb2]
	Cho hàm số $f(x)=\heva{&\dfrac{x^2-1}{x-1}&\text{khi}&\,\,x<3,x\ne  1\\&4&\text{khi}&\,\,x=1\\&\sqrt{x+1}&\text{khi}&\,\,x\ge 3}$ . Hàm số $f(x)$ liên tục tại:
	\choice 
		{mọi điểm thuộc  $\mathbb{R}$}
		{mọi  điểm trừ $x=1$}
		{mọi điểm trừ $x=3$} 
		{\True mọi  điểm trừ $x=1$ và $x=3$}
	\loigiai{ 
		 Tập xác định $\mathscr{D}=\mathbb{R}$.\\
		Hàm  số liên tục trên mỗi khoảng $(-\infty;1)$, $(1;3)$ và $(3;+\infty)$.
		\begin{enumerate}[$\bullet$]
			\item Tại $x=1$, ta  có $\heva{&f(1)=4\\&\lim\limits_{x\to 1}f(x)=\lim\limits_{x\to 1}\dfrac{x^2-1}{x-1}=\lim\limits_{x\to 1}(x+1)=2.}$\\
			Vậy hàm số $f(x)$ gián đoạn tại  $x=1$.
			\item Tại $x=3$, ta có $\heva{&f(3)=2\\&\lim\limits_{x\to  3^{-}}f(x)=\lim\limits_{x\to 3^{-}}\dfrac{x^2-1}{x-1}=\lim\limits_{x\to 3^{-}}(x+1)=4.}$\\
			Vậy hàm số  $f(x)$ gián đoạn tại $x=3$.
		\end{enumerate}
	} 
\end{ex}

% nguồn: Ok_TracNghiemLop11_Lan1
\begin{ex}%[1D3H3-4]%[ma:cf8084f7]
	Hàm số nào sau đây liên tục  trên $\mathbb{R}$?
	 \choice
		 {$y=\dfrac{1}{x}$}
		{$y=\tan x$} 
		{$y=\cot x$} 
		{\True $y=\sin x$} 
	\loigiai{Vì hàm số $y=\sin  x$ có tập xác định là $\mathbb{R}$ nên nó liên tục trên $\mathbb{R}$.
	 }
\end{ex}

% nguồn: Ok_TracNghiemLop11_Lan1
\begin{ex}%[1D3H3-4]%[ma:7cc2419e]
	 Cho hàm số $f(x)=\heva{&\dfrac{\sqrt{x}-2}{x-4}&\text{khi}\,\,x\ne 4\\&\dfrac{1}{4}&\text{khi}\,\,x=4}$. Khẳng định nào sau đây \textbf{đúng}?
	\choice 
		{Hàm số có tập xác định $\mathscr{D}=\mathbb{R}\setminus\{4\}$} 
		{Hàm số liên tục tại mọi điểm trên tập xác định nhưng gián đoạn tại $x=4$} 
		{Hàm số không liên  tục tại $x=4$}
		{\True Hàm  số liên tục tại $x=4$}
	\loigiai{ 
		 Tập xác định $\mathscr{D}=[0;+\infty)$.\\
		Ta có $\heva{&f(4)=\dfrac{1}{4}\\&\lim\limits_{x\to  4} f(x)=\lim\limits_{x\to 4} \dfrac{1}{\sqrt{x}+2}=\dfrac{1}{4}.}$\\
		Vậy hàm số $f(x)$ liên  tục tại $x=4$.
	 }
\end{ex}

% nguồn: Ok_TracNghiemLop11_Lan1
\begin{ex}%[1D3H3-4]%[ma:de4f7ee4]
	Biết rằng $f(x)=\heva{&\dfrac{x^2-1}{\sqrt{x}-1}&\text{khi}&\,\,x\ne 1\\&a&\text{khi}&\,\,x=1}$ liên tục trên đoạn $[0;1]$ (với $a$ là tham số). Khẳng định nào dưới đây về giá trị  $a$ là \textbf{đúng}?
	\choice 
		{\True $a$ là một số  nguyên}
		{$a$ là một số vô tỉ} 
		 {$a>5$}
		 {$a<0$}
	\loigiai{ 
		Hàm số xác định và liên tục  trên $[0;1)$.\\
		 Ta có $\heva{&f(1)=a\\&\lim\limits_{x\to 1^{-}}f(x)=\lim\limits_{x\to 1^{-}}\dfrac{x^2-1}{\sqrt{x}-1}=\lim\limits_{x\to 1^{-}}\left[(x+1)\left(\sqrt{x}+1 \right)\right]=4.}$\\
		Hàm số $f(x)$ liên tục  trên $[0;1]$ khi và chỉ khi $\lim\limits_{x\to 1^{-}}f(x)=f(1)\Leftrightarrow a=4$.
	 }
\end{ex}

% nguồn: Ok_TracNghiemLop11_Lan1
\begin{ex}%[1D3H3-4]%[ma:a334124e]
	Xét tính liên tục của hàm số $f(x)=\heva{&1-\cos x&\text{khi}&\,\,x\le 0\\&\sqrt{x+1}&\text{khi}&\,\,x>0}$. Khẳng định nào  sau đây \textbf{đúng}?
	\choice 
		{$f(x)$ liên  tục tại $x=0$}
		{$f(x)$ liên tục trên $(-\infty;1)$} 
		{\True  $f(x)$ không liên tục trên $\mathbb{R}$}
		{$f(x)$  gián đoạn tại $x=1$}
	 \loigiai{
		Tập xác định $\mathscr{D}=\mathbb{R}$.\\ 
		Ta có $f(x)$  liên tục trên $(-\infty;0)$ và $(0;+\infty)$,\\
		Mặt khác $\heva{&f(0)=0\\&\lim\limits_{x\to 0^{-}}f(x)=\lim\limits_{x\to 0^{-}}(1-\cos x)=0\\&\lim\limits_{x\to  0^{+}}f(x)=\lim\limits_{x\to 0^{+}}\sqrt{x+1}=1}.$\\
		Vậy hàm số $f(x)$  không liên tục tại $x=0$ nên $f(x)$ không liên tục trên $\mathbb{R}$.
	 }
\end{ex}

% nguồn: Ok_Toan11_TracNghiem_Lan2
\begin{ex}%[1D3V3-4]%[ma:7e04ed77]
	Cho biết  hàm số $f(x)=\heva{
	&\dfrac{x^3 - 3x^2 + 2x}{x(x - 2)}  && \text{nếu } x(x - 2) \neq 0 \\
	&a && \text{nếu  } x = 0 \\
	&b && \text{nếu } x = 2}$ liên tục  trên $\mathbb{R}$.
	Tính $T = a^2 + b^2$. 
	\choice 
		{\True  $T = 2$}
		{$T =  122$}
		{$T =  101$}
		 {$T = 145$}
	\loigiai{ 
		Vì hàm số $f(x)$ liên tục trên $\mathbb{R}$ nên hàm số cũng liên tục tại $x =  0$ và $x = 2$.\\
		Xét tại $x  = 0$, ta có
		\begin{eqnarray*}
			&&\lim\limits_{x \to 0} f(x) = \lim\limits_{x  \to 0} \dfrac{x^3 - 3x^2 + 2x}{x(x - 2)} = \lim\limits_{x \to 0} \dfrac{x(x - 1)(x - 2)}{x(x - 2)} = f(0).\\
			&\Rightarrow& \lim\limits_{x \to 0} \dfrac{(x - 1)(x - 2)}{x - 2} = a \Rightarrow  a = -1.
		\end{eqnarray*}
		Xét tại $x =  2$, ta có
		\begin{eqnarray*}
			&&\lim\limits_{x \to 2} f(x) = \lim\limits_{x \to 2} \dfrac{x^3 - 3x^2 + 2x}{x(x - 2)} = \lim\limits_{x \to 2}  \dfrac{x(x - 1)(x - 2)}{x(x - 2)} = f(2).\\
			&\Rightarrow& \lim\limits_{x \to 2} \dfrac{x(x - 1)}{x} = b \Rightarrow  b = 1.
		\end{eqnarray*}
		Vậy $T = a^2 + b^2 = (-1)^2 + 1^2 = 1 +  1 = 2$.
	 }
\end{ex}

% nguồn: Ok_Toan11_TracNghiem_Lan2
\begin{ex}%[1D3V3-4]%[ma:88e15042]
	Tìm $m$ để hàm  số $y=\heva{&\dfrac{2\sqrt[3]x-x-1}{x-1}, &x \neq 1\\
	&m x+1, &x=1}$ liên  tục trên $\mathbb{R}$.
	 \choice
		 {\True $m=-\dfrac{4}{3}$}
		{$m=-\dfrac{1}{3}$} 
		 {$m=\dfrac{4}{3}$}
		 {$m=\dfrac{2}{3}$}
	\loigiai{Xét $x \neq 1$, hàm số $y=\dfrac{2\sqrt[3]x-x-1}{x-1}$ liên  tục trên khoảng $(-\infty; 1)$ và $(1;+\infty)$.\\
		 Xét $x=1$, ta có $y(1)=m+1$ và
		\[\lim\limits_{x \to 1} y=\lim\limits_{x \to 1} \dfrac{2\sqrt[3]{x}-x-1}{x-1}=\lim\limits_{x \to 1} \dfrac{2(\sqrt[3]{x}-1)-(x-1)}{x-1}=\lim\limits_{x \to 1} \dfrac{2}{\sqrt[3]{x^2}+\sqrt[3]{x}+1}-1=\dfrac{2}{3}-1=-\dfrac{1}{3}.\] 
		Hàm số liên tục tại $x=1$ khi và chỉ  khi $\lim\limits_{x \to 1} y=y(1) \Leftrightarrow m+1=-\dfrac{1}{3} \Leftrightarrow m=-\dfrac{4}{3}$.\\
		Vậy với $m=-\dfrac{4}{3}$ thì hàm số liên tục trên  $\mathbb{R}$.
	} 
\end{ex}

% nguồn: Ok_Toan11_TracNghiem_Lan2
\begin{ex}%[1D3V3-4]%[ma:d16f91f5]
	Cho biết hàm số $f(x)=\heva{&\dfrac{x^3-3x^2+2x}{x\left(x-2\right)} &\text{ khi }& x\left(x-2\right)\ne 0 \\ &a &\text{  khi }& {x=0} \\
	 & b &\text{ khi }& x=2} $ liên tục trên $\mathbb{R}$.
	Tính $T=a^2+b^2$. 
	\choice 
		 {\True $T=2$}
		{$T=122$} 
		{$T=101$} 
		 {$T=145$}
	 \loigiai{
		Vì hàm số $f(x)$ liên tục trên $\mathbb{R}$ suy ra hàm  số cũng liên tục tại $x=0$ và $x=2$. Do đó
		\begin{eqnarray*}
			&&\lim\limits_{x\to 0} f(x)=\lim\limits_{x\to 0}  \dfrac{x^3-3x^2+2x}{x\left(x-2\right)}=\lim\limits_{x\to 0} \dfrac{x\left(x-1\right)\left(x-2\right)}{x\left(x-2\right)}=f(0)\Leftrightarrow \lim\limits_{x\to 0} \dfrac{\left(x-1\right)\left(x-2\right)}{x-2}=a\Leftrightarrow a=-1.\\
			&&\lim\limits_{x\to 2} f(x)=\lim\limits_{x\to 2} \dfrac{x^3-3x^2+2x}{x\left(x-2\right)}=\lim\limits_{x\to  2} \dfrac{x\left(x-1\right)\left(x-2\right)}{x\left(x-2\right)}=f(2)\Leftrightarrow \lim\limits_{x\to 2} \dfrac{x\left(x-1\right)}{x}=b\Leftrightarrow b=1.
		\end{eqnarray*}
		 Vậy $T=a^2+b^2=1+1=2$.
	} 
\end{ex}

% nguồn: Ok_Toan11_TracNghiem_Lan2
\begin{ex}%[1D3V3-4]%[ma:3e86c53a]
	Tìm $m$ để hàm số $f(x)=\heva{&\dfrac{x^2+4x+3}{x+1} &\text{khi }  x >-1\\
	&m x+2  &\text{khi } x \leq-1}$ liên tục tại điểm $x=-1$.
	\choice 
		{$m=2$} 
		 {\True $m=0$}
		 {$m=-4$}
		 {$m=4$}
	\loigiai{Ta có 
		\begin{align*}
			&\lim\limits_{x \to(-1)^{+}} f(x)=\lim\limits_{x \to(-1)^{+}}  \dfrac{x^2+4x+3}{x+1}=\lim\limits_{x \to(-1)^{+}} \dfrac{(x+1)(x+3)}{x+1}=\lim\limits_{x \to(-1)^{+}}(x+3)=2\\
			&\lim\limits_{x \to(-1)^{-}}  f(x)=\lim\limits_{x \to(-1)^{-}}(m x+2)=-m+2\\
			 &f(-1)=-m+2.
		\end{align*}
	 Hàm số đã cho liên tục tại điểm $x=-1$ khi và chỉ khi \[\lim\limits_{x \to(-1)^{+}} f(x)=\lim\limits_{x \to(-1)^{-}} f(x)=f(-1)\Leftrightarrow 2=-m+2\Leftrightarrow m=0.\]}
\end{ex}

% nguồn: Ok_Toan11_TracNghiem_Lan2
\begin{ex}%[1D3V3-4]%[ma:66f38a93]
	Tìm tất  cả các giá trị thực của tham số $m$ để hàm số $f(x)=\heva{& \dfrac{x^2-16}{x-4} &\text{ khi } x > 4 \\
	&mx+1 &\text{ khi } x\le 4}$  liên tục trên $\mathbb{R}$.
	 \choice
		 {$m=8$ hoặc $m=-\dfrac{7}{4}$}
		{\True $m=\dfrac{7}{4}$} 
		 {$m=-\dfrac{7}{4}$}
		{$m=-8$  hoặc $m=\dfrac{7}{4}$}
	\loigiai{ 
		Với $x > 4$ thì $f(x)=\dfrac{x^2-16}{x-4}$ là hàm phân thức nên liên tục trên TXĐ  của nó $\Rightarrow f(x)$ liên tục trên $\left(4;+\infty \right)$.
		Với $x < 4$ thì $f(x)=mx+1$ là hàm đa thức nên liên tục  trên $\mathbb{R}$ $\Rightarrow f(x)$ liên tục trên $\left(-\infty;4\right)$.
		Do vậy hàm số  $f(x)$ đã liên tục trên các khoảng $\left(4;+\infty \right)$, $\left(-\infty;4\right)$.
		Suy  ra: Hàm số $f(x)$ liên tục trên $\mathbb{R}$ $\Leftrightarrow$ $f(x)$ liên tục tại $x=4$.
		\begin{eqnarray*}
			&&\Leftrightarrow {\lim\limits_{x\to 4^{+}}}  f(x)=\lim\limits_{x\to 4^{-}} f(x)=f(4)\\
			&&\Leftrightarrow {\lim\limits_{x\to 4^{+}}} \dfrac{x^2-16}{x-4}=\lim\limits_{x\to  4^{-}} \left(mx+1\right)=4m+1\\
			 &&\Leftrightarrow {\lim\limits_{x\to 4^{+}}} \left(x+4\right)=4m+1\\
			&&\Leftrightarrow 4m+1=8\Leftrightarrow  m=\dfrac{7}{4}.
		\end{eqnarray*}
	} 
\end{ex}

% nguồn: Ok_Toan11_TracNghiem_Lan2
\begin{ex}%[1D3V3-4]%[ma:95ad8237]
	Cho hàm số $f(x)=\heva{&\dfrac{x^3-8}{x-2} &\text{khi  } x \neq 2\\
	&2m+1 &\text{khi  } x=2}$. Tìm $m$ để hàm số liên tục tại điểm $x_0=2$.
	 \choice
		 {$m=\dfrac{3}{2}$}
		 {$m=\dfrac{13}{2}$}
		{\True $m=\dfrac{11}{2}$} 
		 {$m=-\dfrac{1}{2}$}
	 \loigiai{Ta có
		\begin{align*}
			&f(2)=2m+1\\ 
			&\lim\limits_{x \to 2} f(x)=\lim\limits_{x \to 2} \dfrac{x^3-8}{x-2}=\lim\limits_{x \to 2} \dfrac{(x-2)\left(x^2+2x+4\right)}{x-2}=\lim\limits_{x \to 2}\left(x^2+2x+4\right)=12. 
		\end{align*}
		Hàm số liên tục  tại $x_0=2$ khi và chỉ khi $f(2)=\lim\limits_{x \to 2} f(x) \Leftrightarrow 2m+1=12\Leftrightarrow m=\dfrac{11}{2}$.
	} 
\end{ex}

% nguồn: Ok_Toan11_TracNghiem_Lan2
\begin{ex}%[1D3V3-4]%[ma:6b6df8bf]
	Cho hàm số $y=f(x)=\heva{& \dfrac{\sqrt{x^2+16}-5}{x-3} &\text{ khi } x\ne 3  \\ &a &\text{ khi } x=3}$. Tập các giá trị của $a$ để hàm số đã cho liên tục trên $S$ là:
	\choice 
		 {$\left\{\dfrac{2}{5} \right\}$}
		 {$\left\{\dfrac{1}{5} \right\}$}
		{$\{0\}$} 
		{\True  $\left\{\dfrac{3}{5} \right\}$}
	\loigiai{ 
		Tập xác  định $\mathscr{D}= \mathbb{R}$.
		Khi $x\ne 3$ thì $f(x)=\dfrac{\sqrt{x^2+16}-5}{x-3}$ xác định và liên tục trên  các khoảng $\left(-\infty;3\right)$ và $\left(3;+\infty \right)$.
		Khi $x=3$  thì $f(3)=a$ và $\lim\limits_{x\to 3} f(x)=\lim\limits_{x\to 3} \dfrac{\sqrt{x^2+16}-5}{x-3}$ $=\lim\limits_{x\to 3} \dfrac{x+3}{\sqrt{x^2+16}+5}=\dfrac{3}{5}$.
		Hàm số đã cho liên tục trên $ \mathbb{R}$ khi và chỉ khi nó liên tục tại điểm  $x=3\Leftrightarrow a=\dfrac{3}{5}$.
	 }
\end{ex}

% nguồn: Ok_Toan11_TracNghiem_Lan2
\begin{ex}%[1D3V3-4]%[ma:246ae8f6]
	 Cho hàm số $f(x)=\heva{&\dfrac{-x^2+2x+8}{x+2} &\text{khi } x \neq-2\\
	&m^2x^2+5m x &\text{khi } x=-2}(m \in \mathbb{R})$. Biết hàm số $f(x)$ liên tục tại $x_0=-2$. Số giá trị nguyên của $m$ thỏa mãn yêu cầu  bài toán là
	\choice 
		{$3$} 
		{$2$} 
		{\True  $1$}
		 {$0$}
	\loigiai{TXĐ:  $\mathscr{D}=\mathbb{R}$. Ta có $\lim\limits_{x \to-2} f(x)=\lim\limits_{x \to-2} \dfrac{-x^2+2x+8}{x+2}=6$, $f(2)=4m^2-10m$.\\
		Hàm số  liên tục tại $x_0=-2$ khi và chỉ khi
		$4m^2-10m=6\Leftrightarrow  4m^2-10m-6=0\Leftrightarrow\hoac{&m=3\\
		 &m=-\dfrac{1}{2}.
		 }$\\
		Mà $m$ là số nguyên nên $m=3$. Vậy có $1$ giá trị  nguyên của $m$ thỏa mãn bài toán.
	 }
\end{ex}

% nguồn: Ok_Toan11_TracNghiem_Lan2
\begin{ex}%[1D3V3-4]%[ma:73fe8618]
	Có bao nhiêu  giá trị thực của tham số $m$ để hàm số $f(x)=\heva{&m^2x^2 &\text{ khi } x\le 2 \\&\left(1-m\right)x &\text{ khi } x > 2}$ liên tục trên $\mathbb{R}$?
	 \choice
		 {$0$}
		{\True  $2$}
		{$3$} 
		 {$4$}
	 \loigiai{
		 Ta có
		\begin{eqnarray*}
			 &&\lim\limits_{x\to 2^{+}} f(x)=\lim\limits_{x\to 2^{-}} \left(1-m\right)x=\left(1-m\right)2\\
			&&\lim\limits_{x\to 2^{-}} f(x)=\lim\limits_{x\to  2^{-}} (m^2x^2)=4m^2; f(2)=4m^2.
		\end{eqnarray*}
		Hàm số liên tục tại $x=2$ khi và chỉ  khi
		\[\lim\limits_{x\to 2^{+}} f(x)=\lim\limits_{x\to  2^{+}} f(x)=f(2)\Leftrightarrow 4m^2=\left(1-m\right)2\Leftrightarrow 4m^2+2m-2=0\]
		Phương  trình luôn có hai nghiệm thực phân biệt. Vậy có hai giá trị của $m$.
	 }
\end{ex}

% nguồn: Ok_Toan11_TracNghiem_Lan2
\begin{ex}%[1D3V3-4]%[ma:3902287e]
	Tìm $P$ để  hàm số $y=\heva{& \dfrac{x^2-4x+3}{x-1} &\text{ khi } x > 1 \\ & 6Px-3 &\text{ khi } x\le 1}$ liên tục trên $\mathbb{R}$.
	\choice 
		{$P=\dfrac{5}{6}$} 
		{$P=\dfrac{1}{2}$} 
		{\True  $P=\dfrac{1}{6}$}
		 {$P=\dfrac{1}{3}$}
	\loigiai{ 
		Hàm số $y=f(x)$ liên  tục trên $\mathbb{R}$ $\Rightarrow$ $y=f(x)$ liên tục tại $x=1$.
		\begin{eqnarray*}
			 \lim\limits_{x\to 1^{+}} f(x)&=&\lim\limits_{x\to 1^{-}} f(x)=f(1);\\
			\lim\limits_{x\to 1^{+}}  f(x)&=&\lim\limits_{x\to 1^{+}} \dfrac{x^2-4x+3}{x-1}=\lim\limits_{x\to 1^{+}} \left(x-3\right)=-2;\\
			\lim\limits_{x\to 1^{-}} f(x)&=&\lim\limits_{x\to 1^{-}}  \left(6Px-3\right)=6P-3;\\
			f(1)&=&6P-3. 
		\end{eqnarray*}
		Do đó $\lim\limits_{x\to 1^{+}} f(x)=\lim\limits_{x\to 1^{-}} f(x)=f(1)$  $\Leftrightarrow 6P-3=-2\Leftrightarrow P=\dfrac{1}{6}$.
	 }
\end{ex}

% nguồn: Ok_Toan11_TracNghiem_Lan2
\begin{ex}%[1D3V3-4]%[ma:31c126a8]
	Hàm số $f(x)=\heva{ &ax+b+1, &\text{ khi } x > 0  \\ &a\cos x+b\sin x, &\text{ khi } x\le 0}$ liên tục trên $\mathbb{R}$ khi và chỉ khi
	 \choice
		 {\True $a-b=1$}
		{$a-b=-1$} 
		{$a+b=1$} 
		{$a+b=1$} 
	 \loigiai{
		Khi $x < 0$  thì $f(x)=a\cos x+b\sin x$ liên tục với $x < 0$.
		Khi $x > 0$ thì $f(x)=ax+b+1$ liên tục với mọi  $x > 0$.
		Tại  $x=0$ ta có $f(0)=a$.
		\begin{eqnarray*}
			\lim\limits_{x\to  0^{+}}f(x)&=&\lim\limits_{x\to 0^{+}}\left(ax+b+1\right)=b+1.\\
			 \lim\limits_{x\to 0^{-}}f(x)&=&\lim\limits_{x\to 0^{-}}\left(a\cos x+b\sin x\right)=a.
		\end{eqnarray*}
		Để hàm số liên tục tại $x=0$ thì  $\lim\limits_{x\to 0^{+}}f(x)$ $=\lim\limits_{x\to 0^{-}}f(x)$, $=f(0)$ $\Leftrightarrow a=b+1$ $\Leftrightarrow a-b=1$.
	} 
\end{ex}

% nguồn: Ok_Toan11_TracNghiem_Lan2
\begin{ex}%[1D3V3-4]%[ma:f53341db]
	Cho hàm số $f(x)=\heva{&\sin  x &\text{nếu } \cos x \geq 0\\
	&1+\cos x &\text{nếu } \cos x < 0}$. Hỏi hàm số $f$ có tất cả bao nhiêu điểm gián đoạn trên  khoảng $(0; 2018)$?
	 \choice
		 {$2018$}
		{$1009$} 
		 {$642$}
		{\True $321$} 
	\loigiai{ 
		Ta có  $f(x)$ liên tục trên từng khoảng làm cho $\cos x>0$ hoặc $\cos x<0$.\\
		Tại những  điểm làm cho $\cos x=0\Leftrightarrow x=\dfrac{\pi}{2}+k\pi$ $(k\in\mathbb{Z})$, ta có $\lim\limits_{\cos x \to 0^-}f(x)=\lim\limits_{\cos x \to 0^-}(1+\cos x)=1$.\\
		Tuy nhiên, khi $k$ chẵn, ta có khi $\lim\limits_{\cos x \to 0^+}f(x)=\lim\limits_{\cos x \to  0^+}\sin x=1$.\\
		Khi $k$ lẻ,  $\lim\limits_{\cos x \to 0^+}f(x)=\lim\limits_{\cos x \to 0^+}\sin x=-1$.\\
		Do  đó, hàm số sẽ gián đoạn tại những điểm $x=\dfrac{\pi}{2}+k\pi$ ($k\in\mathbb{Z}$, $k$ lẻ).\\
		Vì ta chỉ tìm những điểm gián đoạn trên $(0;2018)$  nên $0<\dfrac{\pi}{2}+k\pi<2018\Leftrightarrow 0\leq k < 641$.\\
		Số số  lẻ từ $0$ đến $641$ là $321$ số.
	} 
\end{ex}

% nguồn: Ok_Toan11_TracNghiem_Lan2
\begin{ex}%[1D3V3-4]%[ma:0caad967]
	Cho  hàm số $f(x)=\heva{&2m x^2-4 &\text{khi } x \leq 3\\
	&5 &\text{khi } x > 3}$ ($m$ là tham số). Tìm giá trị của $m$ để hàm số liên tục trên $\mathbb{R}$. 
	\choice 
		{\True $\dfrac{1}{2}$} 
		 {$\dfrac{1}{18}$}
		 {$18$}
		{$2$} 
	\loigiai{Khi  $x < 3$ thì $f(x)=2m x^2-4$ là hàm sơ cấp nên liên tục trên $(-\infty; 3)$.\\
		Khi $x > 3$ thì $f(x)=5$  là hàm hằng nên liên tục trên $(+\infty; 3)$.\\
		Vậy hàm số  $f(x)$ liên tục trên $\mathbb{R}$ khi và chỉ khi hàm số $f(x)$ liên tục tại $x=3$ hay
		\[\lim\limits_{x \to 3^{-}} f(x)=\lim\limits_{x \to  3^{+}} f(x)=f(3)\Leftrightarrow \lim\limits_{x \to 3^{-}}\left(2m x^2-4\right)=\lim\limits_{x \to 3^{+}} 5=18m-4 \Leftrightarrow 18m-4=5\Leftrightarrow m=\dfrac{1}{2}.\]
	} 
\end{ex}

% nguồn: Ok_Toan11_TracNghiem_Lan2
\begin{ex}%[1D3V3-4]%[ma:9f68f017]
	Cho hàm số $f(x)=\heva{&3x+a-1&\text{ khi } x\le 0 \\&\dfrac{\sqrt{1+2x}-1}{x} &\text{ khi }  x > 0}$. Tìm tất cả giá trị thực của a để hàm số đã cho liên tục trên $\mathbb{R}$.
	 \choice
		 {$a=1$}
		{$a=3$} 
		{$a=4$} 
		 {\True $a=2$}
	\loigiai{ 
		Hàm số liên  tục tại mọi điểm $x\ne 0$ với bất kỳ $a$.
		Với $x=0$ Ta có $f(0)=a-1$; 
		\begin{eqnarray*}
			\lim\limits_{x\to 0^{-}}f(x)&=&\lim\limits_{x\to 0^{-}}(3x+a-1) =a-1;\\ 
			\lim\limits_{x\to 0^{+}}f(x)&=&\lim\limits_{x\to 0^{+}}\dfrac{\sqrt{1+2x}-1}{x} =\lim\limits_{x\to 0^{+}}\dfrac{2x}{x\left(\sqrt{1+2x}+1\right)} =\lim\limits_{x\to  0^{+}}\dfrac{2}{\sqrt{1+2x}+1}=1;
		\end{eqnarray*}
		Hàm số  liên tục trên $\mathbb{R}$ khi và chỉ khi hàm số liên tục tại $x=0\Leftrightarrow a-1=1\Leftrightarrow a=2$.
	} 
\end{ex}

% nguồn: Ok_Toan11_TracNghiem_Lan2
\begin{ex}%[1D3V3-4]%[ma:37acb0fa]
	 Biết hàm số $f(x)=\heva{&\dfrac{x^3-7x+6}{x^2-3x+2}, &\text{khi } x \neq 1\text{ và } x \neq 2\\
	&2a,  &\text{khi } x=1\\
	&b-3, &\text{khi } x=2}$ liên tục trên $\mathbb{R}$.  Tính $P=a^2+b^2$.
	 \choice
		{\True $P=68$} 
		{$P=45$} 
		 {$P=41$}
		{$P=10$} 
	\loigiai{Khi $x \neq 1$ và $x \neq  2$ thì hàm số $f(x)=\dfrac{x^3-7x+6}{x^2-3x+2}$ liên tục trên từng khoảng xác định.\\
		 Tại $x=1$, $x=2$, ta có
		\begin{align*}
			&\lim\limits_{x \to 1} f(x)=\lim\limits_{x \to 1} \dfrac{x^3-7x+6}{x^2-3x+2}=\lim\limits_{x  \to 1} \dfrac{(x-1)(x-2)(x+3)}{(x-1)(x-2)}=\lim\limits_{x \to 1}(x+3)=4\\
			&\lim\limits_{x \to 2} f(x)=\lim\limits_{x \to 2}  \dfrac{x^3-7x+6}{x^2-3x+2}=\lim\limits_{x \to 2} \dfrac{(x-1)(x-2)(x+3)}{(x-1)(x-2)}=\lim\limits_{x \to 2}(x+3)=5\\
			&f(1)=2a;  f(2)=b-3.
		\end{align*}
		Hàm số  liên tục trên $\mathbb{R}$ khi và chỉ khi $\heva{&\lim\limits_{x \to 1}=f(1) \\
		 &\lim\limits_{x \to 2}=f(2)} \Rightarrow\heva{&2a=4\\
		&b-3=5} \Leftrightarrow\heva{&a=2\\ 
		 &b=8.}$\\
		Vậy $P=2^2+8^2=68$. 
	 }
\end{ex}

% nguồn: Ok_Toan11_TracNghiem_Lan2
\begin{ex}%[1D3V3-4]%[ma:59cdf96a]
	Tìm $m$ để hàm số $y=f(x)=\heva{&x^2+2\sqrt{x-2} &\text{ khi } x\ge 2 \\& 5x-5m+m^2&\text{ khi } x < 2}$ liên tục trên  $\mathbb{R}$?
	 \choice
		 {\True $m=2$; $m=3$}
		{$m=-2$; $m=-3$} 
		{$m=1$; $m=6$} 
		 {$m=-1$; $m=-6$}
	 \loigiai{
		Tập xác  định $\mathscr{D}= \mathbb{R}$.
		Xét  trên $\left(2;+\infty \right)$ khi đó $f(x)=x^2+2\sqrt{x-2}$, $\forall x_0 \in \left(2;+\infty \right)$.
		${\lim\limits_{x\to x_0}} \left(x_0 ^2+2\sqrt{x_0-2} \right)=x_0 ^2+2\sqrt{x_0-2}=f\left(x_0 \right) \Rightarrow$ hàm số liên tục  trên $\left(2;+\infty \right)$.
		Xét trên $\left(-\infty; 2\right)$  khi đó $f(x)=5x-5m+m^2$ là hàm số liên tục trên $\left(-\infty; 2\right)$.
		Xét tại $x_0=2$,  ta có: $f(2)=4$.
		\[\lim\limits_{x\to  2^{+}} f(x)=\lim\limits_{x\to 2^{+}} \left(x^2+2\sqrt{x-2} \right)=4; \lim\limits_{x\to 2^{-}} f(x)=\lim\limits_{x\to 2^{-}} \left(5x-5m+m^2\right)=m^2-5m+10.\]
		Để hàm số đã cho liên tục trên  $\mathbb{R}$ thì nó phải liên tục tại $x_0=2$. Khi đó
		\[\lim\limits_{x\to 2^{+}} f(x)=\lim\limits_{x\to 2^{-}} f(x)=f(2)\Leftrightarrow m^2-5m+10=4\Leftrightarrow m^2-5m+6=0\Leftrightarrow \hoac{&m=2 \\&  m=3}.\]
	 }
\end{ex}

% nguồn: Ok_Toan11_TracNghiem_Lan2
\begin{ex}%[1D3V3-4]%[ma:efaf75a4]
	Biết hàm  số $f(x)=\heva{&a x^2+b x-4 &\text{khi } x \leq 1\\
	&2a x-2b &\text{khi } x > 1}$ liên  tục trên $\mathbb{R}$. Tìm giá trị của biểu thức $P=a-3b$.
	 \choice
		{\True $P=-4$} 
		{$P=5$} 
		 {$P=4$}
		{$P=-5$} 
	\loigiai{Ta có $\lim\limits_{x \to 1^{-}}\left(a  x^2+b x-4\right)=a+b-4$; $\lim\limits_{x \to 1^{+}}(2a x-2b)=2a-2b$.
		Hàm số liên tục trên $\mathbb{R}$ khi  và chỉ khi hàm số liên tục tại điểm $x=1$. Điều này tương đương với
		\[\lim\limits_{x \to 1^{-}} f(x)=\lim\limits_{x \to 1^{+}}  f(x) =f(1)\Leftrightarrow a+b-4=2a-2b \Leftrightarrow a-3b=-4.\]
	 }
\end{ex}

% nguồn: Ok_Toan11_TracNghiem_Lan2
\begin{ex}%[1D3V3-4]%[ma:9cbc063f]
	Cho hàm  số $y=\heva{&x+2 &\text{khi } x \geq-1\\
	&x+m &\text{khi } x <-1}$, $m$ là tham số. Tìm $m$ để hàm số  liên tục trên $\mathbb{R}$.
	 \choice
		 {\True $m=2$}
		{$m=-3$} 
		{$m=5$} 
		{$m=3$} 
	\loigiai{Với $x \geq-1$, ta có hàm số $y=x+2$, hàm số  liên tục trên khoảng $(-1;+\infty)$.\\
		Với $x  <-1$, ta có hàm số $y=x+m$, hàm số liên tục trên khoảng $(-\infty;-1)$.\\
		Do đó để hàm số liên tục trên $\mathbb{R}$ thì hàm số cần liên tục tại $x=-1$. Điều này tương đương với \[\lim\limits_{x \to(-1)^{+}} y=\lim\limits_{x  \to (-1)^{-}} y=f(-1).\]
		Ta có $\lim\limits_{x \to(-1)^{+}} y=\lim\limits_{x \to(-1)^{+}}(x+2)=1$; $\lim\limits_{x \to(-1)^{-}} y=\lim\limits_{x \to(-1)^-}(x+m)=-1+m$; $f(-1)=1$.\\ 
		 Suy ra $1=-1+m \Leftrightarrow m=2$.
	 }
\end{ex}

% nguồn: Ok_TracNghiemLop11_Lan1
\begin{ex}%[1D3V3-4]%[ma:c85fd832]
	Có bao nhiêu giá trị thực của tham số $m$ để hàm số $f(x)=\heva{&m^2x^2&\text{khi}&\,\,x\le 2\\&(1-m)x&\text{khi}&\,\, x>2}$ liên  tục trên $\mathbb{R}$?
	\choice 
		 {\True $2$}
		 {$1$}
		{$0$} 
		{$3$} 
	\loigiai{ 
		Tập  xác định $\mathscr{D}=\mathbb{R}$.\\
		Ta có hàm số liên tục trên mỗi khoảng $(-\infty;2)$ và $(2;+\infty)$.\\ 
		Ta có $\heva{&f(2)=4m^2\\&\lim\limits_{x\to  2^{+}} f(x)=\lim\limits_{x\to 2^{+}}[(1-m)x]=2(1-m)\\&\lim\limits_{x\to 2^{-}} f(x)=\lim\limits_{x\to 2^{-}}m^2x^2=4m^2}$\\
		Hàm số $f(x)$ liên  tục trên $\mathbb{R}$ khi và chỉ khi $f(x)$ liên tục tại $x=2$\\
		$\Leftrightarrow\lim\limits_{x\to 2^{+}}  f(x)=\lim\limits_{x\to 2^{-}} f(x)=f(2)\Leftrightarrow 4m^2=2(1-m)\Leftrightarrow 2m^2+m-1=0\Leftrightarrow \hoac{&m=-1\\&m=\dfrac{1}{2}.}$
	 }
\end{ex}

% nguồn: Ok_TracNghiemLop11_Lan1
\begin{ex}%[1D3V3-4]%[ma:7d3921e2]
	Biết rằng hàm số $f(x)=\heva{&\sqrt{x}&\text{khi}&\,\,x\in  [0;4]\\&1+m&\text{khi}&\,\, x\in (4;6]}$ liên tục trên $[0;6]$. Khẳng định nào sau đây \textbf{đúng}?
	\choice 
		 {\True $m<2$}
		 {$2\le m<3$}
		{$3<m<5$} 
		{$m  \ge 5$}
	\loigiai{ 
		Tập xác định $\mathscr{D}=\mathbb{R}$.\\ 
		Hàm số $f(x)$ liên tục trên mỗi khoảng $(0;4)$ và $(4;6)$.\\ Khi đó hàm số liên tục trên đoạn $[0;6]$ khi và chỉ khi hàm số liên  tục tại $x=0$, $x=4$, $x=6$.\\
		Tức là ta cần có $\heva{&\lim\limits_{x\to 0^{+}}f(x)=f(0)\\&\lim\limits_{x\to 6^{-}}f(x)=f(6)\\&\lim\limits_{x\to 4^{-}}f(x)=\lim\limits_{x\to  4^{+}}f(x)=f(4).}\quad (\mathrm{I})$
		\begin{enumerate}[$\bullet$]
			\item  $\heva{&\lim\limits_{x\to 0^{+}}f(x)=\lim\limits_{x\to 0^{+}}\sqrt{x}=0\\&f(0)=\sqrt{0}=0}$
			 \item $\heva{&\lim\limits_{x\to 6^{-}}f(x)=\lim\limits_{x\to 6^{-}}(1+m)=1+m\\&f(6)=1+m}$
			\item $\heva{&\lim\limits_{x\to 4^{-}}f(x)=\lim\limits_{x\to  4^{-}}\sqrt{x}=2\\&\lim\limits_{x\to 4^{+}}f(x)=\lim\limits_{x\to 4^{+}}(1+m)=1+m\\&f(4)=2}$
		\end{enumerate}
		Khi đó $(\mathrm{I})$ trở  thành $1+m=2\Leftrightarrow m=1$.
	 }
\end{ex}

% nguồn: Ok_TracNghiemLop11_Lan1
\begin{ex}%[1D3V3-4]%[ma:46549671]
	Cho hàm số $f(x)=\heva{&\cos\dfrac{\pi x}{2}&\text{khi}&\,\,\left|x\right|\le 1\\&x-1&\text{khi}&\left|x\right| >1 }$.  Mệnh đề nào sau đây là \textbf{sai}?
	 \choice
		{\True Hàm số  liên tục tại $x=-1$}
		{Hàm số liên tục trên  các khoảng $(-\infty;-1)$; $(1;+\infty)$}
		{Hàm số liên tục tại  $x=1$}
		{Hàm số liên tục trên  khoảng $(-1;1)$}
	\loigiai{ 
		 Tập xác định $\mathscr{D}=\mathbb{R}$.\\
		Ta có $f(x)$ liên tục trên các khoảng  $(-\infty;-1)$, $(-1;1)$, $(1;+\infty)$.\\
		Mặt khác $\heva{&f(-1)=\cos\left(-\dfrac{\pi}{2}\right)=0\\&\lim\limits_{x\to (-1)^{+}}f(x)=\lim\limits_{x\to  (-1)^{+}}(x-1)=-2.}$ \\
		Vậy hàm số $f(x)$ không liên  tục tại $x=-1$.
	} 
\end{ex}

% nguồn: Ok_TracNghiemLop11_Lan1
\begin{ex}%[1D3V3-4]%[ma:a1cf761a]
	Có bao nhiêu giá trị của tham số $a$ để hàm số $f(x)=\heva{&\dfrac{x^2-3x+2}{\left|x-1\right|}&\text{khi}&\,\,x\ne  1\\&a&\text{khi}&\,\,x=1}$ liên tục trên $\mathbb{R}$?
	\choice 
		{$1$} 
		 {$2$}
		{\True  $0$}
		{$3$} 
	 \loigiai{
		Tập xác  định $\mathscr{D}=\mathbb{R}$.\\
		Hàm số $f(x)$ liên tục trên các khoảng $(-\infty;1)$ và  $(1;+\infty)$.\\ Khi đó hàm số đã cho liên tục trên $\mathbb{R}$ khi và chỉ khi nó liên tục tại $x=1$, tức là ta cần có $\lim\limits_{x\to 1}f(x)=f(1)\Leftrightarrow \lim\limits_{x\to 1^{+}}f(x)=\lim\limits_{x\to 1^{-}}f(x)=f(1)\quad (\mathrm{I})$\\
		Ta có $f(x)=\heva{&x-2&\text{khi}\,\,&x>1\\&a&\text{khi}\,\,&x=1\\&2-x&\text{khi}\,\,&x<1}$\\ 
		Nên $\heva{&\lim\limits_{x\to 1^{-}}f(x)=\lim\limits_{x\to 1^{-}}(2-x)=1\\&\lim\limits_{x\to 1^{+}}f(x)=\lim\limits_{x\to 1^{+}}(x-2)=-1}\Rightarrow (\mathrm{I})$ không thỏa  mãn với mọi $x\in \mathbb{R}$.\\
		Không tồn tại giá trị $a$  thỏa yêu cầu.
	 }
\end{ex}

% nguồn: Ok_TracNghiemLop11_Lan1
\begin{ex}%[1D3V3-4]%[ma:8c76c031]
	Cho hàm số $f(x)$ liên tục trên đoạn $[-1;4]$ sao cho $f(-1)=2$, $f(4)=7$. Có thể kết luận gì về số  nghiệm của phương trình $f(x)=5$ trên đoạn $[-1;4]$.
	 \choice
		{Vô nghiệm} 
		{\True Có ít  nhất một nghiệm}
		{Có đúng một  nghiệm}
		{Có đúng hai nghiệm} 
	 \loigiai{
		 Ta có $f(x)=5\Leftrightarrow f(x)-5=0$.\\
		Đặt  $g(x)=f(x)-5$. Khi đó $\heva{&g(-1)=f(-1)-5=2-5=-3\\&g(4)=f(4)-5=7-5=2.}$ \\
		 $\Rightarrow g(-1)\cdot g(4)=(-3)\cdot 2=-6<0$.\\
		Vậy phương trình $g(x)=0$ có ít nhất một nghiệm thuộc khoảng $(-1;4)$ hay phương trình $f(x)=5$ có ít nhất một nghiệm  thuộc khoảng $(-1;4)$.
	} 
\end{ex}

% nguồn: Ok_TracNghiemLop11_Lan1
\begin{ex}%[1D3V3-4]%[ma:25091fa2]
	 Hàm số $y=f(x)=\heva{&3&\text{khi}&\,\,x=-1\\
	&\dfrac{x^4+x}{x^2+x}&\text{khi}&\,\,x\ne-1, x\ne 0\\ 
	&1&\text{khi}&\,\,x=0}$  liên tục tại:
	 \choice
		{mọi  điểm trừ $x=0$, $x=1$}
		{\True mọi điểm  $x\in \mathbb{R}$}
		{mọi điểm trừ $x=-1$} 
		{mọi  điểm trừ $x=0$}
	\loigiai{ 
		 Tập xác định $\mathscr{D}=\mathbb{R}$.\\
		Ta  có hàm số liên tục trên mỗi khoảng $(-\infty;-1)$, $(-1;0)$ và $(0;+\infty)$.
		\begin{enumerate}[(i)]
			\item Xét tại $x=-1$,  ta có
			\begin{eqnarray*}
				\lim\limits_{x\to  -1} f(x)
				&=& 
				\lim\limits_{x\to -1} \dfrac{x^4+x}{x^2+x}\\ 
				 &=&
				 \lim\limits_{x\to -1} \dfrac{x(x+1)(x^2-x+1)}{x(x+1)}\\
				 &=&
				\lim\limits_{x\to  -1}(x^2-x+1)=3=f(-1)
			\end{eqnarray*}
			Vậy  hàm số $y=f(x)$ liên tục tại $x=-1$.
			\item Xét tại $x=0$,  ta có
			\begin{eqnarray*}
				\lim\limits_{x\to 0} f(x) 
				 &=&
				 \lim\limits_{x\to 0} \dfrac{x^4+x}{x^2+x}\\
				&=& 
				\lim\limits_{x\to  0} \dfrac{x(x+1)(x^2-x+1)}{x(x+1)}\\
				 &=&
				\lim\limits_{x\to  0}(x^2-x+1)=1=f(0)
			\end{eqnarray*}
			Vậy hàm số $y=f(x)$ liên tục  tại $x=0$.
		\end{enumerate}
	} 
\end{ex}

% nguồn: Ok_Toan11_TracNghiem_Lan2
\begin{ex}%[1D3C3-4]%[ma:6f7e31ad]
	Tìm tất cả các giá trị thực của $m$ để hàm số $f(x)=\heva{&\dfrac{\sqrt{x+1}-1}{x} &\text{ khi }x >  0 \\
	&\sqrt{x^2+1}-m &\text{ khi  } x \le 0}$ liên tục trên $\mathbb{R}$.
	\choice 
		 {$m=\dfrac{3}{2}$}
		{\True $m=\dfrac{1}{2}$} 
		 {$m=-2$}
		{$m=-\dfrac{1}{2}$} 
	 \loigiai{
		Khi $x > 0$ ta có: $f(x)=\dfrac{\sqrt{x+1}-1}{x}$ liên tục trên khoảng $\left(0;+\infty \right)$. 
		 Khi $x < 0$ ta có: $f(x)=\sqrt{x^2+1}-m$ liên tục trên khoảng $\left(-\infty;0\right)$.
		Hàm số liên  tục trên $\mathbb{R}$ khi và chỉ khi hàm số liên tục tại $x=0$.
		Ta có: 
		\begin{eqnarray*}
			 \lim\limits_{x\to 0^{+}} f(x)&=&\lim\limits_{x\to 0^{+}}\dfrac{\sqrt{x+1}-1}{x}=\lim\limits_{x\to 0^{+}}\dfrac{1}{\sqrt{x+1}+1}=\dfrac{1}{2}\\
			\lim\limits_{x\to  0^{-}}f(x)&=&\lim\limits_{x\to 0^{-}}\left(\sqrt{x^2+1}-m\right)=1-m=f(0).
		\end{eqnarray*}
		Do đó  hàm số liên tục tại $x=0$ khi và chỉ khi $\dfrac{1}{2}=1-m\Leftrightarrow m=\dfrac{1}{2}$.
	} 
\end{ex}

% nguồn: Ok_Toan11_TracNghiem_Lan2
\begin{ex}%[1D3C3-4]%[ma:0ec9ed02]
	Nếu hàm số $f(x)=\heva{&x^2+ax+b &\text{ khi }& x <-5 \\&  x+17 &\text{ khi }&-5\le x\le 10 \\ &ax+b+10&\text{ khi }& x > 10}$ liên tục trên $ \mathbb{R}$ thì $a+b$ bằng
	 \choice
		{\True  $-1$}
		 {$0$}
		{$1$} 
		{$2$} 
	 \loigiai{
		Với $x <-5$ ta có $f(x)=x^2+ax+b$, là hàm đa thức nên liên tục trên  $\left(-\infty;-5\right)$.
		Với $-5< x  < 10$ ta có $f(x)=x+7$, là hàm đa thức nên liên tục trên $\left(-5;10\right)$.
		Với $x > 10$ ta có $f(x)=ax+b+10$, là hàm đa thức nên  liên tục trên $\left(10;+\infty \right)$.
		Để hàm  số liên tục trên $ \mathbb{R}$ thì hàm số phải liên tục tại $x=-5$ và $x=10$.
		Ta  có: $f\left(-5\right)=12$; $f(10)=17$.
		\begin{eqnarray*}
			 &&\lim\limits_{x\to-5^{-}} f(x)=\lim\limits_{x\to-5^{-}} \left(x^2+ax+b\right)=-5a+b+25.\\
			&&\lim\limits_{x\to-5^{+}} f(x)=\lim\limits_{x\to-5^{+}} \left(x+17\right)=12.\\ 
			&&\lim\limits_{x\to 10^{-}} f(x)=\lim\limits_{x\to 10^{-}}  \left(x+17\right)=27.\\
			&&\lim\limits_{x\to 10^{+}} f(x)=\lim\limits_{x\to 10^{+}} \left(ax+b+10\right)=10a+b+10. 
		\end{eqnarray*}
		Hàm số liên tục tại $x=-5$ và $x=10$  khi
		\[ 
		\heva{  &{5a+b+25=12} \\ &{10a+b+10=27}} \Leftrightarrow
		\heva{  &{-5a+b=-13} \\ &{10a+b=17} } \Leftrightarrow
		\heva{ &{a=2} \\  &{b=-3}} \Rightarrow a+b=-1.
		 \]
	} 
\end{ex}

% nguồn: Ok_TracNghiemLop11_Lan1
\begin{ex}%[1D3K3-4]%[ma:1208e0c9]
	Biết các  số $\mathrm{C}_n^1$, $\mathrm{C}_n^2$, $\mathrm{C}_n^3$ theo thứ tự lập thành một cấp số cộng với $n>3$. Tìm $n$.
	\choice 
		{ $n=5$ } 
		 { $n=11$ }
		{  $n=9$ }
		{ \True  $n=7$ }
	 \loigiai{
		Ba số $\mathrm{C}_n^1$, $\mathrm{C}_n^2$, $\mathrm{C}_n^3$ theo thứ tự ${{u}_{1}},\,\,{{u}_{2}},\,\,{{u}_{3}}$ lập thành cấp số cộng nên 
		\begin{eqnarray*}
			&&{{u}_{1}}+{{u}_{3}}=2{{u}_{2}}\Leftrightarrow \mathrm{C}_n^1+\mathrm{C}_n^3=2\mathrm{C}_n^1\,\,(n\ge 3)\\ 
			 &\Leftrightarrow& n+\dfrac{(n-2)(n-1)n}{6}=2\cdot \dfrac{(n-1)n}{2}\Leftrightarrow 1+\dfrac{{{n}^2}-3n+2}{6}=n-1\\
			&\Leftrightarrow& {{n}^2}-9n+14=0\Leftrightarrow \hoac{ 
			& n=2  \\
			& n=7  \\
			}\Leftrightarrow  n=7\,(n\ge 3).
		\end{eqnarray*}
	\textbf{Nhận xét:}\\ Nếu ${{u}_{k-1}},\,\,{{u}_{k}},\,\,{{u}_{k+1}}$ là ba số hạng liên tiếp của một cấp số cộng thì  ta có ${{u}_{k-1}}+{{u}_{k+1}}=2{{u}_{k}}$.}
\end{ex}

% nguồn: Ok_TracNghiemLop11_Lan1
\begin{ex}%[1D3K3-4]%[ma:b7c6b891]
	Cho $a;b;c$ theo thứ tự lập thành cấp  số cộng. Mệnh đề nào sau đây đúng?
	\choice 
		 { \True $a^2+{{c}^2}+2ac=4b^2$}
		{  $a^2+{{c}^2}=2ab-2bc$}
		{ $a^2-{{c}^2}=2ab-2bc$} 
		{  $a^2-{{c}^2}=ab-bc$}
	\loigiai{ 
	Ta có $a+c=2b\Rightarrow  {{(a+c)}^2}=4b^2\Leftrightarrow a^2+{{c}^2}+2ac=4b^2$.}
\end{ex}

% nguồn: Ok_TracNghiemLop11_Lan1
\begin{ex}%[1D3K3-4]%[ma:7b20d66d]
	Nếu $a;b;c$ theo thứ tự lập thành cấp số cộng thì dãy số nào sau  đây lập thành cấp số cộng?
	\choice 
		{  $2b;a;c$ }
		{  \True $-2b;-2a;-2c$ }
		{ $2b;-a;-c$} 
		{ $2b^2;a^2;{{c}^2}$  }
	 \loigiai{
	 Ta có $c+a=2b\Rightarrow -2(c+a)=-2(2b)\Leftrightarrow (-2c)+(-2a)=2(-2b)$.}
\end{ex}


%%%%%%%%%%%%%%%% ĐÚNG/SAI %%%%%%%%%%%%%%%%

% nguồn: Ok_dungsai_lop11_Lan1
\begin{ex}%[1D3H3-4]%[ma:6d5d6f54]
 Cho hàm số $y=f(x)=\heva{&\dfrac{2-\sqrt{x+3}}{3x+6}-3\  & \text{khi}&\ x>1\\
 &\dfrac{\sqrt{x^2}}{4}+3mx\ & \text{khi}&\ x\le 1}$. Các mệnh  đề sau đúng hay sai?
  \choiceTF
  {\True Tập  xác định $\mathscr{D}=\mathbb{R}$}
   {$\displaystyle\lim\limits_{x\to 1^+}f(x)=3$}
  {\True $\displaystyle\lim\limits_{x\to 1^-}f(x)=\dfrac{1}{4}+3m$} 
  {Với $m=-\dfrac{11}{12}$ thì  hàm số $f(x)$ liên tục trên $\mathbb{R}$}
  \loigiai{
  \begin{itemchoice}
   \itemch \textbf{Đúng}.\\ Tập xác định $\mathscr{D}=\mathbb{R}$. 
   \itemch \textbf{Sai}.\\ Ta  có
   $\displaystyle\lim\limits_{x\to 1^+}f(x)=\displaystyle\lim\limits_{x\to  1^+}\left[\dfrac{2-\sqrt{x+3}}{3x+6}-3\right]=\dfrac{2-\sqrt{1+3}}{3\cdot 1+6}-3=-3$.
   \itemch \textbf{Đúng}.\\$\displaystyle\lim\limits_{x\to  1^-}f(x)=\displaystyle\lim\limits_{x\to 1^-}\left[\dfrac{\sqrt{x^2}}{4}+3mx\right] =\dfrac{1}{4}+3m$.
   \itemch \textbf{Sai}.\\ Với $x<1$  thì $f(x)=\dfrac{\sqrt{x^2}}{4}+3mx$ liên tục trên $(-\infty;1)$ và với $x>1$ thì $f(x)=\dfrac{2-\sqrt{x+3}}{3x+6}-3$ liên tục trên $(1;+\infty)$.\\
   Do đó hàm số $y=f(x)$ liên tục trên $\mathbb{R}$ khi và chỉ khi $y=f(x)$ liên tục tại $x=1$. Điều này tương đương  với
   \[\displaystyle\lim\limits_{x\to 1^+}f(x)=\displaystyle\lim\limits_{x\to 1^-}f(x)=f(1)\Leftrightarrow 3=\dfrac{1}{4}+3m\Leftrightarrow m=\dfrac{11}{12}.\] 
   Vậy hàm số  hàm số $y=f(x)$ liên tục trên $\mathbb{R}$ khi và chỉ khi $m=\dfrac{11}{12}$.
  \end{itemchoice}
  }
\end{ex}

% nguồn: Ok_dungsai_lop11_Lan1
\begin{ex}%[1D3H3-4]%[ma:8796dbf1]
 Xét tính liên tục của các hàm số sau trên tập xác  định của nó
  \choiceTF
  {\True $f(x)=x^3-x^2+8x$ là hàm số  liên tục trên $\mathbb{R}$}
  {$f(x)=\dfrac{x^2}{x^2-3x}$ là hàm số  liên tục trên khoảng $(-\infty;+\infty)$}
  {$f(x)=\dfrac{\sin x+1}{x+1}$ là hàm số liên tục trên các khoảng  $(-\infty;0)$ và $(0;+\infty)$}
  {\True $f(x)=\sqrt{x-2}$ là hàm số liên  tục trên nửa khoảng $[2;+\infty)$}
 \loigiai{ 
  \begin{itemchoice}
   \itemch Vì $f(x)=x^3-x^2+8x$ là hàm đa thức nên hàm số liên  tục trên $\mathbb{R}$.
   \itemch Vì $f(x)=\dfrac{x^2}{x^2-3x}$ là hàm phân thức xác định trên các khoảng $(-\infty;0)$, $(0;3)$, $(3;+\infty)$ nên hàm số liên tục trên các khoảng $(-\infty;0)$, $(0;3)$,$(3;+\infty)$.\\ 
   Ta  có $\heva{&\lim\limits_{x\to 3^{-}}f(x)=\lim\limits_{x\to 3^{-}}\dfrac{x^2}{x^2-3x}=-\infty\\&\lim\limits_{x\to 3^{+}}f(x)=\lim\limits_{x\to 3^{+}}\dfrac{x^2}{x^2-3x}=+\infty.}$\\
    Vậy hàm số $f(x)$ gián đoạn tại điểm $x=3$.
   \itemch Tập xác định $\mathscr{D}=(-\infty;-1)\cup  (-1;+\infty)$.\\
   Các hàm số $y=\sin x+1$ và $y=x+1$ đều liên tục trên các khoảng  $(-\infty;-1)$,$ (-1;+\infty)$.\\
   Nên $f(x)$ liên tục trên các khoảng $(-\infty;-1)$,$  (-1;+\infty)$.
   \itemch Hàm số xác định trên  nửa khoảng $[2;+\infty)$ nên liên tục trên khoảng $(2;+\infty)$.\\
   Xét $x_0=2$, ta có  $\heva{&f(2)=0\\&\lim\limits_{x\to 2^{+}}f(x)=\lim\limits_{x\to 2^{+}}\sqrt{x-2}=0.}$\\
   Vậy hàm số  $f(x)$ liên tục tại điểm $x_0=2$.
 \end{itemchoice}}
\end{ex}

% nguồn: Ok_dungsai_lop11_Lan1
\begin{ex}%[1D3H3-4]%[ma:016aba8b]
 Cho hàm số $f(x)=\dfrac{x-1}{x+1}$. Hãy cho biết tính đúng sai của mỗi khẳng  định sau?
  \choiceTF
   {Hàm số liên tục trên $\mathbb{R}$}
  {\True  Hàm số gián đoạn tại $x=-1$}
   {Hàm số liên tục trong khoảng $(-3;1)$}
  {\True Hàm số liên tục trên khoảng $(-3;-1)$} 
 \loigiai{ 
  Đây là hàm phân thức hữu tỉ có tập xác định là  $\mathbb{R}\setminus \{-1\}$. Do đó
  \begin{itemchoice}
   \itemch Hàm số không liên tục trên  $\mathbb{R}$.
   \itemch Hàm số gián đoạn tại  $x=-1$.
   \itemch Hàm số không liên tục trong khoảng  $(-3;1)$ vì $-1 \in (-3;1)$.
   \itemch Hàm số liên tục trên khoảng  $(-3;-1)$.
  \end{itemchoice}
 } 
\end{ex}

% nguồn: Ok_dungsai_lop11_Lan1
\begin{ex}%[1D3H3-4]%[ma:b47c018d]
 Cho hàm số $f(x)=\heva{&\dfrac{x^2+2x-3}{x-1}&\text{ nếu } x\ne 1 \\& x-3m &\text{ nếu }  x=1}$. Hãy cho biết tính đúng sai của mỗi khẳng định sau?
 \choiceTF 
   {\True $\lim\limits_{x\to 1} f(x)=4$}
   {\True Hàm số liên tục trên khoảng $(-\infty;0)$}
  {\True Hàm  số liên tục trên $\mathbb{R}$ khi $m=-1$}
   {Hàm số không xác định tại $x=1$}
 \loigiai{ 
  \begin{itemchoice}
   \itemch $\lim\limits_{x\to  1} f(x)=\lim\limits_{x\to 1}\dfrac{x^2+2x-3}{x-1}=\lim\limits_{x\to 1}\dfrac{(x-1)(x+3)}{x-1}= \lim\limits_{x\to 1} (x+3)= 4$.
   \itemch Khi $x\in (-\infty; 0)$ thì $f(x)=\dfrac{x^2+2x-3}{x-1}$. Hàm số phân thức hữu  tỉ liên tục trên từng khoảng xác định nên nó liên tục trên khoảng $(-\infty; 0)$.
   \itemch Với  $m=-1$. Ta có $f(x)=\heva{&\dfrac{x^2+2x-3}{x-1}&\text{ nếu } x\ne 1 \\& 4 &\text{ nếu } x=1}$. Hàm số liên tục trên $\mathbb{R}$.
   \itemch Ta có $f(1)=1-3m$. Do đó hàm số  xác định tại $x=1$.
  \end{itemchoice}
  }
\end{ex}

% nguồn: Ok_dungsai_lop11_Lan1
\begin{ex}%[1D3H3-4]%[ma:59a4044d]
 Cho hàm số $f(x)=\heva{ & \dfrac{\sqrt{x+3}-x-1}{x^2-5 x+4} \text { khi } x>1 \\ & 3 x+m \text { khi  } x \leq 1}$ với $m$ là tham số và hàm số $g(x)=x^2-x-2$.
 \choiceTF 
  {\True Hàm số $f(x)$  liên tục trên khoảng $(0; 1)$}
  {Hàm số $f(x)$ liên tục trên  khoảng $(1;+\infty)$}
  {Hàm số $\dfrac{f(x)}{g(x)}$ liên tục trên khoảng  $(-\infty; 1)$}
  {Với $m=\dfrac{11}{4}$ thì hàm số $f(x)$ liên tục trên tập xác định} 
 \loigiai{ 
  \begin{itemchoice}
   \itemch Đúng. Xét $x \leq 1$ ta có $f(x)=3 x+m$  xác định với mọi với $x \leq 1$.\\
   $\Rightarrow$ Hàm số $f(x)$  liên tục trên khoảng $(-\infty; 1)$\\
   $\Rightarrow$ Hàm số  $f(x)$ liên tục trên khoảng $(0; 1)$.
   \itemch Sai. Hàm số $f(x)$ không xác định  tại $x=4$ nên hàm số $f(x)$ gián đoạn tại $x=4$.\\
   Vậy hàm  số $f(x)$ không liên tục trên khoảng $(1;+\infty)$.
    \itemch Sai. Ta có $g(x)=0 \Leftrightarrow \hoac{&x=-1 \\ & x=2}$\\
   Hàm số $\dfrac{f(x)}{g(x)}$ không xác định tại $x=-1 \in(-\infty; 1)$  nên hàm số $f(x)$ gián đoạn tại $x=-1$.\\
   Vậy hàm số  $\dfrac{f(x)}{g(x)}$ không liên tục trên khoảng $(-\infty; 1)$.
   \itemch Sai. Xét $x \leq 1$ ta có $f(x)=3 x+m$ xác định với mọi với $x \leq  1$\\
   Xét $x>1$ ta có  $f(x)=\dfrac{\sqrt{x+3}-x-1}{x^2-5 x+4}$\\
   Hàm số xác định khi $\heva{&x+3 \geq 0 \\  & x^2-5 x+4 \neq 0} \Leftrightarrow \heva{&x \geq-3 \\ & x \neq 1 \\ & x \neq 4}$
    Mà $x>1 \Rightarrow x \in(1;+\infty) \backslash\{4\}$
   Từ (1) và $(2)$ suy  ra hàm số có tâp xác định là $D=\mathbb{R} \backslash\{4\}$.
   Hàm số xác định trên các khoảng $(-\infty; 1);(1; 4);(4;+\infty)$ nên hàm số  liên tục trên các khoảng này.
   Xét  tại $x=1$
   \begin{align*}
    \lim \limits_{x \rightarrow 1^{+}} f(x)=&\lim \limits_{x \rightarrow  1^{+}} \dfrac{\sqrt{x+3}-x-1}{x^2-5 x+4}\\
    =&\lim \limits_{x  \rightarrow 1^{+}} \dfrac{x+3-(x+1)^2}{(x-1)(x-4)(\sqrt{x+3}+x+1)} \\
    =&\lim \limits_{x  \rightarrow 1^{+}} \dfrac{-x-2}{(x-4)(\sqrt{x+3}+x+1)}=\dfrac{1}{4} .
   \end{align*}
   $ \lim \limits_{x \rightarrow 1^{-}} f(x)=\lim \limits_{x \rightarrow 1^{-}}(3 x+m)=3+m $.  \\
   $ f(1)=3+m $. 
   Hàm số liên tục trên tập xác định khi  hàm số liên tục tại $x=1$
    $$
   \Leftrightarrow \lim \limits_{x  \rightarrow 1^{+}} f(x)=\lim \limits_{x \rightarrow 1^{-}} f(x)=f(1) \Leftrightarrow 3+m=\dfrac{1}{4} \Leftrightarrow m=-\dfrac{11}{4} \text {. }
    $$
  \end{itemchoice}
  }
\end{ex}

% nguồn: Ok_dungsai_lop11_Lan2
\begin{ex}%[1D3H3-4]%[ma:066239b3]
	[Dự án  11HVA 2024-2025]
	Cho  hàm số $f(x)=\heva{&\dfrac{x-1}{\sqrt{3x-2}-x} \text{khi } x > 1\\
	&x+m \text{khi  } x \leq 1}$ với $m$ là tham số và hàm số $g(x)=3x-1$.
	\choiceTF 
		{\True Hàm số $f(x)$  gián đoạnn tại $x=-2$}
		{\True Hàm số $f(x)$ gián  đoạn tại $x=2$}
		{\True Hàm số $f(x)+g(x)$ liên  tục tại $x=-3$}
		{\True Hàm số $f(x)$ liên tục tại $x=1$ khi $m=1$} 
	\loigiai{\begin{itemchoice}
		 \itemch \textbf{Đúng.} \\
		Với $x \leq 1$ ta có $f(x)=x+m$ luôn xác định  với mọi $x \leq 1$.\\
		Suy ra hàm số liên tục trên khoảng $(-\infty; 1)$ nên hàm  số $f(x)$ liên tục tại $x=-2$.
		\itemch \textbf{Đúng.}  \\
		Xét  $x > 1$ ta có $f(x)=\dfrac{x-1}{\sqrt{3x-2}-x}\cdot$\\
		Hàm số không xác định tại $x=2$ nên hàm số gián đoạn  tại $x=2$.
		 \itemch \textbf{Đúng.} \\
		Ta  có $g(x)=3x-1$ luôn xác định với mọi $x \in \mathbb{R}$.\\
		Suy ra hàm số $g(x)$  liên tục trên $\mathbb{R}$ nên hàm số $g(x)$ liên tục tại $x=-3$.\\
		Mà hàm số $f(x)$ cũng liên tục tại $x=-3$ suy ra hàm  số $f(x)+g(x)$ liên tục tại $x=-3$.
		\itemch  \textbf{Đúng.} \\
		$\lim\limits_{x \to 1^{+}} f(x)=\lim\limits_{x \to 1^{+}}  \dfrac{x-1}{\sqrt{3x-2}-x}=\lim\limits_{x \to 1^{+}} \dfrac{(x-1)(\sqrt{3x-2}+x)}{3x-2-x^2}=\lim\limits_{x \to 1^{+}} \dfrac{\sqrt{3x-2}+x}{2-x}=2$.\\
		$\lim\limits_{x \to 1^{-}} f(x)=1+m; f(1)=1+m$.\\ 
		Hàm số $f(x)$  liên tục tại $x=1$ khi và chỉ khi
		$\lim\limits_{x \to 1^{+}} f(x)=\lim\limits_{x \to 1^{-}} f(x)=f(1) \Leftrightarrow 1+m=2\Leftrightarrow m=1$. 
	\end{itemchoice}
 }
\end{ex}

% nguồn: Ok_dungsai_lop11_Lan2
\begin{ex}%[1D3H3-4]%[ma:6a5f233c]
	[Dự án 11HVA 2024-2025] 
	Cho hàm số $f(x)=\heva{&\dfrac{1-\sqrt{5x+11}}{2x^2-5x-18} \text{khi } x >-2\\ 
	&4-x^2\text{khi } x \leq-2}$ và  $g(x)=\heva{&\dfrac{x^2-x-6}{x+2} \text{khi } x \neq-2\\
	&2x+a \text{khi  } x=-2\cdot
	 }$
	\choiceTF 
		{\True Ta có $\lim\limits_{x \to-2^{+}} f(x)=\dfrac{5}{26}$} 
		{Hàm số  $f(x)$ liên tục tại điểm $x_0=-2$}
		{Để hàm số $g(x)$ liên tục tại điểm $x_0=-2$  thì $a=1$}
		 {\True Khi $a=-1$ thì hàm số $y=f(x) \cdot g(x)$ gián đoạn tại điểm $x_0=-2$}
	\loigiai{\begin{itemchoice}
		\itemch \textbf{Đúng.} \\ 
		Ta có $\lim\limits_{x \to  x_0^{+}} f(x)=\lim\limits_{x \to-2^{+}} \dfrac{1-\sqrt{5x+11}}{2x^2-5x-18}=\dfrac{5}{26}$.
		\itemch  \textbf{Sai.} \\
		 Ta có $f(-2)=\lim\limits_{x \to-2^{-}} f(x)=0\neq \lim\limits_{x \to-2^{+}} f(x)$.\\
		Vậy không tồn tại $\lim\limits_{x \to-2} f(x)$ hay hàm số gián đoạn tại điểm $x_0=-2$. 
		\itemch \textbf{Sai.} \\ 
		Ta có $g(x_0)=g(-2)=-4+a$.\\ 
		 $\lim\limits_{x \to x_0} g(x)=\lim\limits_{x \to-2} \dfrac{x^2-x-6}{x+2}=-5$.\\
		Để hàm số liên tục tại điểm $x_0=-2$ thì  $\lim\limits_{x \to-2} g(x)=g(-2)$.\\
		 Suy ra $-4+a=-5\Leftrightarrow a=-1$.
		\itemch \textbf{Đúng.} \\ 
		Khi  $a=-1$, $g(x)$ liên tục tại điểm $x_0=-2$, $f(x)$ gián đoạn tại điểm $x_0=-2$.\\
		Do đó hàm số $y=f(x)  \cdot g(x)$ gián đoạn tại điểm $x_0=-2$.
	\end{itemchoice}
 }
\end{ex}

% nguồn: Ok_dungsai_lop11_Lan2
\begin{ex}%[1D3H3-4]%[ma:03adf9de]
	[Dự  án 11HVA 2024-2025]
	Cho hàm số $f(x)=\heva{&\dfrac{1}{x-1}-\dfrac{3}{x^3-1}  \text{khi } x > 1\\
	&m x+2\quad \text{khi } x \leq 1\cdot 
	 }$
	\choiceTF 
		{\True Hàm số $f(x)$ xác định tại  $x=1$}
		{\True Hàm  số $f(x)$ liên tục tại $x=-2$}
		{\True Khi $m=-2$, hàm số $f(x)$ gián đoạn tại $x=1$} 
		{Hàm số  $f(x)$ liên tục trên $\mathbb{R}$ khi $m=1$}
	\loigiai{Tập xác định: $\mathscr{D}=\mathbb{R}$.\\ 
		Với $x \in(-\infty; 1): f(x)=m x+2$  là hàm đa thức nên liên tục trên $(-\infty; 1)$.\\
		Với $x \in(1;+\infty):  f(x)=\dfrac{1}{x-1}-\dfrac{3}{x^3-1}$ là hiệu hai hàm phân thức nên liên tục trên $(1;+\infty)$.\\
		Với  $x=1$:
		\begin{itemize}
			\item $\lim\limits_{x \to 1^{+}} f(x)=\lim\limits_{x \to 1^{+}}\left(\dfrac{1}{x-1}-\dfrac{3}{x^3-1}\right)=\lim\limits_{x \to 1^{+}} \dfrac{x^2+x-2}{x^3-1}\\=\lim\limits_{x  \to 1^{+}} \dfrac{(x-1)(x+2)}{(x-1)\left(x^2+x+1\right)}=\lim\limits_{x \to 1^{+}} \dfrac{x+2}{x^2+x+1}=1$.
			 \item $\lim\limits_{x \to 1^{-}}(m x+2)=2+m$.
			 \item $f(1)=2+m$.
		\end{itemize}
		Hàm số $f(x)$ liên tục trên $\mathbb{R} \Leftrightarrow$ hàm  số $f(x)$ liên tục tại $x=1$ $\Leftrightarrow \lim\limits_{x \to 1^{+}} f(x)=\lim\limits_{x \to 1^{-}} f(x)=f(1) \Leftrightarrow 2+m=1\Leftrightarrow m=-1$.
		\begin{itemchoice}
			\itemch \textbf{Đúng.} 
			\itemch  \textbf{Đúng.} \\
			Vì $x \in(-\infty; 1): f(x)=m x+2$  là hàm đa thức nên liên tục trên $(-\infty; 1)$.
			\itemch \textbf{Đúng.} \\ 
			Vì hàm  số $f(x)$ gián đoạn tại $x=1$ khi $2+m \neq 1\Leftrightarrow m \neq-1$.
			\itemch \textbf{Sai.}  \\
			Vì hàm số $f(x)$ liên tục  trên $\mathbb{R}$ khi $m=-1$.
		\end{itemchoice}
	} 
\end{ex}

% nguồn: Ok_dungsai_lop11_Lan2
\begin{ex}%[1D3H3-4]%[ma:b003e282]
	[Dự án 11HVA 2024-2025] 
	Cho  hàm số $h(x)=\heva{&2x \text{khi } x < 0\\
	&x^2+1\text{khi } 0\leq x \leq 2\\ 
	&3x-1\text{khi } x >  2\cdot
	 }$
	\choiceTF 
		{\True  Hàm số liên tục tại $x=-1$}
		{\True Hàm số  liên tục tại $x=3$}
		{Hàm số liên tục tại $x=0$} 
		 {\True Hàm số liên tục tại $x=2$}
	\loigiai{Hàm số  $y=h(x)$ có TXĐ: $\mathscr{D}=\mathbb{R}$.\\
		Dễ thấy hàm số $y=h(x)$ liên  tục trên mỗi khoảng $(-\infty; 0),(0; 2)$ và $(2;+\infty)$. \\
		Vì ta  có $\heva{&h(0)=1\\
		&\lim\limits_{x \to 0^{-}} h(x)=\lim\limits_{x \to 0^{-}} 2x=0}  \Rightarrow f(x)$ không liên tục tại $x=0$.\\
		Vì  ta có $\heva{&h(2)=5\\
		&\lim\limits_{x \to x^{-}} h(x)=\lim\limits_{x \to 2^{-}}\left(x^2+1\right)=5\Rightarrow f(x) \text{liên  tục tại } x=2\cdot \\
		&\lim\limits_{x \to 2^{+}} h(x)=\lim\limits_{x \to 2^{+}}(3x-1)=5\cdot 
		 }$
		\begin{itemchoice}
			\itemch  \textbf{Đúng.} \\
			Hàm số liên tục tại $x=-1$. 
			\itemch \textbf{Đúng.} \\ 
			Hàm số liên tục tại  $x=3$.
			\itemch  \textbf{Sai.} \\
			 Hàm số liên tục tại $x=0$.
			\itemch \textbf{Đúng.} \\ 
			Hàm  số liên tục tại $x=2$.
		\end{itemchoice}
	} 
\end{ex}

% nguồn: Ok_dungsai_lop11_Lan2
\begin{ex}%[1D3H3-4]%[ma:c365f283]
	 [Dự án 11HVA 2024-2025]
	Cho phương trình $2x^4-5x^2+x+1=0$. Xét tính đúng sai của các mệnh  đề sau
	\choiceTF 
		{Phương trình không có nghiệm trong khoảng  $(-1; 1)$}
		{Phương trình không có nghiệm trong  khoảng $(-2; 0)$}
		{Phương trình chỉ có một nghiệm trong khoảng  $(-2; 1)$}
		{\True Phương trình có ít nhất hai nghiệm trong khoảng $(0;  2)$}
	\loigiai{Hàm số $f(x)=2x^4-5x^2+x+1$ là hàm đa thức có tập xác định là $\mathbb{R}$ nên liên tục trên  $\mathbb{R}$.
		Ta có 
		\begin{itemize}
			\item $\heva{&f(0)=1\\ 
			&f(-1)=-3} \Rightarrow f(-1) \cdot f(0) < 0\Rightarrow f(x)=0$ có ít nhất một nghiệm $x_1$ thuộc $(-1; 0)$. 
			 \item $\heva{&f(0)=1\\
			&f(1)=-1} \Rightarrow f(0) \cdot f(1) < 0\Rightarrow f(x)=0$ có ít nhất một nghiệm  $x_2$ thuộc $(0; 1)$.
			\item $\heva{&f(1)=-1\\ 
			&f(2)=15} \Rightarrow f(1) \cdot f(2) < 0\to f(x)=0$ có ít nhất một nghiệm $x_3$ thuộc  $(1; 2)$.
		\end{itemize}
		Vậy phương trình $f(x)=0$ đã cho có các nghiệm $x_1, x_2, x_3$ thỏa $-1< x_1< 0< x_2<  1< x_3< 2$.
		\begin{itemchoice}
			\itemch  \textbf{Sai.} \\
			Phương trình  không có nghiệm trong khoảng $(-1; 1)$.
			\itemch \textbf{Sai.}  \\
			Phương  trình không có nghiệm trong khoảng $(-2; 0)$.
			 \itemch \textbf{Sai.} \\
			Phương trình chỉ có một  nghiệm trong khoảng $(-2; 1)$.
			 \itemch \textbf{Đúng.} \\
			Phương  trình có ít nhất hai nghiệm trong khoảng $(0; 2)$.
		\end{itemchoice}
	 }
\end{ex}

% nguồn: Ok_dungsai_lop11_Lan2
\begin{ex}%[1D3H3-4]%[ma:d4e01c97]
	[Dự án 11HVA  2024-2025]
	Cho hàm số $f(x)=\heva{&\dfrac{\sqrt{1+2x}-1}{x} \text{khi  } x > 0\\
	&3x+m+1\text{khi } x \leq 0}$, $m$ là tham số. 
	\choiceTF 
		{\True Hàm số $f(x)$ liên  tục tại $x=-3$}
		{Hàm  số $f(x)$ gián đoạn tại $x=1$}
		{\True Khi $m=1$, hàm số $f(x)$ gián  đoạn tại $x=0$}
		{\True Hàm số $f(x)$  liên tục trên $\mathbb{R}$ khi $m=0$}
	\loigiai{Hàm số $y=f(x)$ xác định trên $\mathbb{R}$.\\ 
		Với $x  > 0$ ta có hàm số $f(x)=\dfrac{\sqrt{1+2x}-1}{x}$ liên tục trên khoảng $(0;+\infty)$.\\
		Với  $x < 0$ ta có $f(x)=3x+m+1$ liên tục trên khoảng $(-\infty; 0)$.\\
		Với  $x=0$ ta có $f(0)=1$.\\
		$\lim\limits_{x \to 0^{-}} f(x)=\lim\limits_{x \to  0^{-}}(3x+m+1)=m+1$.\\
		$\lim\limits_{x \to 0^{+}} f(x)=\lim\limits_{x \to  0^{+}}\left(\dfrac{\sqrt{1+2x}-1}{x}\right)=\lim\limits_{x \to 0^{+}}\left(\dfrac{2x}{x(\sqrt{1+2x}+1)}\right)=\lim\limits_{x \to 0^{+}}\left(\dfrac{2}{(\sqrt{1+2x}+1)}\right)=1$.\\
		Hàm số $f(x)$ liên tục trên $\mathbb{R}$ khi hàm số $f(x)$ liên tục tại $x=0$ 
		$\Leftrightarrow \lim\limits_{x  \to 0^{-}} f(x)=\lim\limits_{x \to 0^{+}} f(x)=f(0) \Leftrightarrow m+1=1\Leftrightarrow m=0$.\\
		Vậy hàm số $f(x)$ liên  tục trên $\mathbb{R}$ khi $m=0$.
		\begin{itemchoice}
			\itemch \textbf{Đúng.} 
			Vì với $x < 0$ ta có $f(x)=3x+m+1$ liên tục  trên khoảng $(-\infty; 0),-3\in(-\infty; 0)$.
			 \itemch \textbf{Sai.} \\
			Vì với $x > 0$ ta có hàm số $f(x)=\dfrac{\sqrt{1+2x}-1}{x}$ liên tục trên khoảng $(0;+\infty)$  nên $f(x)$ liên tục tại $x=1$.
			 \itemch \textbf{Đúng.} \\
			Vì hàm số $f(x)$ gián đoạn tại  $x=0$ khi $m+1\neq 1\Leftrightarrow m \neq 0$.
			\itemch  \textbf{Đúng.} \\
		\end{itemchoice}
	} 
\end{ex}

% nguồn: Ok_dungsai_lop11_Lan2
\begin{ex}%[1D3H3-4]%[ma:c2fb7bf9]
	[Dự án  11HVA 2024-2025]
	Cho hàm số $f(x)=\heva{&\dfrac{\sqrt{x+4}-2}{x}  \text{khi } x > 0\\
	&m x^2+2m+\dfrac{1}{4} \text{khi } x \leq 0}$, với $m$ là tham số. 
	\choiceTF 
		{Tập xác định của  hàm số $f(x)$ là $\mathscr{D}=[-4;+\infty)$}
		{\True  $f(0)=2m+\dfrac{1}{4}$}
		 {$\lim\limits_{x \to 0^{-}} f(x)=\dfrac{1}{4}$}
		{Hàm số $f(x)$  liên tục trên $\mathbb{R}$ khi $m=\dfrac{1}{4}$}
	\loigiai{\begin{itemchoice}
		\itemch  \textbf{Sai.} \\
		Tập xác định: $\mathscr{D}=\mathbb{R}$. 
		\itemch \textbf{Đúng.} \\ 
		Ta có: 
		 $f(0)=2m+\dfrac{1}{4}\cdot$
		\itemch \textbf{Sai.}  \\
		Ta có $\lim\limits_{x \to 0^{-}} f(x)=\lim\limits_{x  \to 0^{-}}\left(m x^2+2m+\dfrac{1}{4}\right)=2m+\dfrac{1}{4}\cdot$
		\itemch \textbf{Sai.}  \\
		Ta có $\lim\limits_{x \to 0^{+}} f(x)=\lim\limits_{x \to 0^{+}} \dfrac{\sqrt{x+4}-2}{x}=\lim\limits_{x  \to 0^{+}} \dfrac{x+4-2^2}{x(\sqrt{x+4}+2)}=\lim\limits_{x \to 0^{+}}\left(\dfrac{1}{\sqrt{x+4}+2}\right)=\dfrac{1}{4}\cdot$\\
		Hàm số $f(x)$ liên tục trên các khoảng  $(-\infty; 0)$ và $(0;+\infty)$.\\
		 Nên để hàm số $f(x)$ liên tục trên $\mathbb{R}$, khi hàm số $f(x)$ liên tục tại $x=0$.\\
		Mà hàm số $f(x)$ liên tục tại  $x=0$ khi:\\
		$f(0)=\lim\limits_{x \to 0^{-}} f(x)=\lim\limits_{x \to 0^{+}} f(x) \Leftrightarrow 2m+\dfrac{1}{4}=\dfrac{1}{4}  \Leftrightarrow m=0$.
	\end{itemchoice}
 }
\end{ex}

% nguồn: Ok_dungsai_lop11_Lan2
\begin{ex}%[1D3H3-4]%[ma:485e156b]
	 [Dự án 11HVA 2024-2025]
	Cho hàm số $f(x)=\heva{&-\dfrac{x}{2} \text{khi  } x \leq 1\\
	&\dfrac{x^2-3x+2}{x^2-1} \text{khi } x >  1}$ và $g(x)=x^2-3x+1$. Khi đó:
	 \choiceTF
		 {Giới hạn $\lim\limits_{x \to 1^{+}} f(x)=\dfrac{1}{2}$}
		{\True  Hàm số $f(x)$ liên tục tại điểm $x_0=1$}
		{\True  Hàm số $g(x)$ liên tục tại điểm $x_0=1$}
		{\True Hàm số $y=f(x)+g(x)$  liên tục tại điểm $x_0=1$}
	\loigiai{\begin{itemchoice}
		\itemch \textbf{Sai.} \\ 
		Ta có $\lim\limits_{x \to 1^{+}}  f(x)=\lim\limits_{x \to 1^{+}} \dfrac{x^2-3x+2}{x^2-1}=\lim\limits_{x \to 1^{+}} \dfrac{(x-1)(x-2)}{(x-1)(x+1)}=\lim\limits_{x \to 1^{+}} \dfrac{x-2}{x+1}=-\dfrac{1}{2}\cdot$
		\itemch  \textbf{Đúng.} \\
		Ta có $f(1)=-\dfrac{1}{2}$; $\lim\limits_{x  \to 1^{-}} f(x)=\lim\limits_{x \to 1^{-}} \dfrac{-x}{2}=-\dfrac{1}{2}$; $\lim\limits_{x \to 1^{+}} f(x)=-\dfrac{1}{2}$.\\
		Vậy $f(1)=\lim\limits_{x \to 1} f(x)=-\dfrac{1}{2}$ nên hàm số $f(x)$ liên tục  tại điểm $x_0=1$.
		\itemch \textbf{Đúng.} \\ 
		Ta có $g(1)=-1$ và $\lim\limits_{x \to 1} g(x)=1^2-3\cdot  1+1=-1$ nên $g(1)=\lim\limits_{x \to 1} g(x)$.\\
		Vậy hàm số $g(x)$ liên  tục tại điểm $x_0=1$.
		\itemch \textbf{Đúng.} \\ 
		Ta có hàm số $f(x)$, $g(x)$ liên tục tại điểm $x_0=1$ nên hàm số  $y=f(x)+g(x)$ liên tục tại điểm $x_0=1$
	\end{itemchoice}
} 
\end{ex}

% nguồn: Ok_dungsai_lop11_Lan2
\begin{ex}%[1D3H3-4]%[ma:24a8b1ab]
	 [Dự án 11HVA 2024-2025]
	Cho hàm số $f(x)=-4x^3+4x-1$. Xét tính đúng sai của các mệnh đề  sau
	 \choiceTF
		{\True Hàm số đã cho liên tục  trên $\mathbb{R}$}
		{Phương trình $f(x)=0$ không có nghiệm trên khoảng $(-\infty;  1)$}
		{\True Phương  trình $f(x)=0$ có nghiệm trên khoảng $(-2; 0)$}
		{\True Phương trình $f(x)=0$ có ít nhất  hai nghiệm trên khoảng $\left(-3; \dfrac{1}{2}\right)$}
	\loigiai{Hàm $f(x)$ là hàm đa thức  nên liên tục trên $\mathbb{R}$.\\
		Ta có $\heva{&f(-1)=-1< 0\\ 
		&f(-2)=23>  0} \to f(x)=0$ có nghiệm $x_1$ trên $(-2;-1)$.\\
		Mà $(-2;-1) \subset(-2; 0) \subset(-\infty; 1)$. 
		 Ta có $\heva{&f(0)=-1< 0\\
		&f\left(\dfrac{1}{2}\right)=\dfrac{1}{2}  > 0} \to f(x)=0$ có nghiệm $x_2$ thuộc $\left(0; \dfrac{1}{2}\right)$. Kết hợp với (1) suy ra $f(x)=0$ có các nghiệm $x_1, x_2$ thỏa: $-3< x_1<-1< 0< x_2< \dfrac{1}{2}$.
		\begin{itemchoice}
			\itemch  \textbf{Đúng.} \\
			 Hàm số đã cho liên tục trên $\mathbb{R}$.
			\itemch  \textbf{Sai.} \\
			Phương trình $f(x)=0$ không  có nghiệm trên khoảng $(-\infty; 1)$.
			 \itemch \textbf{Đúng.} \\
			Phương trình $f(x)=0$ có nghiệm trên khoảng ($-2; 0$). 
			\itemch  \textbf{Đúng.} \\
			Phương trình $f(x)=0$ có ít nhất  hai nghiệm trên khoảng $\left(-3; \dfrac{1}{2}\right)$.
		\end{itemchoice}
	} 
\end{ex}

% nguồn: Ok_dungsai_lop11_Lan2
\begin{ex}%[1D3H3-4]%[ma:c0fcef3d]
	 [Dự án 11HVA 2024-2025]
	Cho  các hàm số $f(x)=\heva{&\dfrac{\sqrt{4x-7}-1}{x^2-4} \text{khi } x > 2\\
	&\dfrac{5x-9}{2} \text{khi }  x \leq 2}$ và $g(x)=\heva{&\dfrac{\sqrt{x+2}-2}{2-x} \text{khi } x > 2\\
	&\dfrac{1-x}{4} \text{khi } x \leq  2\cdot}$
	 \choiceTF
		 {\True Hàm số $f(x)$ liên tục tại điểm $x_0=2$}
		{Giới hạn $\lim\limits_{x \to 2^{+}} g(x)=\dfrac{1}{4}$} 
		{Hàm số  $g(x)$ gián đoạn tại điểm $x_0=2$}
		{\True Hàm số $y=\dfrac{f(x)}{g(x)}$ liên tục tại điểm $x_0=2$} 
	\loigiai{\begin{itemchoice}
		 \itemch \textbf{Đúng.} \\
		Ta  có $f(x_0)=f(2)=\dfrac{1}{2}=\lim\limits_{x \to 2^{-}} f(x)$.\\
		$\lim\limits_{x \to x_0^{+}} f(x)=\lim\limits_{x \to 2^{+}} \dfrac{\sqrt{4x-7}-1}{x^2-4}=\dfrac{1}{2}=\lim\limits_{x \to  2^{-}} f(x) \Rightarrow \lim\limits_{x \to 2^{-}} f(x)=\lim\limits_{x \to 2^{+}} f(x)=\dfrac{1}{2}=\lim\limits_{x \to 2} f(x)=f(2)$.\\
		Vậy hàm số $f(x)$ liên tục  tại điểm $x_0=2$.
		\itemch \textbf{Sai.}  \\
		Ta có $\lim\limits_{x \to 2^{+}} g(x)=\lim\limits_{x  \to 2^{+}}\left(\dfrac{\sqrt{x+2}-2}{2-x}\right)=\lim\limits_{x \to 2^{+}} \dfrac{x+2-4}{(2-x)(\sqrt{x+2}+2)}=\lim\limits_{x \to 2^{+}} \dfrac{-1}{\sqrt{x+2}+2}=-\dfrac{1}{4}\cdot$
		 \itemch \textbf{Sai.} \\
		Ta có $g(2)=\dfrac{1-2}{4}=-\dfrac{1}{4}; \lim\limits_{x \to 2^{-}} g(x)=\lim\limits_{x \to 2^{-}}\left(\dfrac{1-x}{4}\right)=-\dfrac{1}{4}$; $\lim\limits_{x \to 2^{+}} g(x)=-\dfrac{1}{4}\cdot$\\ 
		Suy ra $\lim\limits_{x \to 2} g(x)=-\dfrac{1}{4}=g(2)$.\\ 
		Vậy hàm số $g(x)$ liên  tục tại điểm $x_0=2$.
		\itemch \textbf{Đúng.} \\ 
		Ta có hàm số $f(x)$, $g(x)$ liên tục tại điểm $x_0=2$ và $g(2)=-\dfrac{1}{4}\neq 0$ do đó hàm số $y=\dfrac{f(x)}{g(x)}$ liên  tục tại điểm $x_0=2$.
	\end{itemchoice}
} 
\end{ex}

% nguồn: Ok_dungsai_lop11_Lan2
\begin{ex}%[1D3H3-4]%[ma:e6b0e57c]
	[Dự  án 11HVA 2024-2025]
	Cho hàm số $y=f(x)$ là hàm  đa thức thoả điều kiện $f(-1)=-2$, $f(1)=2$, $f(5)=10$.
	 \choiceTF
		{\True Phương trình $f(x)=0$ có ít nhất một  nghiệm trên khoảng $(-1; 1)$}
		{Phương trình $f(x)=0$ có ít nhất một  nghiệm trên khoảng (1;5)}
		{\True Phương trình $f(x)=7$ có ít  nhất một nghiệm trên khoảng $(1; 5)$}
		{\True Phương trình $m[f(x)-1](x+1)-m=0$ luôn có nghiệm với mọi giá trị  của tham số $m \neq 0$}
	\loigiai{\begin{itemchoice}
		 \itemch \textbf{Đúng.} \\
		Ta có  hàm số $y=f(x)$ liên tục trên đoạn $[-1; 1]$ và $f(-1) \cdot f(1)=-2\cdot 2=-4< 0$ nên phương trình $f(x)=0$ có ít nhất một nghiệm trên khoảng $(-1; 1)$.
		 \itemch \textbf{Sai.} \\
		Ta có hàm số $y=f(x)$ liên tục trên đoạn $[1;  5]$ và $f(1) \cdot f(5)=2\cdot 10=20> 0$ nên chưa thể kết luận được phương trình $f(x)=0$ có ít nhất một nghiệm trên khoảng $(1; 5)$.
		\itemch \textbf{Đúng.}  \\
		Do hàm số $y=f(x)$ liên tục trên đoạn $[1; 5]$ và $f(1)=2, f(5)=10$ nên phương trình $f(x)=a$ với $a \in(2; 10)$ luôn có ít nhất một nghiệm thuộc khoảng $(1; 5)$. Vì $7\in(2; 10)$ nên phương trình $f(x)=7$ có  ít nhất một nghiệm trên khoảng $(1; 5)$.
		\itemch \textbf{Đúng.} \\ 
		Xét hàm số $h(x)=m[f(x)-1](x+1)-m$ liên tục trên đoạn $[-1; 1]$ có:\\ 
		$h(1)=m[f(1)-1](1+1)-m=m[2-1](1+1)-m=m;$ \\ 
		$h(-1)=m[f(-1)-1](-1+1)-m=-m$;\\ 
		Suy ra $h(1)\cdot h(-1)=-m^2< 0$ với mọi giá trị của tham số $m \neq  0$.
		Vậy phương trình $m[f(x)-1](x+1)-m=0$ luôn có ít nhất một nghiệm thuộc khoảng  $(-1; 1)$ với mọi $m \neq 0$.
	\end{itemchoice}
 }
\end{ex}

% nguồn: Ok_dungsai_lop11_Lan2
\begin{ex}%[1D3H3-4]%[ma:32336825]
	[Dự án 11HVA  2024-2025]
	Cho các hàm số sau: $f(x)=\heva{&-\dfrac{x}{2} \text{khi } x \leq  1\\
	 &[7pt] \dfrac{x^2-3x+2}{x^2-1} \text{khi } x > 1}, g(x)=x^2-3x+1$ và $h(x)=\sin \dfrac{\pi x}{4}$
	\choiceTF 
		{\True Hàm số  $f(x)$ liên tục tại điểm $x_0=1$}
		{\True Hàm số $g(x)$ liên tục tại  điểm $x_0=1$}
		{Hàm số $h(x)$ không liên  tục tại điểm $x_0=2$}
		{Hàm số $y=f(x) \cdot g(x)$ không liên  tục tại điểm $x_0=1$}
	\loigiai{\begin{itemchoice}
		\itemch  \textbf{Đúng.} \\
		Ta có $f(1)=-\dfrac{1}{2}$ và $\lim\limits_{x \to 1^{-}}  f(x)=\lim\limits_{x \to 1^{-}} \dfrac{-x}{2}=-\dfrac{1}{2}\cdot$\\
		$\lim\limits_{x \to 1^{+}} f(x)=\lim\limits_{x \to 1^{+}} \dfrac{x^2-3x+2}{x^2-1}=\lim\limits_{x \to  1^{+}} \dfrac{(x-1)(x-2)}{(x-1)(x+1)}=\lim\limits_{x \to 1^{+}} \dfrac{x-2}{x+1}=-\dfrac{1}{2}\cdot$\\
		Vậy $f(1)=\lim\limits_{x \to 1} f(x)=-\dfrac{1}{2}$ nên  hàm số $f(x)$ liên tục tại điểm $x_0=1$.
		\itemch \textbf{Đúng.}  \\
		Ta có $g(1)=-1$ và $\lim\limits_{x  \to 1} g(x)=1^2-3\cdot 1+1=-1$ nên $g(1)=\lim\limits_{x \to 1} g(x)$.\\
		Vậy hàm  số $g(x)$ liên tục tại điểm $x_0=1$.
		\itemch  \textbf{Sai.} \\
		Ta có $h(2)=\sin \dfrac{2\pi}{4}=1$ và $\lim\limits_{x \to 2} h(x)=\sin \dfrac{2\pi}{4}=1$ nên $h(2)=\lim\limits_{x \to 2} h(x)$.\\ 
		Vậy hàm số  $h(x)$ liên tục tại điểm $x_0=2$.
		 \itemch \textbf{Sai.} \\
		Hàm số $f(x)$ và  $g(x)$ liên tục tại điểm $x_0=1$ nên hàm số $y=f(x) \cdot g(x)$ liên tục tại điểm $x_0=1$.
	\end{itemchoice}
 }
\end{ex}

% nguồn: Ok_dungsai_lop11_Lan2
\begin{ex}%[1D3H3-4]%[ma:dc72f61e]
	[Dự án  11HVA 2024-2025]
	Chi phí sản xuất $C$ phụ thuộc lượng đầu ra  của sản phẩm $Q$, chẳng hạn:
	\begin{align*}
		C=f(Q)=\begin{cases}Q+2& \text{khi } 0< Q\leq 3\\ \dfrac{Q^2-2Q-3}{Q-3} & \text{khi } Q> 3.
	\end{cases}
\end{align*}
\choiceTF 
	{$\lim\limits_{Q\to 0^{+}}  f(Q)$ không tồn tại}
	 {\True $\lim\limits_{Q\to 3^{-}} f(Q)=5$}
	{$\lim\limits_{Q\to  3^{+}} f(Q)=1$}
	{Hàm số $f(Q)$ liên tục  trên $(0;+\infty)$}
\loigiai{\begin{itemchoice}
	\itemch  \textbf{Sai.} \\
	Hàm số $f(Q)$ xác định trong  $(0;+\infty)$. \\
	Ta có $\lim\limits_{Q\to 0^{+}} f(Q)=\lim\limits_{Q\to 0^{+}}(Q+2)=2$. 
	\itemch \textbf{Đúng.} \\ 
	Ta có 
	$\lim\limits_{Q\to 3^{-}} f(Q)=\lim\limits_{Q\to  3^{-}}(Q+2)=5$.
	 \itemch \textbf{Sai.} \\
	Ta có $\lim\limits_{Q\to 3^{+}} f(Q)=\lim\limits_{Q\to 3^{+}} \dfrac{Q^2-2Q-3}{Q-3}=\lim\limits_{Q\to 3^{+}} \dfrac{(Q+1)(Q-3)}{Q-3}=\lim\limits_{Q\to 3^{+}}(Q+1)=4$. 
	\itemch \textbf{Sai.}  \\
	Do $\lim\limits_{Q\to 3^{-}} f(Q) \neq \lim\limits_{Q\to 3^{+}} f(Q)$ nên hàm số $f(Q)$ gián đoạn tại $Q=3$. Suy ra hàm  số không liên tục trên $(0;+\infty)$.
	\end{itemchoice}
} 
\end{ex}

% nguồn: Ok_dungsai_lop11_Lan1
\begin{ex}%[1D3V3-4]%[ma:57cdf2b8]
 Cho hàm số $f(x)=\heva{&\dfrac{2-\sqrt{x+5}}{x^2-5x-4}&\text{khi}&\,\,x>-1\\&x^2-9x&\text{khi}&\,\,x\le -1}$  và $g(x)=\heva{&\dfrac{x^2-1}{x+1}&\text{khi}&\,\,x\ne-1\\&2a+1&\text{khi}&\,\,x=-1}$. Khi đó:
  \choiceTF
  {\True Ta có  $\lim\limits_{x\to {-1}^{+}}f(x)=\dfrac{1}{8}$}
  {\True Hàm  số $f(x)$ liên tục trên khoảng $(-\infty;-1)$}
  {Hàm số  $g(x)$ liên tục tại điểm $x_0=-1$ thì $a=\dfrac{1}{2}$}
  { Khi $a=-\dfrac{1}{2}$ thì hàm  số $y=f(x)-g(x)$ liên tục tại điểm $x_0=-1$}
 \loigiai{ 
  \begin{itemchoice}
   \itemch Ta có 
   \begin{eqnarray*}
     \lim\limits_{x\to{-1}^{+}}f(x)&=&\lim\limits_{x\to{-1}^{+}}\dfrac{2-\sqrt{x+5}}{x^2-5x-4}\\
     &=& \lim\limits_{x\to{-1}^{+}}\left[-\dfrac{1}{(x-4)\left(2+\sqrt{x+5} \right) } \right]\\
     &=&\dfrac{1}{8}
   \end{eqnarray*}
    \itemch
   Xét $x_0\in(-\infty;-1)$, ta có $f(x_0)=\lim\limits_{x\to  x_0}f(x)=x_0^2-9x_0$.\\
   Vậy hàm số $f(x)$ liên tục trên  khoảng $(-\infty;-1)$.
   \itemch Ta có $\heva{&g(-1)=2a+1\\&\lim\limits_{x\to {-1}}g(x)=\lim\limits_{x\to {-1}}\dfrac{x^2-1}{x-1}=\lim\limits_{x\to {-1}}(x+1)=0.}$\\ 
   Để  hàm số $g(x)$ liên tục tại điểm $x_0=-1\Leftrightarrow 2a+1=0\Leftrightarrow a=-\dfrac{1}{2}$.
   \itemch 
   Ta  có $\heva{&f(-1)=\lim\limits_{x\to {-1}^{-}}f(x)=10\\&\lim\limits_{x\to {-1}^{+}}f(x)=\dfrac{1}{8}}$.\\
   Vậy hàm số  $f(x)$ gián đoạn tại $x_0=-1$.\\
   Do hàm số $f(x)$ gián đoạn tại $x_0=-1$ nên hàm số $y=f(x)-g(x)$ cũng gián  đoạn tại điểm $x_0=-1$.
  \end{itemchoice}
  }
\end{ex}

% nguồn: Ok_dungsai_lop11_Lan1
\begin{ex}%[1D3V3-4]%[ma:5079fd45]
 Cho các hàm số $f(x)=\heva{&\dfrac{x^2-4}{x-2}&\text{khi}&\,\,x\ne 2\\&4{,}5&\text{khi}&\,\,x=2}$ và $g(x)=\dfrac{2}{x-1}$. Khi đó. 
  \choiceTF
  {\True Hàm số $g(x)$ liên  tục trên khoảng $(-\infty;1)$}
  {\True Giới hạn $\lim\limits_{x\to  2}f(x)=4$}
  {Hàm  số $f(x)$ liên tục tại điểm $x_0=2$}
   {Hàm số $\dfrac{f(x)}{g(x)}$ liên tục tại điểm $x_0=2$}
  \loigiai{
  \begin{itemchoice}
   \itemch 
   Xét $x_0\in (-\infty;1)$, ta có  $g(x_0)=\lim\limits_{x\to x_0}\dfrac{2}{x-1}=\dfrac{2}{x_0-1}$\\
   Vậy hàm số $g(x)$ liên  tục trên khoảng $(-\infty;1)$.
   \itemch Ta có $\lim\limits_{x\to 2}f(x)=\lim\limits_{x\to  2}\dfrac{x^2-4}{x-2}=\lim\limits_{x\to 2}(x+2)=4$.
   \itemch Ta  có $\heva{&f(2)=4{,}5\\&\lim\limits_{x\to 2}f(x)=\lim\limits_{x\to 2}\dfrac{x^2-4}{x-2}=\lim\limits_{x\to 2}(x+2)=4.}$\\
   Vậy hàm  số $f(x)$ gián đoạn tại điểm $x_0=2$.
   \itemch Do hàm số $f(x)$ gián  đoạn tại điểm $x_0=2$ nên hàm số $\dfrac{f(x)}{g(x)}$ cũng gián đoạn tại điểm $x_0=2$.
  \end{itemchoice}
 } 
\end{ex}

% nguồn: Ok_dungsai_lop11_Lan1
\begin{ex}%[1D3V3-4]%[ma:0e06e9b8]
 Cho hàm  số $f(x)=\heva{&\dfrac{1-\sqrt{5x+11}}{2x^2-5x-18}&\text{khi}&\,\,x>-2\\&4-x^2&\text{khi}&\,\,x\le -2}$ và $g(x)=\heva{&\dfrac{x^2-x-6}{x+2}&\text{khi}&\,\,x\ne-2\\&2x+a&\text{khi}&\,\,x=-2}$. Khi đó:
 \choiceTF 
  {\True Hàm số $f(x)$ liên  tục trên khoảng $(-\infty;-2)$}
  {Hàm số $f(x)$ liên tục tại điểm  $x_0=-2$}
  {Để hàm số $g(x)$ liên tục tại điểm  $x_0=-2$ thì $a=1$}
  {\True  Khi $a=-1$ thì hàm số $y=f(x)\cdot g(x)$ gián đoạn tại điểm $x_0=-2$}
  \loigiai{
  \begin{itemchoice}
    \itemch Xét $x_0\in (-\infty;-2)$, ta có $f(x_0)=\lim\limits_{x\to x_0}f(x)=\lim\limits_{x\to x_0}(4-x^2)=4-x_0^2$.\\
   Vậy hàm số  $f(x)$ liên tục trên khoảng $(-\infty;-2)$.
   \itemch 
   Ta có $f(-2)=\lim\limits_{x\to {-2}^{-}}f(x)=\lim\limits_{x\to {-2}^{-}}(4-x^2)=0$.\\ 
    Mặt khác
   \begin{eqnarray*}
     \lim\limits_{x\to{-2}^{+}}f(x)&=&\lim\limits_{x\to{-2}^{+}}\dfrac{1-\sqrt{5x+11}}{2x^2-5x-18}\\
    &=& \lim\limits_{x\to{-2}^{+}}\left[-\dfrac{5}{(2x-9)\left(1+\sqrt{5x+11} \right)  } \right]\\
     &=&\dfrac{5}{26}
   \end{eqnarray*}
   Vậy hàm số $f(x)$ gián  đoạn tại điểm $x_0=-2$.
   \itemch Ta có $\heva{&g(-2)=-4+a\\&\lim\limits_{x\to {-2}}g(x)=\lim\limits_{x\to {-2}}\dfrac{x^2-x-6}{x+2}=\lim\limits_{x\to  {-2}}(x-3)=-5.}$\\
   Để hàm  số $g(x)$ liên tục tại điểm $x_0=-2\Leftrightarrow -4+a=-5\Leftrightarrow a=-1$.
   \itemch 
   Do hàm số $f(x)$  gián đoạn tại $x_0=-2$ nên hàm số $y=f(x)\cdot g(x)$ cũng gián đoạn tại điểm $x_0=-2$.
  \end{itemchoice}
 } 
\end{ex}

% nguồn: Ok_dungsai_lop11_Lan1
\begin{ex}%[1D3V3-4]%[ma:b65ffd4b]
 Cho các hàm số $f(x)=\heva{&\dfrac{\sqrt{4x-7}-1}{x^2-4}&\text{khi}&\,\,x> 2\\&\dfrac{5x-9}{2}&\text{khi}&\,\,x\le 2}$ và $g(x)=\heva{&\dfrac{\sqrt{x+2}-2}{2-x}&\text{khi}&\,\,x> 2\\&\dfrac{1-x}{4}&\text{khi}&\,\,x\le 2}$. Khi  đó
  \choiceTF
  {\True Hàm số $f(x)$ liên  tục tại điểm $x_0=2$}
  {Hàm số $g(x)$  không liên tục trên khoảng $(2;+\infty)$}
   {Giới hạn $\lim\limits_{x\to 2^{+}}g(x)=\dfrac{1}{4}$}
  {\True Hàm số $y=\dfrac{f(x)}{g(x)}$  liên tục tại điểm $x_0=2$}
  \loigiai{
  \begin{itemchoice}
   \itemch Ta có Ta có $\heva{&f(2)=\dfrac{1}{2}\\&\lim\limits_{x\to 2^{-}}f(x)=\lim\limits_{x\to 2^{-}}\left(\dfrac{5x-9}{2}\right)=\dfrac{1}{2}\\&\lim\limits_{x\to 2^{+}}f(x)=\lim\limits_{x\to 2^{+}}\dfrac{\sqrt{4x-7}-1}{x^2-4}=\lim\limits_{x\to 2^{+}}\dfrac{4}{(x+2)\left(\sqrt{4x-7}+1 \right) }=\dfrac{1}{2}.}$\\ 
   Vậy hàm số $f(x)$ liên tục tại  $x_0=2$.
    \itemch
   Xét $x_0\in  (2;+\infty)$,
    ta có\\ $g(x_0)=\lim\limits_{x\to x_0}g(x)=\lim\limits_{x\to x_0}\dfrac{\sqrt{x+2}-2}{2-x}=\lim\limits_{x\to x_0}\left( -\dfrac{1}{\sqrt{x+2}+2}\right)=-\dfrac{1}{\sqrt{x_0+2}+2}$. \\
   Vậy hàm  số $g(x)$ liên tục trên khoảng $(2;+\infty)$.
   \itemch Ta có $\lim\limits_{x\to 2^{+}}g(x)=\lim\limits_{x\to 2^{+}}\dfrac{\sqrt{x+2}-2}{2-x}=\lim\limits_{x\to  2^{+}}\left( -\dfrac{1}{\sqrt{x+2}+2}\right) =-\dfrac{1}{4}$.
    \itemch
    Ta có $g(2)=\lim\limits_{x\to 2^{-}}g(x)=\lim\limits_{x\to 2^{-}}\dfrac{1-x}{4}=-\dfrac{1}{4}$.\\
   Do đó hàm số $g(x)$  liên tục tại $x_0=2$.\\
   Do các hàm số  $f(x)$ và $g(x)$ đều liên tục tại điểm $x_0=2$ và $g(2)=-\dfrac{1}{4}\ne 0$ nên hàm số $y=\dfrac{f(x)}{g(x)}$ cũng liên tục tại điểm $x_0=2$.
  \end{itemchoice}
 } 
\end{ex}


%%%%%%%%%%%%%%%% TRẢ LỜI NGẮN %%%%%%%%%%%%%%%%

% nguồn: Ok_TraLoiNgan_Lop11_Lan1
\begin{ex}%[1D3H3-4]%[ma:9764eb76]
	Tìm giá trị của tham số $m$ để hàm số  $f(x)=\heva{&\dfrac{x^2-16}{x-4}&\text{khi}&\,\,x>4\\&mx+1&\text{khi}&\,\,x\le 4}$ liên tục trên $\mathbb{R}$.
	\shortans{1,75} 
	 \loigiai{
		Tập  xác định $\mathscr{D}=\mathbb{R}$.\\
		 Xét $f(x)=\dfrac{x^2-16}{x-4}$ là hàm phân thức nên liên tục trên khoảng $x\in (4;+\infty)$.\\
		$f(x)=mx+1$ là hàm đa thức nên liên  tục trên khoảng $x\in (-\infty;4)$.\\
		Vậy $f(x)$ trên $\mathbb{R}$  khi và chỉ khi $f(4)=\lim\limits_{x\to 4^{-}}f(x)=\lim\limits_{x\to 4^{+}}f(x)\quad (1)$\\
		Ta có $\heva{&f(4)=\lim\limits_{x\to 4^{-}}f(x)=4m+1\\&\lim\limits_{x\to 4^{+}}f(x)=\lim\limits_{x\to  4^{+}}\dfrac{x^2-16}{x-4}=\lim\limits_{x\to 4^{+}}(x+4)=8.}$\\
		Hàm số $f(x)$ liên tục trên $\mathbb{R}$ khi  và chỉ khi $f(4)=\lim\limits_{x\to 4^{+}}f(x)=\lim\limits_{x\to 4^{-}}f(x)$\\
		$\Leftrightarrow  4m+1=8\Leftrightarrow m=\dfrac{7}{4}=1{,}75$.
	} 
\end{ex}

% nguồn: Ok_TraLoiNgan_Lop11_Lan1
\begin{ex}%[1D3H3-4]%[ma:ba9984ad]
	Tìm giá  trị của tham số $m$ để hàm số $f(x)=\heva{&\dfrac{x^2-x-2}{x-2}&\text{khi}&\,\,x\ne 2\\&m&\text{khi}&\,\,x=2}$ liên tục trên $\mathbb{R}$.
	\shortans{3} 
	 \loigiai{
		Tập xác định $\mathscr{D}=\mathbb{R}$.\\ 
		Hàm số $f(x)$ liên tục trên  các khoảng $(-\infty;2)$, $(2;+\infty)$.\\
		Ta có $\heva{&f(2)=m\\&\lim\limits_{x\to 2}f(x)=\lim\limits_{x\to 2}\dfrac{x^2-x-2}{x-2}=\lim\limits_{x\to  2}(x+1)=3.}$\\
		Hàm số $f(x)$ liên tục trên $\mathbb{R}$ khi  và chỉ khi $f(2)=\lim\limits_{x\to 2}f(x)\Leftrightarrow m=3$.
	 }
\end{ex}

% nguồn: Ok_TraLoiNgan_Lop11_Lan1
\begin{ex}%[1D3H3-4]%[ma:873c72d8]
	Tìm giá trị của tham số $m$ để hàm số $f(x)=\heva{&x^2+x&\text{khi}&\,\,x<1\\&2&\text{khi}&\,\,x=1\\&mx+1&\text{khi}&\,\,x>1}$ liên tục trên  $\mathbb{R}$.
	\shortans{1} 
	 \loigiai{
		Tập xác  định $\mathscr{D}=\mathbb{R}$.\\
		Hàm số $f(x)$ liên tục  trên các khoảng $(-\infty;1)$, $(1;+\infty)$.\\
		Ta có $\heva{&f(1)=2\\&\lim\limits_{x\to 1^{-}}f(x)=\lim\limits_{x\to 1^{-}}(x^2+x)=2\\&\lim\limits_{x\to 1^{+}}f(x)=\lim\limits_{x\to  1^{+}}(mx+1)=m+1.}$\\
		Hàm số $f(x)$ liên tục trên $\mathbb{R}$  khi và chỉ khi $f(1)=\lim\limits_{x\to 1^{-}}f(x)=\lim\limits_{x\to 1^{+}}f(x)$.\\
		$\Leftrightarrow m+1=2\Leftrightarrow m=1$. 
	} 
\end{ex}

% nguồn: Ok_TraLoiNgan_Lop11_Lan1
\begin{ex}%[1D3H3-4]%[ma:65ee43c5]
	Tìm giá trị của tham số $a$  để hàm số $f(x)=\heva{&x^2+2x+2&\text{khi}&\,\,x\ge 0\\&x^2+a&\text{khi}&\,\,x<0}$ liên tục trên $\mathbb{R}$.
	\shortans{2} 
	 \loigiai{
		 Tập xác định $\mathscr{D}=\mathbb{R}$.\\
		Hàm số $f(x)$ liên tục trên các khoảng $(-\infty;0)$,  $(0;+\infty)$.\\
		Ta có $\heva{&f(0)=\lim\limits_{x\to 0^{+}}f(x)=2\\&\lim\limits_{x\to 0^{-}}f(x)=\lim\limits_{x\to 0^{-}}(x^2+a)=a.}$\\ 
		Hàm số $f(x)$ liên tục trên $\mathbb{R}$ khi và chỉ khi  $f(0)=\lim\limits_{x\to 0^{+}}f(x)=\lim\limits_{x\to 0^{-}}f(x)\Leftrightarrow a=2$.
	 }
\end{ex}

% nguồn: Ok_TraLoiNgan_Lop11_Lan2
\begin{ex}%[1D3H3-4]%[ma:f3b2e84c]
	 Biết hàm số $f(x)=\heva{&{ax^2+bx-4}{\text{khi}} {x\le 1} \\
	&{2ax-2b}{\text{khi}}  {x > 1}.
	}$ liên tục trên $\mathbb{R}$. Tính giá trị của  biểu thức $P=a-3b$.
	 \shortans{-4}
	\loigiai{Tập xác đinh:  $\mathscr{D}=\mathbb{R}$.\\
		Ta có: $f(1)=\lim\limits_{x\to 1^{-}} f(x)=\lim\limits_{x\to  1^{-}} \left(ax^2+bx-4\right)=a+b-4$.\\
		Mặt khác: $ \lim\limits_{x\to 1^{+}} f(x)=\lim\limits_{x\to  1^{+}} \left(2ax-2b\right)=2a-2b$.\\
		Để hàm số liên tục trên $\mathbb{R}$ thì 
		\begin{eqnarray*}
			f(1)&=&\lim\limits_{x\to 1^{-}} f(x)=\lim\limits_{x\to 1^{+}}  f(x) \\
			&\Leftrightarrow&  a+b-4=2a-2b\\
			 &\Leftrightarrow& a-3b=-4.
		\end{eqnarray*}
		Vậy  $P=-4$.
	 }
\end{ex}

% nguồn: Ok_TraLoiNgan_Lop11_Lan2
\begin{ex}%[1D3H3-4]%[ma:4106aac1]
	Gọi $S$ là tập các giá  trị của tham số $a$ để hàm số $y=f(x)=\heva{
	& x^2-3x  & \text{khi } x\neq1\\
	& a^2+a-8 &  \text{khi } x=1
	}$ liên  tục trên $\mathbb{R}$. Tính tích các phần tử của $S$.
	\shortans{-6} 
	\loigiai{ 
		 Khi $x\neq1$, $f(x)=x^2-3x$ là hàm đa thức nên liên tục trên $(-\infty;1)$ và $(1;+\infty)$.\\
		Hàm số liên tục tại $x=1$ khi và chỉ khi 
		\begin{eqnarray*}
			 \lim\limits_{x\to1} f(x)=f(1) &\Rightarrow& \lim\limits_{x\to1}(x^2-3x)=a^2+a-8 \\
			&\Rightarrow& 1-3 = a^2+a-8 \Rightarrow a^2+a-6=0. 
		\end{eqnarray*}
		Giải phương trình $a^2+a-6=0$  ta được
		 $\hoac{
		 & a=2\\
		&  a=-3.
		 }$\\
		Vậy tích các phần tử của $S$ là $2\cdot(-3)=-6$. 
	} 
\end{ex}

% nguồn: Ok_TraLoiNgan_Lop11_Lan2
\begin{ex}%[1D3H3-4]%[ma:91008817]
	Tìm tất cả  các giá trị của tham số thực $m$ sao cho hàm số $f(x)=\heva{&{2\sqrt{x}-m\text{khi} x\ge 0} \\
	&{mx+2\text{khi} x <  0}
	}$ liên  tục trên $\mathbb{R}$.
	\shortans{-2} 
	\loigiai{Trên khoảng $\left(0;+\infty \right)$ hàm số $f(x)=2\sqrt{x}-m$ là hàm số liên tục.\\ 
		Trên khoảng $\left(-\infty;0\right)$ hàm số $f(x)=mx+2$ là hàm số  liên tục.\\
		Ta có $\lim\limits_{x\to 0^{+}} f(x)=\lim\limits_{x\to 0^{+}} \left(2\sqrt{x}-m\right)=-m=f(0)$ và $\lim\limits_{x\to 0^{-}} f(x)=\lim\limits_{x\to  0^{-}} \left(mx+2\right)=2$.\\
		Hàm số $f(x)$ liên tục  trên $\mathbb{R}$ $\Leftrightarrow$ $\lim\limits_{x\to 0^{+}} f(x)=\lim\limits_{x\to 0^{-}} f(x)=f(0)$ $\Leftrightarrow-m=2\Leftrightarrow m=-2$.
	 }
\end{ex}

% nguồn: Ok_TraLoiNgan_Lop11_Lan2
\begin{ex}%[1D3H3-4]%[ma:f1736ef4]
	Một chất điểm chuyển động với tốc độ được cho  bởi hàm số $v(t)=\heva{&10+a & \text{khi }0\leqslant t\leqslant5\\ &t^2-5t+10 & \text{khi }t>5}$, trong đó $v(t)$ được tính theo đơn vị m/s và $t$ được tính theo giây. Tìm $a$ để hàm $v(t)$ liên tục tại điểm $t=5$.
	\shortans{0} 
	 \loigiai{
		Ta  có $v(5)=10+a$ và $\lim\limits_{t\to5^-}(10+a)=10+a$.\\
		 $\lim\limits_{t\to5^+}\left(t^2-5t+10\right)=10$.\\
		Hàm số $v(t)$ liên tục tại $t=5$  khi và chỉ khi $10+a=10\Rightarrow a=0$.
	} 
\end{ex}

% nguồn: Ok_TraLoiNgan_Lop11_Lan2
\begin{ex}%[1D3H3-4]%[ma:30f99a07]
	Biết hàm số  $f(x)=\heva{
	& \dfrac{x^2-16}{x-4}  & \text{khi }x>4\\
	 & mx+1 & \text{khi }x\leqslant4
	}$ liên tục  trên $\mathbb{R}$ khi tham số $m=\dfrac{a}{b}$ (phân số tối giản). Khi đó $2a+b$ bằng?
	\shortans{18} 
	\loigiai{ 
		Xét $x\in(4;+\infty)$ thì $f(x)=\dfrac{x^2-16}{x-4}$ là hàm phân thức nên  liên tục.\\
		Xét  $x\in(-\infty;4)$ thì $f(x)=mx+1$ là hàm đa thức nên liên tục.\\
		Xét $x=4$ 
		\begin{itemize}
			 \item $\lim\limits_{x\to4^+}f(x)=\lim\limits_{x\to4^+}\dfrac{x^2-16}{x-4}=\lim\limits_{x\to4^+}(x+4)=8$.
			 \item $\lim\limits_{x\to4^-}f(x)=\lim\limits_{x\to4^-}(mx+1)=4m+1=f(4)$.
		\end{itemize}
		Hàm số liên tục tại $x=4 \Leftrightarrow  4m+1=8 \Rightarrow m=\dfrac{7}{4}$.\\
		Vậy  $2a+b=2\cdot7+4=18$.
	 }
\end{ex}

% nguồn: Ok_TraLoiNgan_Lop11_Lan2
\begin{ex}%[1D3H3-4]%[ma:836dc9ab]
	Giá trị của tham số $m$ để  hàm số $f(x)=\heva{
	& \dfrac{x^2-5x+6}{\sqrt{4x-3}-x} & \text{khi  }x>3\\
	 & 1-m^2x & \text{khi }x\leqslant3
	}$ liên tục tại  $x=3$ viết dưới dạng $m=-\dfrac{a}{\sqrt{b}}, (m>0, a\in\mathbb{N}, b\in\mathbb{N}^*)$. Hiệu của $a-b$ bằng?
	 \shortans{-1}
	\loigiai{ 
		Điều kiện  liên tục tại $x=3$ là $\lim\limits_{x\to3^+}f(x)=\lim\limits_{x\to3^-}f(x)=f(3)$.\\
		 Ta có $\heva{
		 & f(3)=1-3m^2\\
		& \lim\limits_{x\to3^+}f(x)=\lim\limits_{x\to3^+}\dfrac{x^2-5x+6}{\sqrt{4x-3}-x}=\lim\limits_{x\to3^+}\dfrac{(x-2)(\sqrt{4x-3}+x)}{1-x}=-3\\ 
		&  \lim\limits_{x\to3^-}f(x)=\lim\limits_{x\to3^-}(1-m^2x)=1-3m^2
		 }$.\\
		 Do đó $1-3m^2=-3 \Rightarrow m=\pm\dfrac{2}{\sqrt{3}}$.\\
		Vì $m>0$, giá trị thỏa  yêu cầu là $m=\dfrac{2}{\sqrt{3}}$.\\
		Hiệu  số $a-b=-1$.
	 }
\end{ex}

% nguồn: Ok_TraLoiNgan_Lop11_Lan2
\begin{ex}%[1D3H3-4]%[ma:c5f9da5e]
	Cho hàm số $f(x)=\heva{&{\dfrac{x^2+3x-4}{x-1}\text{khi} x  > 1} \\
	 &{-2ax+1\text{khi} x\le 1}.
	}$ Xác định $a$ để hàm số liên tục trên $\mathbb{R}$. 
	 \shortans{-2}
	\loigiai{Tập xác định $\mathscr{D}={\bf \mathbb{R}}$.\\ 
		Ta thấy hàm số $f(x)=\dfrac{x^2+3x-4}{x-1}$  liên tục trên $\left(1;+\infty \right)$.\\
		Hàm  số $f(x)=-2ax+1$ liên tục trên $\left(-\infty;1\right)$.\\
		Ta có $f(1)=\lim\limits_{x\to 1^{-}} f(x)=\lim\limits_{x\to 1^{-}} \left(-2ax+1\right)=1-2a$.\\ 
		$\lim\limits_{x\to 1^{+}}  f(x)=\lim\limits_{x\to 1^{+}} \dfrac{x^2+3x-4}{x-1}=\lim\limits_{x\to 1^{+}} \left(x+4\right)=5$.\\
		Hàm số đã cho liên tục tại $x=1\Leftrightarrow f(1)=\lim\limits_{x\to 1^{-}} f(x)=\lim\limits_{x\to  1^{+}} f(x)\Leftrightarrow 1-2a=5\Leftrightarrow a=-2$.\\
		 Vậy với $a=-2$ thì hàm số liên tục trên $\mathbb{R}$.
	} 
\end{ex}

% nguồn: Ok_TraLoiNgan_Lop11_Lan2
\begin{ex}%[1D3H3-4]%[ma:f1cb3af7]
	Cho hàm  số $f(x)=\heva{&{3x+a-1\text{khi} x\le 0} \\
	&{\dfrac{\sqrt{1+2x}-1}{x}\text{khi}  x > 0}.
	}$ Tìm tất cả giá  trị thực của $a$ để hàm số đã cho liên tục trên $\mathbb{R}$.
	\shortans{2} 
	\loigiai{Tập xác định $\mathscr{D}=\mathbb{R}$.\\ 
		 Hàm số liên tục tại mọi điểm $x\ne 0$ với bất kỳ $a$.\\
		Ta có $f(0)=\lim\limits_{x\to  0^{-}} f(x)=\lim\limits_{x\to 0^{-}} \left(3x+a-1\right)=a-1$.\\
		 Mặt khác: $\lim\limits_{x\to 0^{+}} f(x)=\lim\limits_{x\to 0^{+}} \dfrac{\sqrt{1+2x}-1}{x}=\lim\limits_{x\to 0^{+}} \dfrac{2x}{x\left(\sqrt{1+2x}+1\right)}=\lim\limits_{x\to 0^{+}} \dfrac{2}{\sqrt{1+2x}+1}=1$.\\
		Hàm số liên tục trên $\mathbb{R}$ khi và chỉ khi hàm số liên  tục tại $x=0\Leftrightarrow a-1=1\Leftrightarrow a=2$.
	 }
\end{ex}

% nguồn: Ok_TraLoiNgan_Lop11_Lan2
\begin{ex}%[1D3H3-4]%[ma:18a0f90c]
	Một hãng  taxi đưa ra giá cước $T(x)$ (đồng) khi đi quãng đường $x$ (km) cho loại xe $4$ chỗ như sau:
	 $T(x)=\heva{
	& 10\,000+a & \text{khi }0<x\leqslant0{,}7\\ 
	& 11\,000+15\,100\cdot(x-0{,}7) & \text{khi  }0{,}7<x\leqslant30\\
	&  45\,3430+12\,000\cdot(x-30) & \text{khi }x>30
	 }$.
	Tìm $a$ để hàm số  $T(x)$ liên tục tại $x=0{,}7$.
	 \shortans{1000}
	 \loigiai{
		Tại $x=0{,}7$ ta có $T(0{,}7)=10\,000+a$ và $\lim\limits_{x\to0{,}7^-}T(x) =10\,000+a$.\\ 
		 $\lim\limits_{x\to0{,}7^+}T(x)=11\,000+15\,100\cdot(0{,}7-0{,}7)=11\,000$.\\
		Hàm số liên tục tại $x=0{,}7$ khi và chỉ khi $\lim\limits_{x\to0{,}7^-}T(x)=\lim\limits_{x\to0{,}7^+}T(x)=T(0{,}7)\Rightarrow a=1\,000$. 
	} 
\end{ex}

% nguồn: Ok_TraLoiNgan_Lop11_Lan2
\begin{ex}%[1D3H3-4]%[ma:da9c3bff]
	Tìm giá trị của $a$ để hàm số $f(x)=\heva{&{\dfrac{\sqrt[3]{3x+2}-2}{x-2} \text{khi} x > 2}  \\
	&{a^2x+\dfrac{1}{4}  \text{khi} x\le 2}
	}$ liên  tục trên $\mathbb{R}$.
	\shortans{0} 
	\loigiai{Ta thấy hàm  số $f(x)=\dfrac{\sqrt[3]{3x+2}-2}{x-2}$ xác định trên $\left(2;+\infty \right)$ nên hàm số liên tục trên $\left(2;+\infty \right)$.\\
		Hàm số $f(x)=a^2x+\dfrac{1}{4}$ xác định trên $\left(-\infty;2\right)$ nên hàm số  liên tục trên $\left(-\infty;2\right)$.
		Hàm số liên tục  trên $\mathbb{R}$ khi và chỉ khi hàm số liên tục tại $x=2$.\\
		Điều đó tương  đương với $\lim\limits_{x\to 2^{+}} f(x)=\lim\limits_{x\to 2^{-}} f(x)=f(2)$.\\
		Ta  có: $f(2)=2a^2+\dfrac{1}{4}$.\\
		\begin{eqnarray*}
			\lim\limits_{x\to 2^{+}} f(x)&=&\lim\limits_{x\to 2^{+}} \dfrac{\sqrt[3]{3x+2}-2}{x-2}=\lim\limits_{x\to  2^{+}} \dfrac{(3x+2)-8}{\left(x-2\right)\left(\sqrt[3]{\left(3x+2\right)^2}+2\sqrt[3]{3x+2}+4\right)}\\
			&=&\lim\limits_{x\to 2^{+}} \dfrac{3}{\left(\sqrt[3]{\left(3x+2\right)^2}+2\sqrt[3]{3x+2}+4\right)}=\dfrac{1}{4}.\\ 
			\lim\limits_{x\to  2^{-}} f(x)&=&\lim\limits_{x\to 2^{-}} \left(a^2x+\dfrac{1}{4} \right)=2a^2+\dfrac{1}{4}.
		\end{eqnarray*}
		Hàm số liên  tục tại $x=2$ $\Leftrightarrow$ $\lim\limits_{x\to 2^{+}} f(x)=\lim\limits_{x\to 2^{-}} f(x)=f(2)\Leftrightarrow 2a^2+\dfrac{1}{4}=\dfrac{1}{4} \Leftrightarrow a=0$.\\
		Vậy với $a=0$ thì hàm số liên tục trên $\mathbb{R}$. 
	 }
\end{ex}

% nguồn: Ok_TraLoiNgan_Lop11_Lan2
\begin{ex}%[1D3H3-4]%[ma:544e5a08]
	Tính tổng các giá trị tìm được của tham  số $m$ để hàm số $f(x)=\heva{
	&  m^2 x^2 & \text{khi }x\leqslant1\\
	 & (2-m)x & \text{khi }x>1
	}$ liên tục  trên $\mathbb{R}$.
	 \shortans{-1}
	\loigiai{ 
		Tập xác định $D=\mathbb{R}$.\\ 
		 Xét $x<1$ thì $f(x)=m^2 x^2$ liên tục trên $\mathbb{R}$ nên liên tục trên $(-\infty;1)$.\\
		Xét $x>1$ thì $f(x)=(2-m)x$ liên tục trên $\mathbb{R}$ nên liên tục trên  $(1;+\infty)$.\\
		Với $x=1$ ta  có $f(1)=m^2$, $\lim\limits_{x\to1^-}f(x)=m^2$, $\lim\limits_{x\to1^+}f(x)=2-m$.\\
		Hàm số  liên tục trên $\mathbb{R} \Rightarrow \lim\limits_{x\to1^-}f(x)=\lim\limits_{x\to1^+}f(x)=f(1) \Rightarrow m^2=2-m$.\\
		Giải phương trình $m^2=2-m$ ta  được
		$\hoac{ 
		&  m=1\\
		&  m=-2.
		 }$\\
		Vậy  tổng các giá trị $m$ là $1+(-2)=-1$.
	 }
\end{ex}

% nguồn: Ok_TraLoiNgan_Lop11_Lan1
\begin{ex}%[1D3V3-4]%[ma:881ca2f5]
	Cho các tham số $a$, $b$ thoả mãn hàm số $f(x) = \heva{& \dfrac{ax^2 - bx + 4}{x-1} & \text{ khi } x\ne 1\\ & 3 & \text{ khi } x = 1}$ liên tục  trên $\mathbb{R}$. Tính giá trị của biểu thức $P = a+b$.
	 \shortans{28}
	\loigiai{ 
		Do hàm số $f(x)$ liên  tục trên $\mathbb{R}$ nên hàm số $f(x)$ liên tục tại $x=1$, điều này tương đương với
		\[ 
		\lim\limits_{x\to 1} \dfrac{ax^2 - bx + 4}{x-1} = 3 \Leftrightarrow \lim\limits_{x\to 1}  \left( ax + a - b + \dfrac{a - b + 4}{x-1} \right) = 3 \Leftrightarrow \heva{& a- b + 4 = 0\\ & 2a - b = 3} \Leftrightarrow \heva{&a=7\\ &b=11.}
		 \]
		Vậy $a+b =  28$.
	} 
\end{ex}

% nguồn: Ok_TraLoiNgan_Lop11_Lan2
\begin{ex}%[1D3V3-4]%[ma:869fa721]
	 Cho hàm số $f(x)=\heva{&{\dfrac{a^2\left(x-2\right)}{\sqrt{x+2}-2} \text{khi} x > 2} \\
	&{\left(1-a\right)x\qquad \text{khi}  x\le 2}.
	}$ Tính tổng tất cả các giá trị của $a$ để hàm số liên tục  trên tập xác định.
	 \shortans{-0,5}
	\loigiai{Hàm số xác định trên  $\mathbb{R}$.\\
		 Với $x > 2$ ta có $f(x)=\dfrac{a^2\left(x-2\right)}{\sqrt{x+2}-2}$ là hàm số liên tục trên từng khoảng xác định.
		Do  đó hàm số $f(x)$ liên tục trên $\left(2;+\infty \right)$.\\
		Với $x < 2$ ta có $f(x)=\left(1-a\right)x$ là hàm số liên tục trên tập xác  định. Do đó hàm số $f(x)$ liên tục trên $\left(-\infty; 2\right)$.\\
		Với $x=2$  ta có $f(2)=2\left(1-a\right)$.\\
		Mặt khác,\\ 
		$\lim\limits_{x\to 2^{-}} f(x)=\lim\limits_{x\to 2^{-}}  \left(1-a\right)x=2\left(1-a\right)$.\\
		$\lim\limits_{x\to 2^{+}} f(x)=\lim\limits_{x\to 2^{+}} \dfrac{a^2\left(x-2\right)}{\sqrt{x+2}-2}=\lim\limits_{x\to 2^{+}}  a^2\left(\sqrt{x+2}+2\right)=4a^2$.\\
		Hàm số  liên tục trên $\mathbb{R}$ khi và chỉ khi hàm số liên tục tại $x=2$, nên
		\[\lim\limits_{x\to 2^{-}}  f(x)=\lim\limits_{x\to 2^{+}} f(x)=f(2)\Leftrightarrow 4a^2=2\left(1-a\right)\Leftrightarrow \hoac{&{a=-1} \\
		 &{a=\dfrac{1}{2}.}
		 }
		\] 
		Tổng các giá trị của tham số  $a$ là $ -1+\dfrac{1}{2}=-0{,}5$.
	 }
\end{ex}


%%%%%%%%%%%%%%%% TỰ LUẬN %%%%%%%%%%%%%%%%

% nguồn: Ok_TuanLuan_Toan11_Lan2
\begin{ex}%[1D3H3-4]%[ma:40da4437]
	Cho hàm số  $f(x)=\heva{
	& \dfrac{3-\sqrt{9-x}}{x}&&\text{khi } 0 < x < 9 \\ &  m&&\text{khi } x=0 \\ & \dfrac{3}{x}&&\text{khi }x\ge 9 \\ }$. Tìm $m$ để $f(x)$ liên tục trên $\left[0;+\infty \right)$.
	 \loigiai{
		\begin{itemize}
			\item Với $x\in (0;9)$, $f(x)=\dfrac{3-\sqrt{9-x}}{x}$  liên tục trên $(0;9)$.
			\item Với $x\in \left[9;+\infty \right)$,  $f(x)=\dfrac{3}{x}$ liên tục trên $\left[9;+\infty \right)$.
			\item Với  $x=0$ ta có $f(0)=m$.
		\end{itemize}
		Ta có $\lim\limits_{x\to0^+}f(x) =\lim\limits_{x\to0^+}\dfrac{3-\sqrt{9-x}}{x}$  $=\lim\limits_{x\to0^+}\dfrac{1}{3+\sqrt{9-x}}$ $=\dfrac{1}{6}$.\\
		Vậy để hàm số liên tục trên $[0;+\infty)$ khi nó phải liên tục tại $x=0$ $\Leftrightarrow$ $\lim\limits_{x\to0^+}f(x)=m$ $\Leftrightarrow  m=\dfrac{1}{6}$.
	} 
\end{ex}

% nguồn: Ok_TuanLuan_Toan11_Lan2
\begin{ex}%[1D3H3-4]%[ma:10359794]
	Cho  hàm số $f(x)=\heva{
	&\dfrac{x^3-4x^2+3}{x-1} && \text{khi } x\ne 1 \\ 
	& ax+\dfrac{5}{2}  && \text{khi } x=1 \\
	}$. Xác định $a$ để hàm số liên  tục trên $\mathbb{R}$.
	\loigiai{ 
		Với $x\ne 1$ ta có $f(x)=\dfrac{x^3-4x^2+3}{x-1}$ liên tục trên tập xác định.\\ 
		 $\lim\limits_{x\to1}\dfrac{x^3-4x^2+3}{x-1} =\lim\limits_{x\to1}\dfrac{(x^2-3x-3)(x-1)}{x-1}=-5$.\\
		 $f(1)=a+\dfrac{5}{2}$.\\
		Để hàm số liên  tục trên $\mathbb{R}$ thì hàm số phải liên tục tại $x=1$. \\
		Điều này xảy  ra khi
		 $\lim\limits_{x\to1}f(x)=f(1)$ $\Leftrightarrow a+\dfrac{5}{2}=-5$ $\Leftrightarrow a=-\dfrac{15}{2}$.
	} 
\end{ex}

% nguồn: Ok_TuanLuan_Toan11_Lan2
\begin{ex}%[1D3H3-4]%[ma:3cddf78a]
	Định $a$ để hàm  số $f(x)=\heva{
	& \dfrac{\sqrt{8x+1}-1}{x}&&\text{khi } x >  0 \\ & x^2+2x-4a&&\text{khi } x\le 0 \\ }$ liên tục trên $\mathbb{R}$.
	\loigiai{ 
		Trên $(0;+\infty)$, $(-\infty;0)$ hàm  số liên tục trên các khoảng đó\\
		Để hàm số  liên tục trên $\mathbb{R}$ ta cần hàm số liên tục tại $x=0$.\\
		Ta có 
		$\lim\limits_{x\to0^-}f(x)  =-4a, \lim\limits_{x\to0^+}f(x)=4, f(0)=-4a$. \\
		Vậy  hàm số liên tục trên $\mathbb{R}$ khi và chỉ khi $a=-1$.
	} 
\end{ex}

% nguồn: Ok_TuanLuan_Toan11_Lan2
\begin{ex}%[1D3H3-4]%[ma:dc491576]
	 Nếu hàm số $f(x)=\heva{
	& x^2+ax+b&&\text{khi }x <-5 \\ &  x+17&&\text{khi }-5\le x\le 10 \\ & ax+b+10&&\text{khi }x > 10 \\ }$ liên tục trên $\mathbb{R}$ thì $a+b$ bằng
	 \loigiai{
		\begin{itemize}
			\item Với $x <-5$ ,  $f(x)=x^2+ax+b$ là hàm đa thức nên liên tục trên $(-\infty;-5)$.
			\item Với $-5 < x < 10$ , $f(x)=x+7$  là hàm đa thức nên liên tục trên $(-5;10)$.
			\item Với $x > 10$ , $f(x)=ax+b+10$ là hàm đa thức nên  liên tục trên $(10;+\infty)$.
		\end{itemize}
		Để hàm số liên tục trên $\mathbb{R}$  thì hàm số phải liên tục tại $x=-5$ và $x=10$.\\
		Ta có: 
		$\begin{aligned}
		& f(-5)=12;  f(10)=17.\\
		 &\lim\limits_{x\to5^-}f(x)= \lim\limits_{x\to5^-}(x^2+ax+b) =-5a+b+25.\\
		& \lim\limits_{x\to5^+}f(x) =\lim\limits_{x\to5^+}(x+17)=12.\\ 
		 & \lim\limits_{x\to10^-}f(x) =\lim\limits_{x\to10^-}(x+17)=27.\\
		&  \lim\limits_{x\to10^+}f(x) =\lim\limits_{x\to10^+}(ax+b+10)=10a+b+10.
	\end{aligned}$\\
	Hàm số liên tục tại $x=-5$ và  $x=10$ khi
	$$\left\{\begin{aligned}
	& 5a+b+25=12 \\ &  10a+b+10=27 \\
	\end{aligned} \right.
	\Leftrightarrow \left\{\begin{aligned}
	&-5a+b=-13 \\ & 10a+b=17  \\
	\end{aligned} \right.
	\Leftrightarrow \left\{\begin{aligned}
	& a=2 \\ & b=-3 \\ 
	\end{aligned} \right.
	 \Rightarrow a+b=-1.$$
	 }
\end{ex}

% nguồn: Ok_TuanLuan_Toan11_Lan2
\begin{ex}%[1D3H3-4]%[ma:f36b99b5]
	Định $a$ để hàm  số $f(x)=\heva{
	& ax+\dfrac{1}{4} &&\text{khi } x\le2\\ & \dfrac{\sqrt[3]{3x+2}-2}{x-2}  &&\text{khi } x>2 \\ }$ liên tục trên $\mathbb{R}$.
	\loigiai{ 
		Hàm  số liên tục trên các khoảng $(2;+\infty)$, $(-\infty;2)$. \\
		Để hàm số liên tục trên $\mathbb{R}$ ta cần hàm số liên tục tại $x=2$.\\ 
		 Ta có
		$\begin{aligned}
		\lim\limits_{x\to2^+}f(x) &=\lim\limits_{x\to2^+}\dfrac{\sqrt[3]{3x+2}-2}{x-2} =\lim\limits_{x\to2^+}\dfrac{3x-6}{(x-2)\left(\sqrt[3]{3x+2}^2+2\sqrt[3]{3x+2}+4  \right)}=\dfrac{1}{4}. \\
		 \lim\limits_{x\to2^-}f(x) &=\dfrac{1}{4}, f(2)=2a+\dfrac{1}{4}. \\
	\end{aligned}$\\
	 Vậy hàm số liên tục trên $\mathbb{R}$ khi và chỉ khi $a=0$.
 }
\end{ex}

% nguồn: Ok_TuanLuan_Toan11_Lan2
\begin{ex}%[1D3H3-4]%[ma:2095b07a]
	 Cho hàm số $f(x)=\heva{
	& a^2x^2&&\text{khi }x\le \sqrt{2},a\in \mathbb{R} \\ & (2-a)x^2&&\text{khi }x > \sqrt{2} \\ }$. Tính giá trị của $a$ để $f(x)$ liên tục trên  $\mathbb{R}$.
	 \loigiai{
		 Tập xác định $\mathscr{D}=\mathbb{R}$.\\
		Với $x < \sqrt{2}$ ta có hàm số $f(x)=a^2x^2$ liên tục trên  khoảng $\left(-\infty;\sqrt{2}\right)$.\\
		Với $x > \sqrt{2}$ ta  có hàm số $f(x)=(2-a)x^2$ liên tục trên khoảng $\left(\sqrt{2};+\infty\right)$.\\
		Với $x=\sqrt{2}$ ta có  $f\left(\sqrt{2} \right)=2a^2$.\\
		$\lim\limits_{x\to\sqrt{2}^+}f(x)  =\lim\limits_{x\to\sqrt{2}^+}(2-a)x^2=2(2-a)$; $\lim\limits_{x\to\sqrt{2}^-}f(x) =\lim\limits_{x\to \sqrt{2}^-}a^2x^2 =2a^2$.\\
		Để hàm số liên tục tại $x=\sqrt{2}$. \\ 
		$\Leftrightarrow \lim\limits_{x\to\sqrt{2}^+}f(x)  =\lim\limits_{x\to\sqrt{2}^-}f(x)=f\left(\sqrt{2} \right)
		\Leftrightarrow  2a^2=2(2-a)
		 \Leftrightarrow a^2+a-2=0
		\Leftrightarrow \left[\begin{aligned}
		& a=1  \\ & a=-2. \\
	\end{aligned} \right.$\\
	Vậy $a=1$ hoặc $a=-2$ thì hàm số liên tục trên  $\mathbb{R}$.
} 
\end{ex}

% nguồn: Ok_TuanLuan_Toan11_Lan2
\begin{ex}%[1D3H3-4]%[ma:e2cd552c]
	Tìm $a$ để hàm số  $f(x)
	 =\heva{
	& x^2+2x+2&&\text{khi }x\ge 0 \\ & x^2+a&&\text{khi }x < 0 \\ 
	}$  liên tục trên $\mathbb{R}$.
	\loigiai{ 
		Hàm số liên tục trên các khoảng  $(0;+\infty)$, $(-\infty;0)$. \\
		 Để hàm số liên tục trên $\mathbb{R}$ ta cần hàm số liên tục tại $x=0$.\\
		Ta có $\lim\limits_{x\to0^+}f(x)=2$, $\lim\limits_{x\to0^-}f(x)=a$,  $f(0)=2$.\\
		Vậy hàm số liên tục tại $x_0=2$ khi và chỉ  khi $a=2$.
	} 
\end{ex}

% nguồn: Ok_TuanLuan_Toan11_Lan2
\begin{ex}%[1D3H3-4]%[ma:53e4fcf7]
	Tìm $a$  để hàm số $f(x)=\heva{
	&\dfrac{\sqrt{2x+1}-\sqrt{x+5}}{x-4} && \text{khi } x >  4 \\
	&\dfrac{\left(a+2 \right)x}{4} && \text{khi }  x\le 4 \\
	}$ liên tục trên tập xác định. 
	\loigiai{ 
		Tập xác định  $\mathscr{D}=\mathbb{R}$.\\
		Nhận xét: Hàm số $f(x)$  liên tục trên các khoảng $(-\infty;4)$ và $(4;+\infty)$. \\
		Do đó, để hàm số liên tục trên $\mathbb{R}$ ta cần tìm $a$ để hàm số liên tục tại $x=4$. \\ 
		Đều  kiện: $\lim\limits_{x\to4^+}f(x) =\lim\limits_{x\to4^-}f(x) =f(4)$. \\
		 $\lim\limits_{x\to4^+}f(x) =\lim\limits_{x\to4^+}\dfrac{\left(\sqrt{2x+1}-\sqrt{x+5} \right)\left(\sqrt{2x+1}+\sqrt{x+5} \right)}{(x-4)\left(\sqrt{2x+1}+\sqrt{x+5} \right)} =\lim\limits_{x\to4^+}\dfrac{1}{\sqrt{2x+1}+\sqrt{x+5}} =\dfrac{1}{6}$.\\
		$\lim\limits_{x\to4^-}f(x) =\lim\limits_{x\to4^-}\dfrac{(a+2)x}{4} =a+2$.  \\
		Cần có: $a+2=\dfrac{1}{6} \Leftrightarrow  a=-\dfrac{11}{6}$.
	 }
\end{ex}

% nguồn: Ok_TuanLuan_Toan11_Lan2
\begin{ex}%[1D3H3-4]%[ma:d5250477]
	Cho hàm số $f(x)=\heva{&\sqrt{2x-4}+3&&\text{khi } x\ge2\\&\dfrac{x+1}{x^2-2mx+3m+2}&&\text{khi } x<2}$.  Tìm tất cả các giá trị của tham số thực $m$ để hàm số liên tục trên $\mathbb{R}$.
	 \loigiai{
		\begin{itemize}[\bfseries Cách 1.]
			\item Hàm số xác định trên $\mathbb{R}$, liên tục trên khoảng  $(2;+\infty)$.\\
			Ta có 
			$\begin{aligned}
			&  f(2)=3 \\
			&  \lim\limits_{x \to 2^+} f(x) =\lim\limits_{x \to 2^+} \left(\sqrt{2x-4}+3\right)=3.
		\end{aligned}$\\
		Nếu $m=6$ thì  $\lim\limits_{x \to 2^+} f(x) =\lim\limits_{x \to 2^+} \dfrac{x+1}{x^2-12x+20}=-\infty$ nên hàm số không liên tục tại $x=2$.\\
		 Nếu $m\ne6$ thì ta có $\lim\limits_{x \to 2^+} f(x) =\lim\limits_{x \to 2^+} \dfrac{x+1}{x^2-2mx+3m+2}=\dfrac{3}{6-m}$ .\\
		Để hàm số liên tục tại $x=2$ thì $\dfrac{3}{6-m}=3 \Leftrightarrow 6-m=1 \Leftrightarrow m=5$.\\ 
		Với $m=5$ thì khi $x<2$, $f(x)=\dfrac{x+1}{x^2-10x+17}$ liên tục trên $(-\infty;2)$.\\ 
		Tóm lại với $m=5$ thì hàm  số đã cho liên tục trên $\mathbb{R}$.
		\item Hàm  số xác định trên $\mathbb{R}$, liên tục trên khoảng $(2;+\infty)$.\\
		Ta có $f(2)=3$; $\lim\limits_{x \to 2^+} f(x) =\lim\limits_{x  \to 2^+} \left(\sqrt{2x-4}+3\right)=3$.\\
	\end{itemize}
} 
\end{ex}

% nguồn: Ok_TuanLuan_Toan11_Lan2
\begin{ex}%[1D3H3-4]%[ma:cf1f7f51]
	Cho hàm số  $f(x)=\heva{
	&3x+a-1 && \text{khi } x\le  0 \\
	&\dfrac{\sqrt{1+2x}-1}{x} && \text{khi  } x > 0 \\
	}$. Tìm tất cả giá trị của $a$ để hàm số đã cho liên  tục trên $\mathbb{R}$.
	 \loigiai{
		Tập xác  định $\mathscr{D}=\mathbb{R}$.\\
		Ta có: Hàm số liên tục  trên các khoảng $(-\infty;0)$ và $(0;+\infty)$.\\
		$\lim\limits_{x\to0^-}f(x)  =\lim\limits_{x\to0^-}(3x+a-1)=a-1$. \\
		$\lim\limits_{x\to0^+}f(x) =\lim\limits_{x\to0^+}\dfrac{\sqrt{1+2x}-1}{x} =\lim\limits_{x\to0^+}\dfrac{2}{\sqrt{1+2x}+1}=1$.  \\
		$f\left(0 \right)=a-1.$  \\
		Hàm số liên tục trên $\mathbb{R} \Leftrightarrow$ Hàm số liên  tục tại điểm $x=0 \Leftrightarrow a-1=1 \Leftrightarrow a=2$.
	} 
\end{ex}

% nguồn: Ok_TuanLuan_Toan11_Lan2
\begin{ex}%[1D3H3-4]%[ma:f3015208]
	Tìm các khoảng liên tục của hàm  số
	\begin{enumEX}[]{4}
		 \item $y=x^3+3x^2+x$;
		 \item $y=\dfrac{x-1}{x+1}$;
		\item  $y=\dfrac{x+1}{x^2+x-2}$;
		\item $y=\tan x+\cos  x$.
	\end{enumEX}
	\loigiai{ 
		\begin{enumerate}
			 \item Hàm số $y=x^3+3x^2+x$ có tập xác định $\mathscr{D}=\mathbb{R}$ nên hàm số liên tục trên $\mathbb{R}$.
			\item Hàm số $y=\dfrac{x-1}{x+1}$ có tập xác định $\mathscr{D}=\mathbb{R}\backslash \{1\}$ nên hàm số liên  tục trên $(-\infty;1)$, $(1;+\infty)$.
			\item Hàm số $y=\dfrac{x+1}{x^2+x-2}$ có tập  xác định $\mathscr{D}=\mathbb{R}\backslash \{1;-2\}$ nên hàm số liên tục trên các khoảng $(-\infty;-2)$, $(-2;1)$, $(1;+\infty)$.
			\item Hàm số $y=\tan  x+\cos x$ có tập xác định $\mathscr{D}=\mathbb{R}\backslash \left\{\dfrac{\pi}{2}+k\pi, k\in \mathbb{Z} \right\}$ nên hàm số liên tục trên các khoảng $\left(-\dfrac{\pi}{2}+k2\pi;\dfrac{\pi}{2}+k2\pi \right)$.
		\end{enumerate}
	 }
\end{ex}

% nguồn: Ok_TuanLuan_Toan11_Lan2
\begin{ex}%[1D3H3-4]%[ma:f62f7527]
	Cho $a$, $b$ là hai số thực sao cho hàm  số $f(x) =\heva{
	&\dfrac{x^2+ax+b}{x-1} && \text{khi } x\ne 1  \\
	 & 2ax-1 && \text{khi } x=1 \\
	}$ liên tục  trên $\mathbb{R}$. Tính $a-b$.
	 \loigiai{
		Ta có $f(1)=2a-1$.\\ 
		Để hàm số liên tục trên $\mathbb{R}$ thì phải tồn tại $\lim\limits_{x\to1}f(x) =\lim\limits_{x\to1}\dfrac{x^2+ax+b}{x-1}$  và $\lim\limits_{x\to1}f(x) =f(1)$.\\
		Để tồn tại $\lim\limits_{x\to1}\dfrac{x^2+ax+b}{x-1}$ thì $(x^2+ax+b)\vdots (x-1) \Rightarrow  1+a+b=0 \Rightarrow b=-a-1$.\\
		Khi đó $\lim\limits_{x\to1}f(x)  =\lim\limits_{x\to1}\dfrac{x^2+ax+b}{x-1} =\lim\limits_{x\to1}\dfrac{(x-1)(x+a+1)}{x-1} =\lim\limits_{x\to1}(x+a+1)=a+2$.\\
		Do đó  để hàm số liên tục trên $\mathbb{R}$ thì $\lim\limits_{x\to1}f(x) =f(1)$. \\
		$\Leftrightarrow  2a-1=a+2 \Leftrightarrow a=3$.\\
		 Suy ra $b=-4$.\\
		Vậy  $a-b=7$.
	 }
\end{ex}

% nguồn: Ok_TuanLuan_Toan11_Lan2
\begin{ex}%[1D3V3-4]%[ma:528ce49c]
	Xét  tính liên tục của hàm số sau:
	\begin{listEX}[3]
		\item $y=x^2-2x+\sin  2x$.
		\item $y=\dfrac{x^2-3x+2}{2x^2+x-3}$. 
		\item $y=\sqrt{4-x}$. 
	\end{listEX}
	 \loigiai{
		\begin{itemize}
			\item[a)]  Ta thấy $x^2-2x$ là hàm đa thức và $\sin 2x$ là hàm số lượng giác có tập xác định $\mathscr{D}=\mathbb{R}$ nên hàm số $y=x^2-2x+\sin 2x$ liên tục trên $\mathbb{R}$.
			\item[b)] Điều kiện $2x^2+x-3\ne 0\Leftrightarrow \heva{&x \ne-\dfrac{3}{2} \\ &x \ne 1}$ nên tập xác định  của hàm số là: $\mathscr{D}=\mathbb{R} \setminus \left\{-\dfrac{3}{2};1\right\}$.\\
			Hàm số $y=\dfrac{x^2-3x+2}{2x^2+x-3}$ là hàm phân  thức hữu tỉ nên hàm số liên tục trên từng khoảng xác định là:
			 $(-\infty,-\dfrac{3}{2}); (-\dfrac{3}{2}, 1); (1,+\infty)$.
			\item[c)] Tập xác định  $\mathscr{D}=(-\infty, 4]$ nên hàm số $y=\sqrt{4-x}$ liên tục trên $(-\infty, 4]$.
		\end{itemize}
	 }
\end{ex}

% nguồn: Ok_TuanLuan_Toan11_Lan2
\begin{ex}%[1D3V3-4]%[ma:485f4f05]
	Tìm tất cả các giá trị của tham số $m$  để hàm số $y=f(x)=\heva{&\dfrac{\sqrt[3]{5x-7}+3}{4-\sqrt{4-3x}} & \text{ khi } x <-4\\ &\\
	&\left|mx \right|-\dfrac{41}{81} &  \text{ khi } x \geq-4}$
	liên tục trên  $(-\infty;+\infty)$.
	\loigiai{ 
		Hàm số $y=f(x)$ liên  tục trên $(-\infty;-4)$ và $(-4;+\infty)$.\\
		Để hàm số liên tục trên $(-\infty;+\infty)$ thì hàm số phải liên tục tại  $x=-4$.
		\begin{align*}
			\lim\limits_{x \to-4^-} f(x) &=\lim\limits_{x \to-4^-} \dfrac{\sqrt[3]{5x-7}+3}{4-\sqrt{4-3x}}  \\
			 &=\lim\limits_{x \to-4^-} \dfrac{5(x+4)\left(4+\sqrt{4-3x}\right)}{3(x+4)\left[\left(\sqrt[3]{5x-7}\right)^2-3\sqrt[3]{5x-7}+9\right]} \\
			 &=\lim\limits_{x \to-4^-} \dfrac{5\left(4+\sqrt{4-3x}\right)}{3\left[\left(\sqrt[3]{5x-7}\right)^2-3\sqrt[3]{5x-7}+9\right]}=\dfrac{40}{81}
		\end{align*}
		$\lim\limits_{x \to-4^+} f(x)=\lim\limits_{x \to-4^+} \left|mx \right|-\dfrac{41}{81}=\left|-4m \right|-\dfrac{41}{81}$ và  $f(-4)=\left|-4m \right|-\dfrac{41}{81}$.\\
		Hàm số  liên tục tại $x=-4$ khi và chỉ khi
		\begin{align*}
			&\lim\limits_{x \to-4^-} f(x) =\lim\limits_{x \to-4^+} f(x)=f(-4)  \\
			&\Leftrightarrow \left|-4m \right|-\dfrac{41}{81} =\dfrac{40}{81}  \\
			 &\Leftrightarrow|-4m|=1\\
			&\Leftrightarrow  4|m|=1\\
			&\Leftrightarrow|m|=\dfrac{1}{4}  \\
			&\Leftrightarrow m=\pm \dfrac{1}{4} 
		\end{align*}
	 }
\end{ex}

% nguồn: Ok_TuanLuan_Toan11_Lan2
\begin{ex}%[1D3V3-4]%[ma:783ccebb]
	Tìm giá trị tham số $m$ để hàm số $f(x)=\heva{&\dfrac{\sqrt{x+3}-2}{x-1}&&\text{khi}\,x > 1\\  &\\
	&mx &&\text{khi}\,x \leq 1}$ liên tục  trên $(-\infty;+\infty)$.
	\loigiai{ 
		Hàm  số $y=f(x)$ liên tục trên $(-\infty; 1)$ và $(1;+\infty)$. Để hàm số liên tục trên $\mathbb{R}$ thì hàm số phải liên tục tại $x=1$.\\
		Ta có \begin{align*}
		\lim\limits_{x \to 1^+} f(x) &=\lim\limits_{x \to 1^+}  \dfrac{\sqrt{x+3}-2}{x-1} \\
		&=\lim\limits_{x \to 1^+} \dfrac{\left(\sqrt{x+3}\right)^2-2^2}{(x-1)\left(\sqrt{x+3}+2\right)} \\ 
		&=\lim\limits_{x \to 1^+}  \dfrac{x-1}{(x-1)\left(\sqrt{x+3}+2\right)}\\
		 &=\lim\limits_{x \to 1^+} \dfrac{1}{\sqrt{x+3}+2}=\dfrac{1}{4}
	\end{align*}
	$\lim\limits_{x \to 1^-} f(x)=\lim\limits_{x \to 1^-} (mx)=m$ và  $f(1)=m$.\\
	Để hàm số liên tục tại $x=1\Leftrightarrow \lim\limits_{x \to 1^+} f(x)=\lim\limits_{x \to 1^-} f(x)=f(1)  \Leftrightarrow m=\dfrac{1}{4}$.
} 
\end{ex}

% nguồn: Ok_TuanLuan_Toan11_Lan2
\begin{ex}%[1D3V3-4]%[ma:858ac0f4]
	Tìm $m$ để hàm số $f(x)=\heva{&\dfrac{x^2-4}{x+2} & \text{ khi } x  \neq-2\\ &\\
	&m & \text{ khi  } x=-2}$
	 liên tục trên $\mathbb{R}$.
	\loigiai{ 
		Trên  $(-\infty;-2)$ và $(-2;+\infty)$, hàm số $f(x)=\dfrac{x^2-4}{x+2}$ xác định nên liên tục trên các khoảng $(-\infty;-2)$ và $(-2;+\infty)$.\\
		 Để hàm số liên tục trên $\mathbb{R}$ thì hàm số phải liên tục tại $x=-2$.\\
		$\lim\limits_{x  \to-2} f(x)=\lim\limits_{x \to-2} \dfrac{x^2-4}{x+2}=\lim\limits_{x \to-2} (x-2)=-4$.\\
		 Để hàm số liên tục tại $x_0=-2$ thì $\lim\limits_{x \to-2} f(x)=f(-2) \Rightarrow m=-4$.
	} 
\end{ex}

% nguồn: Ok_TuanLuan_Toan11_Lan2
\begin{ex}%[1D3V3-4]%[ma:c88cec38]
	Xét tính liên tục  của các hàm số sau trên tập xác định của nó.
	\begin{listEX}[2]
		\item Hàm  số
		$f(x)=\heva{ 
		&\dfrac{x^2-5x+6}{2x^3-16} && x <  2\\
		&2-x && x \ge  2.
		 }$
		\item Hàm số 
		 $f(x)=\heva{
		&\dfrac{x^2-3x+2}{|x-1|}  && x \ne 1\\
		&a  && x=1.
		 }$
	\end{listEX}
	 \loigiai{
		\begin{enumerate}
			\item Tập xác  định: $\mathscr{D}=\mathbb{R}$.\\
			Trên $(-\infty; 2)$ thì $f(x)=\dfrac{x^2-5x+6}{2x^3-16}$ xác  định nên hàm số liên tục trên $(-\infty; 2)$.\\
			Trên $(2;+\infty)$ thì $f(x)=2-x$  là hàm đa thức nên liên tục trên $(2;+\infty)$.\\
			Xét tính liên tục của hàm  số tại $x=2$.\\
			$\lim\limits_{x \to 2^-} f(x)=\lim\limits_{x \to 2^-} \dfrac{x^2-5x+6}{2x^3-16}=\lim\limits_{x \to 2^-} \dfrac{x-3}{2(x^2+2x+4)}=\dfrac{-1}{24}$ và $\lim\limits_{x  \to 2^+} f(x)=\lim\limits_{x \to 2^+} (2-x)=0$.\\
			Ta thấy $\lim\limits_{x \to 2^-} f(x) \neq \lim\limits_{x \to 2^+} f(x)$ nên hàm số  không có giới hạn tại $x=2$.\\
			Do đó hàm  số bị gián đoạn tại $x=2$.\\
			Vậy hàm số liên tục trên $(-\infty; 2)$ và  $(2;+\infty)$. Hàm số bị gián đoạn tại $x=2$.
			\item  Tập xác định: $\mathscr{D}=\mathbb{R}$.\\
			Trên $(-\infty; 1)$ và $(1;+\infty)$ hàm  số $f(x)=\dfrac{x^2-3x+2}{|x-1|}$ xác định nên hàm số liên tục trên $(-\infty; 1)$ và $(1;+\infty)$.\\
			Xét tính liên  tục của hàm số tại $x=1$.\\
			$\lim\limits_{x \to  1^-} f(x)=\lim\limits_{x \to 1^-} \dfrac{x^2-3x+2}{|x-1|}=\lim\limits_{x \to 1^-} \dfrac{(x-1)(x-2)}{x-1}=\lim\limits_{x \to 1^-} (x-2)=-1$.\\
			$\lim\limits_{x \to  1^+} f(x)=\lim\limits_{x \to 1^+} \dfrac{x^2-3x+2}{|x-1|}=\lim\limits_{x \to 1^+} \dfrac{(x-1)(x-2)}{1-x}=\lim\limits_{x \to 1^+} (2-x)=1$.\\
			Ta thấy $\lim\limits_{x \to 1^-} f(x) \neq \lim\limits_{x \to 1^+} f(x)$ nên hàm số không có giới hạn tại  $x=1$.\\
			 Vậy hàm số liên tục trên $(-\infty; 1)$ và $(1;+\infty)$. Hàm số bị gián đoạn tại $x=1$.
		\end{enumerate}
	} 
\end{ex}
